\documentclass[11pt,letterpaper,oneside]{book}
\pdfminorversion=7

\usepackage[margin=1in,headheight=14pt]{geometry}
\usepackage[T1]{fontenc}
\usepackage[utf8]{inputenc}
\usepackage[english]{babel}
\usepackage{lmodern}
\usepackage{microtype}
\usepackage{amsmath,amssymb,amsfonts,amsthm,bm}
\usepackage{booktabs,longtable,array}
\usepackage{graphicx}
\usepackage{placeins}
\usepackage[table]{xcolor}
\usepackage{hyperref}
\usepackage{enumitem}
\usepackage{tikz}
\usepackage{fancyhdr}
\usepackage{titlesec}
\usepackage[font=small,labelfont=bf,labelsep=period]{caption}
\usetikzlibrary{arrows.meta,positioning,shapes.geometric,fit,calc,backgrounds}

\hypersetup{
  colorlinks=true,
  linkcolor=blue!42!black,
  citecolor=blue!42!black,
  urlcolor=blue!55!black,
  pdfauthor={Jorge L. Mayorga Taborda},
  pdftitle={Spectral Grid Dynamics for Converter-Rich Power Systems: Operator-Valued Grid Strength, Rational Self-Energy, and Structural Fragility}
}

\definecolor{uniandesyellow}{HTML}{FFD200}
\definecolor{thesisblue}{HTML}{123A5A}
\definecolor{thesisink}{HTML}{1F2933}
\definecolor{thesisgray}{HTML}{667085}
\definecolor{thesisgreen}{HTML}{2E7D32}
\definecolor{thesisorange}{HTML}{C55A11}
\definecolor{thesisred}{HTML}{A23B3B}
\definecolor{thesisgold}{HTML}{B7791F}
\definecolor{thesisteal}{HTML}{0F766E}
\definecolor{softblue}{HTML}{EAF2F8}
\definecolor{softgreen}{HTML}{EAF7EA}
\definecolor{softorange}{HTML}{FFF0E6}
\definecolor{softgray}{HTML}{F2F4F7}
\definecolor{softyellow}{HTML}{FFF8D8}
\definecolor{softred}{HTML}{FDECEC}
\definecolor{softteal}{HTML}{E6F5F3}
\definecolor{coverink}{HTML}{202833}
\definecolor{coverblue}{HTML}{173B5C}
\definecolor{covergray}{HTML}{D7DCE0}
\definecolor{coverlight}{HTML}{ECEFF1}

\hypersetup{
  linkcolor=thesisink,
  citecolor=thesisblue!88!black,
  urlcolor=thesisorange!75!black,
  pdfborder={0 0 0}
}

\newcommand{\R}{\mathbb{R}}
\newcommand{\C}{\mathbb{C}}
\newcommand{\diag}{\operatorname{diag}}
\newcommand{\spec}{\operatorname{spec}}
\newcommand{\rank}{\operatorname{rank}}
\newcommand{\Real}{\operatorname{Re}}
\newcommand{\Imag}{\operatorname{Im}}
\newcommand{\norm}[1]{\left\lVert #1 \right\rVert}
\newcommand{\Lt}{\widetilde L}
\newcommand{\Dt}{\widetilde D}
\newcommand{\Sigmae}{\Sigma}
\newcommand{\zetaReq}{\zeta_{\rm req}}
\newcommand{\pu}{\mathrm{p.u.}}
\newcommand{\ones}{\mathbf 1}
\newcommand{\Ker}{\operatorname{ker}}
\newcommand{\Range}{\operatorname{range}}
\newcommand{\dist}{\operatorname{dist}}
\newcommand{\Res}{\operatorname{Res}}
\newcolumntype{L}[1]{>{\raggedright\arraybackslash}p{#1}}
\newcommand{\thead}[1]{\textbf{\color{thesisink}#1}}

\newcommand{\thesisbox}[3]{%
\par\smallskip
\noindent\begin{tikzpicture}
\node[
  draw=#1!70!black,
  fill=#1!8,
  rounded corners=3pt,
  line width=0.55pt,
  inner sep=9pt,
  text width=0.94\textwidth
] {\textbf{\color{#1!55!black}#2}\par\vspace{0.22em}#3};
\end{tikzpicture}
\par\smallskip
}

\newcommand{\statuspill}[2]{%
\tikz[baseline=(txt.base)]\node[
  rounded corners=7pt,
  draw=#1!70!black,
  fill=#1!12,
  inner xsep=5pt,
  inner ysep=1.8pt,
  line width=0.35pt
] (txt) {\scriptsize\bfseries #2};%
}

\newtheorem{definition}{Definition}[chapter]
\newtheorem{proposition}{Proposition}[chapter]
\newtheorem{theorem}{Theorem}[chapter]
\newtheorem{remark}{Remark}[chapter]
\newtheorem{conjecture}{Conjecture}[chapter]
\newtheorem{example}{Example}[chapter]

\setlength{\parindent}{0pt}
\setlength{\parskip}{0.48em plus 0.12em minus 0.08em}
\setlength{\tabcolsep}{5.4pt}
\renewcommand{\arraystretch}{1.16}
\setlist[itemize]{leftmargin=*,topsep=0.35em,itemsep=0.2em}
\setlist[enumerate]{leftmargin=*,topsep=0.35em,itemsep=0.2em}

\titleformat{\chapter}[display]
  {\normalfont\huge\bfseries\color{thesisink}}
  {\filleft\Large\color{thesisgray}\chaptertitlename\ \thechapter}
  {0.8ex}
  {\titlerule[0.6pt]\vspace{1.2ex}\filright}
  [\vspace{1.1ex}{\color{uniandesyellow}\titlerule[2pt]}]
\titlespacing*{\chapter}{0pt}{0.25in}{0.45in}

\titleformat{name=\chapter,numberless}[display]
  {\normalfont\huge\bfseries\color{thesisink}}
  {}
  {0pt}
  {\filright}
  [\vspace{1.1ex}{\color{uniandesyellow}\titlerule[2pt]}]

\titleformat{\section}[hang]
  {\normalfont\Large\bfseries\color{thesisink}}
  {\color{thesisblue}\thesection}
  {0.8em}
  {}
  [\vspace{0.15ex}{\color{uniandesyellow}\titlerule[1.1pt]}]
\titlespacing*{\section}{0pt}{1.45em plus 0.25em minus 0.15em}{0.65em}

\titleformat{\subsection}[hang]
  {\normalfont\large\bfseries\color{thesisblue}}
  {\thesubsection}
  {0.75em}
  {}
\titlespacing*{\subsection}{0pt}{1.1em plus 0.2em minus 0.1em}{0.35em}

\titleformat{\subsubsection}[hang]
  {\normalfont\normalsize\bfseries\color{thesisink}}
  {\thesubsubsection}
  {0.65em}
  {}
\titlespacing*{\subsubsection}{0pt}{0.9em plus 0.15em minus 0.1em}{0.25em}

\pagestyle{fancy}
\fancyhf{}
\renewcommand{\chaptermark}[1]{\markboth{Chapter \thechapter}{}}
\fancyhead[L]{\small\color{thesisgray}Spectral Hosting in Converter-Rich Grids}
\fancyhead[R]{\small\color{thesisblue}\leftmark}
\fancyfoot[C]{\thepage}
\renewcommand{\headrulewidth}{0.35pt}
\renewcommand{\footrulewidth}{0pt}

\fancypagestyle{plain}{
  \fancyhf{}
  \fancyfoot[C]{\thepage}
  \renewcommand{\headrulewidth}{0pt}
}

\begin{document}

\hypersetup{pageanchor=false}
\begin{titlepage}
\thispagestyle{empty}
% Redesigned institutional cover. The former decorative triangle/concept-box
% treatment was replaced by a restrained graph--matrix--spectrum composition.
\begin{tikzpicture}[remember picture,overlay]
  \fill[white] (current page.south west) rectangle (current page.north east);
  \fill[uniandesyellow] (current page.north west) rectangle
    ([xshift=0.23in]current page.south west);
  \draw[coverink,line width=0.45pt]
    ([xshift=0.23in]current page.north west) --
    ([xshift=0.23in]current page.south west);
  \draw[coverink,line width=0.5pt]
    ([xshift=0.44in,yshift=-0.24in]current page.north west) --
    ([xshift=-0.42in,yshift=-0.24in]current page.north east);
  \draw[coverink,line width=0.5pt]
    ([xshift=0.44in,yshift=0.43in]current page.south west) --
    ([xshift=-0.42in,yshift=0.43in]current page.south east);
  \node[
    rotate=90,
    anchor=center,
    font=\tiny\bfseries,
    text=coverink,
    opacity=0.72
  ] at ([xshift=0.135in]current page.west) {SPECTRAL GRID DYNAMICS};
  \foreach \y in {0.18,0.29,0.40,0.51,0.62,0.73,0.84} {
    \draw[coverink!32,line width=0.35pt]
      ([xshift=0.075in,yshift=-\y\paperheight]current page.north west) --
      ([xshift=0.155in,yshift=-\y\paperheight]current page.north west);
  }
  \node[
    anchor=center,
    opacity=0.016,
    inner sep=0pt
  ] at ([xshift=0.46\paperwidth,yshift=-0.03\paperheight]current page.center) {
    \includegraphics[
      width=2.20\paperwidth,
      keepaspectratio
    ]{assets/uniandes_logo.pdf}
  };
\end{tikzpicture}
\centering
\vspace*{0.34in}
\includegraphics[width=0.285\textwidth,keepaspectratio]{assets/uniandes_logo_color.pdf}\par
\vspace{0.22in}
{\large\color{coverink}Faculty of Engineering\par}
{\large\color{coverink}Department of Electrical and Electronic Engineering\par}
\vspace{0.24in}
{\color{uniandesyellow}\rule{0.38\textwidth}{1.35pt}\par}
\vspace{0.20in}
{\fontsize{32}{38}\selectfont\bfseries\color{coverink}
Spectral Grid Dynamics\par
\vspace{0.02in}
for Converter-Rich Power Systems\par}
\vspace{0.20in}
\begin{minipage}{0.79\textwidth}
\centering
{\Large\bfseries\color{coverblue}
Operator-Valued Grid Strength,\par
Rational Self-Energy, and Structural Fragility\par}
\end{minipage}\par
\vspace{0.25in}
\begin{tikzpicture}[
  x=1cm,
  y=1cm,
  modulelabel/.style={font=\scriptsize\scshape\color{coverink}},
  thinlink/.style={covergray,line width=0.45pt},
  flow/.style={-{Latex[length=2.2mm,width=1.4mm]},coverink!58,line width=0.55pt},
  nodebase/.style={circle,draw=coverink!45,fill=white,line width=0.45pt,inner sep=0pt,minimum size=4.2pt},
  converter/.style={circle,draw=uniandesyellow!92!black,fill=white,line width=1.05pt,inner sep=0pt,minimum size=6.2pt},
  eig/.style={circle,fill=coverink!50,inner sep=1.15pt}
]
  \draw[coverlight,line width=0.45pt] (-6.25,1.78) -- (6.25,1.78);
  \draw[coverlight,line width=0.45pt] (-6.25,-1.45) -- (6.25,-1.45);
  \node[modulelabel] at (-4.75,1.45) {Network graph};
  \coordinate (n1) at (-5.85,0.90);
  \coordinate (n2) at (-5.10,1.22);
  \coordinate (n3) at (-4.30,0.82);
  \coordinate (n4) at (-3.57,1.17);
  \coordinate (n5) at (-5.62,0.08);
  \coordinate (n6) at (-4.83,-0.22);
  \coordinate (n7) at (-4.02,0.12);
  \coordinate (n8) at (-3.28,0.38);
  \coordinate (n9) at (-5.18,-0.88);
  \coordinate (n10) at (-4.35,-0.66);
  \coordinate (n11) at (-3.62,-0.92);
  \coordinate (n12) at (-2.95,-0.42);
  \draw[thinlink] (n1)--(n2)--(n3)--(n4)--(n8)--(n12)--(n11)--(n10)--(n9)--(n5)--(n1);
  \draw[thinlink] (n2)--(n5) (n2)--(n6) (n3)--(n6) (n3)--(n7)
    (n4)--(n7) (n5)--(n6)--(n7)--(n8) (n6)--(n10) (n7)--(n11);
  \foreach \p in {n1,n2,n4,n5,n7,n9,n10,n12}
    \node[nodebase] at (\p) {};
  \foreach \p in {n3,n6,n11}
    \node[converter] at (\p) {};
  \node[font=\scriptsize\color{covergray!70!coverink}] at (-5.78,-0.58) {AC};

  \draw[flow] (-2.55,0.08) -- (-1.47,0.08);

  \node[modulelabel] at (0,1.45) {Sparse matrix};
  \begin{scope}[shift={(-0.54,-0.94)}]
    \draw[coverink!40,line width=0.35pt] (0,0) rectangle (1.88,1.88);
    \foreach \i in {0,...,8} {
      \draw[coverlight,line width=0.25pt] ({\i*0.209},0) -- ({\i*0.209},1.88);
      \draw[coverlight,line width=0.25pt] (0,{\i*0.209}) -- (1.88,{\i*0.209});
    }
    \foreach \x/\y in {
      0/8,1/7,2/6,3/5,4/4,5/3,6/2,7/1,8/0,
      0/5,1/3,2/8,3/1,4/7,5/0,5/6,6/4,7/8,8/2,
      0/1,2/3,3/8,6/7,8/5}
      \fill[coverink!34] ({\x*0.209+0.045},{\y*0.209+0.045}) rectangle ++(0.119,0.119);
    \foreach \x/\y in {1/7,4/4,6/2}
      \fill[uniandesyellow!86!coverink] ({\x*0.209+0.038},{\y*0.209+0.038}) rectangle ++(0.133,0.133);
  \end{scope}

  \draw[flow] (1.47,0.08) -- (2.55,0.08);

  \node[modulelabel] at (4.75,1.45) {Spectral plane};
  \draw[coverink!55,line width=0.45pt,-{Latex[length=1.7mm]}] (3.28,-0.93) -- (6.10,-0.93)
    node[right,font=\scriptsize\color{coverink!70}] {\(\Real(\lambda)\)};
  \draw[coverink!55,line width=0.45pt,-{Latex[length=1.7mm]}] (4.12,-1.22) -- (4.12,1.07)
    node[above,font=\scriptsize\color{coverink!70}] {\(\Imag(\lambda)\)};
  \foreach \x/\y in {3.62/-0.26,3.70/0.36,4.70/-0.12,4.86/0.54,5.30/-0.42,5.55/0.24,5.82/-0.66}
    \node[eig] at (\x,\y) {};
  \node[circle,draw=uniandesyellow!92!black,line width=1.05pt,inner sep=3.2pt] at (3.70,0.36) {};

  \draw[uniandesyellow!88!coverink,densely dotted,line width=0.9pt]
    (n3) .. controls (-2.95,1.52) and (-0.25,1.62) .. (0.09,0.73)
    .. controls (1.52,-0.02) and (2.72,0.80) .. (3.70,0.36);
\end{tikzpicture}\par
\vspace{0.18in}
{\small\bfseries\color{coverink}
SPECTRAL GRAPHS \(\cdot\) OPERATOR-VALUED STRENGTH \(\cdot\) SELF-ENERGY \(\cdot\) STRUCTURAL FRAGILITY\par}
\begin{tikzpicture}[remember picture,overlay]
  \node[
    anchor=south,
    align=center,
    text=coverink,
    text width=0.84\paperwidth,
    inner sep=0pt
  ] at
    ([yshift=0.88in]current page.south) {%
    {\Large\bfseries Jorge L. Mayorga Taborda\par}
    \vspace{0.13in}
    {\large Doctoral Dissertation Theoretical Draft\par}
    \vspace{0.07in}
    {\large Submitted as a research thesis manuscript in Electrical Engineering\par}
    \vspace{0.07in}
    {\large Bogota, Colombia \quad June 2026\par}
    \vspace{0.09in}
    {\small
    \textbf{Research area:} Power-system dynamics and spectral methods\par
    \textbf{Primary artifact:} Reproducible LaTeX theoretical thesis\par
    \textbf{Repository:} \texttt{paper-P06-spectral-graph-power-dynamics}\par}};
\end{tikzpicture}
\vfill
\end{titlepage}

\cleardoublepage
\thispagestyle{empty}
\vspace*{1.0cm}
\begin{center}
{\Large\bfseries Dissertation Information\par}
\vspace{0.45cm}
{\color{uniandesyellow}\rule{0.32\textwidth}{2pt}\par}
\vspace{1.0cm}
\begin{tabular}{>{\bfseries}p{0.28\textwidth}p{0.55\textwidth}}
Institution & Universidad de los Andes\\
Faculty & Faculty of Engineering\\
Program & Electrical Engineering\\
Author & Jorge L. Mayorga Taborda\\
Document type & Doctoral dissertation theoretical draft\\
Research area & Power-system dynamics, spectral graph methods, and converter-rich grid stability\\
Repository & \texttt{paper-P06-spectral-graph-power-dynamics}\\
\end{tabular}
\end{center}
\vfill
{\small
This theoretical draft is prepared as a professional academic manuscript. The
front matter uses the official Universidad de los Andes logosymbol for
institutional identification. The technical content remains a reproducible
LaTeX artifact under \texttt{reports/thesis/}.
\par}

\cleardoublepage
\hypersetup{pageanchor=true}
\frontmatter

\chapter*{Abstract}
\addcontentsline{toc}{chapter}{Abstract}
This dissertation develops an operator-valued theory of dynamic grid strength
for converter-rich power systems. The central claim is that grid strength is
not a scalar such as short-circuit ratio, algebraic connectivity, or a single
damping ratio. It is a port-dependent dynamic operator whose static limit is
encoded by graph stiffness, whose oscillatory behavior is carried by the Green
function of a quadratic or rational eigenvalue problem, whose
controller-mediated behavior is represented by rational self-energy, whose
fragility is measured by structured pseudospectra, and whose nonlinear
extension appears as memory. In the classical swing-equation model,
\[
  M\ddot\theta+D\dot\theta+L\theta=p
\]
the damping matrix \(D\) is fixed. Exact graph-mode decoupling occurs when the
inertia-normalized stiffness and damping matrices commute. The special case
\(D=\gamma M\) is therefore not the whole decouplable class; it is the uniform
damping-to-inertia case in which all underdamped modes share the same real
decay rate. Rayleigh damping \(D=\alpha M+\beta L\), and more generally
commuting damping operators, define a broader constant-\(D\) floor. If \(D\) is
fixed but non-commuting, the modal damping matrix
\(Q^TM^{-1/2}DM^{-1/2}Q\) is not diagonal; graph eigenvectors still organize
stiffness, but they no longer diagonalize damping. If converter, phase-locked
loop, exciter, governor, stabilizer, or auxiliary control states are condensed,
the reduced object is not a fixed-coefficient quadratic eigenvalue problem. It
is a rational nonlinear eigenvalue problem,
\[
  T(s)=s^2M+sD_0+L+\Pi(s).
\]
Only after the rational correction is separated into
\(\Pi(s)=\Pi_K(s)+s\Pi_D(s)\) does the damping-channel specialization
\(\Sigma(s)=D_0+\Pi_D(s)\) appear.

The original contribution is not a new Laplacian metric or a new eigenvalue
solver. It is a mechanism-resolved operator framework showing when
graph-theoretic strength is a valid damping certificate, when
damping-to-inertia roughness creates modal coupling, when topology actions
must be evaluated through spectral effective resistance, and when controller
self-energy can reverse the dynamic effect predicted by static graph screens.

The thesis is deliberately theoretical and is organized as a doctoral
monograph rather than as an extended article. It develops the physical
framework, mathematical foundations, state of the art, operational graph
dynamics, constant-damping theory, rational self-energy, nonlinear memory, and
applications as separable modules.
Its unifying object is the
dynamic Green function \(G_T(s)=T(s)^{-1}\). For a line \(e=(i,j)\), the scalar
\(\mathcal R_e(s)=a_e^TG_T(s)a_e\), with \(a_e=e_i-e_j\), is a dynamic
effective resistance: at \(s=0\) it recovers classical graph effective
resistance, while away from zero it records inertia, damping, stiffness,
controller poles, residues, non-normality, and modal coupling. A finite line
action then satisfies the exact scalar equation
\(1+\Delta b_e\mathcal R_e(s)=0\); general low-rank actions satisfy the
operator determinant law
\(\det T_u(s)=\det T(s)\det[I+K_u(s)V^*T(s)^{-1}U]\). These identities turn
local eigenvalue sensitivities into first-order limits of exact finite-action
laws. The remaining modules add port-equivalence of self-energy,
model-adequacy certificates, structured fragility radii, directional
passivity, data-identified self-energy, and nonlinear memory. Simulation cases
and internal project drafts are treated as validation gates and traceability
records, not as universal theorems. The main contribution is a coherent
mathematical language for deciding when spectral graph intuition is a
certificate, when it is only a screen, when a network improvement can reduce a
dynamic margin, and when controller-aware rational self-energy or nonlinear
memory is required.

\textbf{Keywords:} operator-valued grid strength, operational graph,
graph-time resonance, dynamic effective resistance, Green function, damping
ratio, inverter-based resources, quadratic eigenvalue problem, Schur
complement, nonlinear eigenvalue problem, rational self-energy, structured
pseudospectrum, dynamic hosting capacity, virtual inertia, Braess-like
reversal, weak nodes, weak links, hidden damping margin, risk vector, ANDES,
IEEE 39-bus system.

\chapter*{Scope and Reproducibility Statement}
\addcontentsline{toc}{chapter}{Scope and Reproducibility Statement}
This document treats the current PDF as version 0 of the intellectual map, not
as a manuscript to be enlarged linearly. It is a doctoral theoretical draft,
not an experimental report. It does not present a final benchmark, a calibrated
IEEE test case, or a claim of empirical closure. Its purpose is to establish
the modular mathematical structure that later numerical experiments must
respect. Claims about constant damping, non-commuting modal coupling, Schur
condensation, rational self-energy, and nonlinear memory are presented as
theoretical statements, derivations, or chapter contracts. Claims about
practical planning are formulated as requirements for future validation.

All source files for this manuscript are kept under \texttt{reports/thesis/}.
The compiled PDF is copied to \texttt{outputs/pdf/thesis\_draft.pdf}. Scratch
material in \texttt{temp/} is preserved as historical context and is not part of
the theoretical thesis source. The integrated claim ledger is kept out of the
main scientific narrative and appears as an appendix-level traceability record.

\chapter*{Graphical Roadmap}
\addcontentsline{toc}{chapter}{Graphical Roadmap}
The thesis is organized as a chain from electrical modeling to dynamic
planning. The diagram below is the reader's map: each block introduces one
object that must be carried forward without changing its meaning.

\begin{figure}[ht]
\centering
\begin{tikzpicture}[
  roadmap/.style={
    rectangle,
    rounded corners=3pt,
    draw=#1!70!black,
    fill=#1!10,
    text width=0.20\textwidth,
    minimum height=1.28cm,
    align=center,
    font=\small
  },
  flow/.style={-{Latex[length=2.2mm]}, thick, draw=thesisgray}
]
\node[roadmap=thesisblue] (pf) {\textbf{AC operating point}\\
\(I=Y_{\rm bus}V\), \(S=VI^*\)};
\node[roadmap=thesisgreen, right=0.55cm of pf] (lin) {\textbf{Operational graph}\\
\(L=B W(z^\star)B^T\)};
\node[roadmap=thesisorange, right=0.55cm of lin] (qep) {\textbf{Graph-time response}\\
\(H_G(\omega)\), \(G(s)\)};
\node[roadmap=thesisred, right=0.55cm of qep] (nep) {\textbf{Controller memory}\\
\(s^2M+sD_0+L+\Pi(s)\)};
\node[roadmap=thesisgold, below=0.8cm of lin, xshift=2.7cm] (plan) {\textbf{Planning margins}\\
\(\mathcal R_e(s)\), fragility radius, certified actions};
\draw[flow] (pf) -- (lin);
\draw[flow] (lin) -- (qep);
\draw[flow] (qep) -- (nep);
\draw[flow] (qep.south) -- ++(0,-0.42) -| ([xshift=-0.32cm]plan.north);
\draw[flow] (nep.south) -- ++(0,-0.42) -| ([xshift=0.32cm]plan.north);
\end{tikzpicture}
\caption{Visual roadmap of the dissertation. The thesis starts from the AC
operating point and the physical network model, derives the operational graph
\(L=BWB^T\), builds graph-time response and the Green function of the
fixed-\(D\) spectral dynamics,
replaces fixed damping by rational self-energy when controllers are condensed,
separates nonlinear memory from NEP theory, and ends with certified action and
fragility planning.}
\label{fig:front_roadmap}
\end{figure}

\thesisbox{thesisblue}{How to read the dissertation}{
The main text follows the modular path: physical framework, mathematical
foundations, operational graph dynamics, constant damping, dynamic Green
functions, rational self-energy, nonlinear memory, and applications. A graph
metric is accepted only after the dynamic operator, protected margin, action
family, and certification gate attached to it are named.}

\FloatBarrier

\chapter*{Acknowledgments}
\addcontentsline{toc}{chapter}{Acknowledgments}
This draft consolidates theory, reproducible experiments, and editorial
material from the project into a thesis-oriented theoretical structure. I am
grateful to the academic and engineering communities whose work on power-system
stability, spectral graph methods, quadratic eigenvalue problems, nonlinear
eigenvalue problems, and converter-dominated dynamics makes this synthesis
possible.

\tableofcontents
\listoffigures
\listoftables

\mainmatter

\part{Foundations and State of the Art}
\chapter{Introduction}

Power-system small-signal stability is a theory of operators before it is a
theory of plots. A linearized power system is stable or unstable because the
spectrum of its differential-algebraic representation lies in particular
regions of the complex plane, because the corresponding eigenvectors project
onto physical states in particular ways, and because perturbations of
parameters move those eigenvalues through mechanisms that can be traced
mathematically. The familiar language of damping ratio, participation factors,
critical modes, and modal frequencies is a compact engineering vocabulary for
that operator theory \cite{kundur1994power,rogers2000oscillations,kundur2004definition}.

The same vocabulary becomes harder to use in converter-rich systems. A
synchronous-machine dominated model permits a useful separation: inertia is
stored in rotating masses, synchronizing stiffness comes from the network and
power-angle relation, and damping can often be represented by a fixed local
coefficient. Converter-interfaced resources alter this separation. A
grid-following converter may contribute phase-locked-loop states, current-loop
states, plant controllers, filters, and protection-related dynamics. A
grid-forming converter may contribute droop, virtual oscillator, virtual
synchronous-machine, voltage-loop, current-loop, and current-limiting dynamics.
The extended IEEE/CIGRE stability taxonomy explicitly recognizes that
power-electronic interfaces introduce faster and device-dependent dynamic
phenomena that do not fit cleanly into the classical electromechanical
classification \cite{hatziargyriou2021stability}.

This dissertation asks a sharper theoretical question:

\begin{quote}
\itshape
What is the reduced dynamic operator that preserves grid strength, fragility,
and actionable pole motion in a converter-rich power system?
\end{quote}

That question has a planning consequence. If a utility asks how much
converter-based generation can be hosted, where virtual inertia should be
placed, or which lines should be reinforced before installing batteries, the
answer cannot be read from voltage limits or graph connectivity alone. A
hosting decision is dynamic only if it states the damping, resolvent, or
protection margin that must survive the conversion. A line decision is robust
only if it separates useful stiffness increase from damping-ratio loss or
controller-resonant reinforcement risk.

The answer developed here is hierarchical, but the hierarchy is not the final
claim. Fixed scalar damping is the uniform-decay member of a broader sequence
of increasingly expressive operators. The exact decoupling condition for fixed
damping is not \(D=\gamma M\) alone. It is commutation of the
inertia-normalized stiffness and damping operators. A fixed but non-commuting
\(D\) keeps the problem quadratic, but it couples graph modes through the modal
damping matrix. A controller-condensed converter-rich model changes the problem
more fundamentally: the reduced correction depends on the spectral parameter
\(s\). The effective operator is then a rational self-energy \(\Pi(s)\),
obtained by a Schur complement of the eliminated controller states. The object
that turns this hierarchy into an action theory is the dynamic Green function
\(G_T(s)=T(s)^{-1}\), whose port projections measure how finite physical
actions move poles.

\section{Thesis Statement}

The thesis statement is:

\begin{quote}
Dynamic grid strength in converter-rich power systems is an operator-valued
property. Its static limit is encoded by graph stiffness, its oscillatory
behavior by the Green function of the QEP, its controller-mediated behavior by
rational self-energy \(\Pi(s)\), its finite-action behavior by low-rank
determinant identities, its fragility by structured pseudospectra, and its
nonlinear extension by memory kernels. The damping-channel symbol
\(\Sigma(s)\) is reserved for the velocity channel after
\(\Pi(s)=\Pi_K(s)+s\Pi_D(s)\) has been declared. Spectral graph methods remain
valuable, but their role changes from certificate to screen whenever damping
and stiffness cease to share an eigenbasis, whenever eliminated controllers
contribute frequency-dependent dynamics, or whenever an action claim requires
a certified finite perturbation rather than a local derivative.
\end{quote}

This statement does not claim novelty for Laplacian linearizations,
generalized graph modes, effective resistance, QEP or NEP sensitivity
formulas, Schur complements, or the existence of Braess-like phenomena in
oscillator and power networks. Those are established mathematical and
engineering ingredients. The novelty sought by this theoretical thesis is the
synthesis: a single operator hierarchy that explains when classical damping
intuition is mathematically justified, when it fails, and what reduced object
must replace it.

\thesisbox{thesisorange}{Originality boundary}{
Known mathematics is used as the foundation; it is not claimed as the
dissertation's invention. The doctoral contribution is the operator-valued
planning synthesis. The strongest original claims are the graph roughness of
\(D_i/M_i\) as a modal-coupling obstruction, the graph-time response atlas,
the spectral effective resistance law for finite topology updates, and the
controller-induced reinforcement-reversal certificate. Claims about universal
weak nodes, completed sheaf-grid theory, or empirical benchmark closure remain
validation targets rather than defended theorems.}

\thesisbox{thesisgreen}{Doctoral reading contract}{
Every planning statement in this dissertation must answer four questions:
which operator is being used, which port or action family is being tested,
which margin is being protected, and which validation gate turns the
calculation from a screen into evidence.}

\section{The Operator Hierarchy}

The hierarchy begins with the second-order swing equation
\begin{equation}
  M\ddot\theta+D\dot\theta+L\theta=p,
  \label{eq:intro_swing}
\end{equation}
where \(M\) is inertia, \(D\) is damping, \(L\) is synchronizing stiffness, and
\(\theta\) is the vector of bus or machine angles after removing the uniform
reference. In inertia-normalized coordinates, the stiffness becomes
\(\Lt=M^{-1/2}LM^{-1/2}\) and the damping becomes
\(\Dt=M^{-1/2}DM^{-1/2}\).

The first theoretical floor is fixed damping. In inertia-normalized
coordinates,
\begin{equation}
  \Lt=M^{-1/2}LM^{-1/2},\qquad
  \Dt=M^{-1/2}DM^{-1/2}.
  \label{eq:intro_normalized_operators}
\end{equation}
Exact modal decoupling occurs when
\begin{equation}
  [\Lt,\Dt]=\Lt\Dt-\Dt\Lt=0.
  \label{eq:intro_commutation_condition}
\end{equation}
For symmetric matrices, this is the condition for simultaneous orthogonal
diagonalization. The familiar \(D=\gamma M\) case gives
\(\Dt=\gamma I\), so it certainly satisfies
\eqref{eq:intro_commutation_condition}; it is the special case in which every
underdamped mode has the same real decay rate \(-\gamma/2\). Rayleigh damping
\(D=\alpha M+\beta L\) also commutes in normalized coordinates and therefore
decouples graph modes, but it does not produce a uniform decay rate.

The second floor keeps \(D\) fixed but removes proportionality. Now
\(Q^T\Dt Q\) is not diagonal in the graph eigenbasis \(Q\). Off-diagonal modal
damping couples the modes. The eigenvalue problem remains a quadratic
eigenvalue problem (QEP), but graph-frequency movement is no longer a damping
certificate \cite{tisseur2001qep}.

The third floor appears when auxiliary states are eliminated. If the full
linearization is partitioned into retained electromechanical states \(x_e\) and
condensed controller states \(x_c\),
\begin{equation}
  A=
  \begin{bmatrix}
    A_{ee} & A_{ec}\\
    A_{ce} & A_{cc}
  \end{bmatrix},
  \label{eq:intro_partition}
\end{equation}
then Schur condensation gives a nonlinear retained-state operator,
\begin{equation}
  S(s)=sI-A_{ee}-A_{ec}(sI-A_{cc})^{-1}A_{ce}.
  \label{eq:intro_schur}
\end{equation}
In second-order retained coordinates the general condensed operator is
\begin{equation}
  T(s)=s^2M+sD_0+\Lt+\Pi(s),
  \qquad
  \Pi(s)=\Pi_K(s)+s\Pi_D(s).
  \label{eq:intro_sigma}
\end{equation}
The rational correction may contain stiffness-like and damping-like channels.
Only when the stiffness-like part is absent or explicitly absorbed does this
reduce to \(T(s)=s^2M+s\Sigma(s)+\Lt\), with
\(\Sigma(s)=D_0+\Pi_D(s)\). The reduced damping object is then rational in
\(s\), with poles inherited from the eliminated controller block.

The same operator also defines the thesis' action calculus. On the resolvent
set of \(T\), let
\begin{equation}
  G_T(s)=T(s)^{-1}.
  \label{eq:intro_green_function}
\end{equation}
For a line \(e=(i,j)\), with incidence vector \(a_e=e_i-e_j\), define
\begin{equation}
  \mathcal R_e(s)=a_e^TG_T(s)a_e .
  \label{eq:intro_dynamic_resistance}
\end{equation}
At \(s=0\), when the reduced operator has the Laplacian limit and the reference
angle is removed, this recovers classical effective resistance
\(a_e^TL^\dagger a_e\). Away from zero it is a dynamic effective resistance:
it records stiffness, inertia, damping, non-normality, controller poles,
residues, and modal coupling through the same port projection. A finite line
change \(\Delta b_e\) then satisfies
\begin{equation}
  \det(T(s)+\Delta b_ea_ea_e^T)
  =
  \det T(s)\left(1+\Delta b_e\mathcal R_e(s)\right).
  \label{eq:intro_line_action_law}
\end{equation}
Thus a local line sensitivity is only the small-\(\Delta b_e\) limit of an
exact scalar finite-action law.

\begin{figure}[ht]
\centering
\begin{tikzpicture}[
  box/.style={rectangle, rounded corners=2pt, draw=black!65, align=center,
    minimum height=1.06cm, text width=0.72\textwidth},
  level0/.style={box, fill=softblue, draw=thesisblue},
  level1/.style={box, fill=softgreen, draw=thesisgreen},
  level2/.style={box, fill=softorange, draw=thesisorange},
  transition/.style={font=\small, align=center, fill=white, inner sep=2pt},
  arrow/.style={-{Latex[length=2.4mm]}, thick}
]
\node[level0] (l0) {Level 0: commuting fixed damping\\
\([\Lt,\Dt]=0\), including \(D=\gamma M\) and Rayleigh damping};
\node[level1, below=1.15cm of l0] (l1) {Level 1: fixed non-commuting damping\\
\(D\) is constant, but \(Q^T\Dt Q\) couples graph modes};
\node[level2, below=1.15cm of l1] (l2) {Level 2: rational self-energy and Green function\\
\(\Pi(s)=\Pi_K(s)+s\Pi_D(s)\), \(G_T(s)=T(s)^{-1}\)};
\draw[arrow] (l0) -- node[transition,right=0.25cm]{non-commuting damping operator} (l1);
\draw[arrow] (l1) -- node[transition,right=0.25cm]{controller-state condensation} (l2);
\end{tikzpicture}
\caption{Damping-operator hierarchy developed in this thesis. Each transition
keeps part of the previous intuition but removes one mathematical guarantee.}
\label{fig:intro_operator_hierarchy}
\end{figure}

\section{Core Calibrated Messages}

The hierarchy leads to five messages that are strong enough to guide the
thesis, but only if their assumptions are stated precisely.

\begin{enumerate}
\item \textbf{Pure inertia is not damping.} In the fixed-mode, fixed-\(D\)
screen,
\begin{equation}
  \frac{\partial\zeta_k}{\partial M_i}
  \approx
  -\frac{1}{2}\zeta_k\phi_k(i)^2.
  \label{eq:intro_inertia_message}
\end{equation}
Thus adding inertia at a modal antinode can reduce the scalar damping-ratio
screen if damping is not co-designed. This does not say inertia is bad. It
says inertia, damping ratio, RoCoF, and nadir are distinct objectives.

\item \textbf{Graph stiffness is not a damping certificate.} For a line
\(e=(i,j)\),
\begin{equation}
  \frac{\partial\zeta_k}{\partial b_e}
  =
  -\frac{\zeta_k}{2\nu_k}(\phi_k(i)-\phi_k(j))^2
  +
  \frac{\zeta_k}{\delta_k}
  \frac{\partial\delta_k}{\partial b_e}.
  \label{eq:intro_line_message}
\end{equation}
If the line action raises stiffness but does not raise modal damping, the
scalar damping-ratio screen can decrease. Therefore algebraic connectivity and
effective-resistance improvements are graph screens, not automatic
small-signal damping certificates.

\item \textbf{Modal reduction and controller condensation share a resolvent
structure.} Eliminating coupled graph modes from a fixed-\(D\) QEP produces a
rational correction such as \eqref{eq:general_modal_correction}. Eliminating
controller states produces the Schur self-energy
\(\Pi(s)=G(sI-A_{cc})^{-1}(B_\theta+sB_\omega)\). The common mechanism is not
the physical origin of the eliminated states; it is the resolvent injected into
the retained operator. The damping-channel notation \(\Sigma(s)\) is used only
after the rational correction has been projected onto the velocity channel.

\item \textbf{Static damping bridges have a spectral validity region.} In the
damping-channel specialization, the NEP sensitivity denominator contains
\begin{equation}
  T_s(s)=2sM+\Sigma(s)+s\Sigma'(s).
  \label{eq:intro_spectral_memory}
\end{equation}
The term \(s\Sigma'(s)\) is the spectral memory of the condensed controller
block. Setting it to zero is acceptable only when the rational self-energy is
locally flat enough over the critical region. Near a controller resonance, a
static bridge can predict the wrong pole-motion direction.
\item \textbf{Finite actions are not only derivatives.} For line, damping,
inertia, and controller actions that enter \(T(s)\) through low-rank port
updates, the pole condition can be reduced to a determinant whose dimension is
the number of action ports. This does not replace full-model validation, but it
does give the right intermediate object between a scalar graph score and a
large NEP recomputation.
\end{enumerate}

None of these messages is presented as an industry-wide theorem without
qualification. They are calibrated statements: exact in the commuting
constant-\(D\) floor, first-order screens in fixed non-commuting damping, and
validation gates after controller condensation.

\section{Research Questions}

The thesis is organized around eight research questions.

\begin{enumerate}
\item What exact assumptions make graph-mode damping analysis valid in the
classical swing model?
\item How does the operating point induce the operational graph
\(L=B W(z^\star)B^T\), and how do power-flow geometry, line loading, cut
structure, and cycle structure shape dynamic modes?
\item How does fixed non-commuting damping create modal coupling, and how
can that coupling be expressed as a rational correction?
\item How does Schur condensation of controller states transform a fixed
damping matrix into a rational self-energy \(\Pi(s)\), and when is its
damping-channel projection \(\Sigma(s)\) legitimate?
\item How does the dynamic Green function \(G_T(s)\) turn effective
resistance, line reinforcement, damping placement, inertia placement, and
controller retuning into one low-rank action calculus?
\item What theoretical requirements must a planning metric satisfy before it
can claim to identify weak buses, weak lines, or useful converter actions in a
converter-rich grid?
\item How can dynamic hosting capacity, virtual-inertia placement, and
line-reinforcement screening be expressed as spectral margin problems rather
than as static graph or voltage heuristics?
\item Which certificates decide when a reduced QEP, a Schur NEP, a
data-identified self-energy, or a full DAE is adequate for a protected
decision?
\end{enumerate}

The first five questions are mathematical. The last three are epistemic,
planning-oriented, and computational: they define what future empirical
validation must show before a screen becomes an engineering certificate.

\section{Contributions}

The dissertation makes one central contribution and five supporting
contributions.

\begin{enumerate}
\item \textbf{Central contribution: operator-valued dynamic grid strength.}
The dissertation treats graph stiffness, damping-to-inertia roughness,
spectral effective resistance, rational controller self-energy, and nonlinear
memory as successive mechanisms in one retained operator
\(T(s)=s^2M+sD+L+\Pi(s)\).
\item \textbf{Damping-to-inertia graph roughness.} It shows how
non-commutation between inertia-normalized stiffness and damping becomes the
graph-gradient energy of \(D_i/M_i\), and how this roughness, together with
small modal gaps, predicts modal mixing.
\item \textbf{Graph-time resonance atlas.} It defines
\(H_G(\omega)=[Q^T\widetilde T(j\omega)Q]^{-1}\) as the joint spatial-mode
and temporal-frequency response used to locate resonant modal conversion.
\item \textbf{Spectral effective resistance.} It defines
\(\mathcal R_e(s)=a_e^TG_T(s)a_e\) and derives the finite line-update law
\(1+\Delta w_e\mathcal R_e(s)=0\), recovering static effective resistance at
the Laplacian limit.
\item \textbf{Controller-induced reinforcement reversal.} It formalizes the
case in which a line or parameter action passes a static graph or fixed-\(D\)
screen but reverses once the controller Schur resolvent is included.
\item \textbf{Certificate/screen/gate planning methodology.} It formulates
weak elements, dynamic hosting, virtual-inertia placement, and line
reinforcement as operator-, margin-, action-, and validation-gate-relative
claims rather than graph-intrinsic labels.
\end{enumerate}

\section{Dissertation Organization}

The dissertation is now organized as a modular monograph. Part~I builds the
problem statement, the physical grid framework, the spectral preliminaries, and
the state of the art. The new physical framework chapter develops the AC
network, phasor, power-flow, swing-equation, DAE, and operator chain before the
more compact mathematical chapters use those objects. Part~II develops the
operational graph theory: power-flow geometry, total derivatives of
\(L=B W(z^\star)B^T\), graph roughness of damping-to-inertia fields,
graph-time response, dynamic effective resistance, and structured fragility
radii. Part~III closes the constant-\(D\) theory: QEP identities, commuting
damping, uniform decay, Rayleigh damping, non-commuting coupling, and
reduced-QEP gates. Part~IV develops rational self-energy \(\Pi(s)\), its
damping-channel specialization \(\Sigma(s)\), Schur condensation,
finite-action laws, NEP sensitivities, port equivalence, passivity screens,
structured fragility, and controller-resolvent reversals. Part~V separates
genuine nonlinear DAE stability from spectral NEP theory by treating
equilibria, nonlinear memory, bifurcation, energy, and regions of attraction.
Part~VI turns the operators into applications: weak elements, planning risk
vectors, dynamic hosting regions, reinforcement envelopes, and reproducible
benchmark protocols.

The integrated source-to-claim ledger is not part of the main scientific
progression. It is kept as an appendix-level traceability record that maps
claims to sources, scripts, figures, and validation gates.

\chapter{Physical and Theoretical Framework}

This chapter explains the physical model behind the operator hierarchy. The
later chapters are intentionally abstract, but the abstraction is not detached
from the grid. Every matrix in the thesis comes from a physical layer: phasor
voltages, line admittances, active-power balance, rotor or converter dynamics,
controller states, and the operating point around which the system is
linearized. The purpose of this chapter is to make that chain explicit before
the thesis uses \(L\), \(M\), \(D\), \(\Pi(s)\), and \(T(s)^{-1}\) as technical
objects.

\thesisbox{thesisblue}{Chapter message}{
The problem is not only to compute eigenvalues. The problem is to identify
which physical mechanism creates, moves, amplifies, or protects those
eigenvalues. That is why the thesis needs a hierarchy from network stiffness to
fixed damping, rational self-energy, Green functions, structured fragility, and
nonlinear memory.}

\section{How to Read This Chapter Step by Step}

The rest of the thesis uses compact notation, so this chapter starts from the
physical chain and builds the operator one piece at a time. The intended
reading order is:
\begin{enumerate}
\item Start with the electrical graph: buses are nodes, lines are edges, and
line admittances create the network equation \(I=Y_{\rm bus}V\).
\item Fix an operating point: the same topology can have different
small-signal stiffness under a different dispatch, voltage profile, or angle
stress.
\item Linearize active-power balance: the derivative of power with respect to
angles becomes the synchronizing stiffness matrix \(L\).
\item Add dynamic storage and damping: inertia gives \(M\ddot\theta\), damping
or damping-like control gives \(D\dot\theta\), and stiffness gives
\(L\theta\).
\item Introduce the spectral parameter \(s\): an oscillatory response
\(\theta(t)=ve^{st}\) converts the differential equation into an operator
\(s^2M+sD+L\).
\item Escalate only when needed: if hidden controller states are eliminated,
the missing frequency-dependent physics appears as \(\Pi(s)\), and the object
becomes \(T(s)=s^2M+sD_0+L+\Pi(s)\).
\end{enumerate}
The key point is that each symbol is a physical commitment. A graph eigenvalue
is not automatically a stability certificate; it is a certificate only when
the dynamic terms that can change the protected margin have been accounted for.

\section{The Grid as a Networked Dynamic System}

A power grid is a networked dynamical system, not a single aggregate machine.
The buses form the electrical graph, the lines and transformers define the
admittance matrix, the generators and converters inject dynamic behavior, and
the loads close the power balance. A useful small-signal model must preserve
all four roles: network geometry, operating-point physics, dynamic storage, and
control memory.

\begin{figure}[ht]
\centering
\resizebox{0.96\textwidth}{!}{%
\begin{tikzpicture}[
  bus/.style={circle,draw=thesisink!65,fill=white,line width=0.45pt,minimum size=6pt,inner sep=0pt},
  gen/.style={circle,draw=thesisblue!85!black,fill=thesisblue!10,line width=0.9pt,minimum size=8pt,inner sep=0pt},
  conv/.style={regular polygon,regular polygon sides=4,draw=thesisorange!85!black,fill=thesisorange!12,line width=0.9pt,minimum size=9pt,inner sep=0pt},
  load/.style={rectangle,draw=thesisgray!85!black,fill=softgray,line width=0.6pt,minimum width=8pt,minimum height=6pt,inner sep=0pt},
  line/.style={draw=thesisgray!76,line width=0.7pt},
  ctrl/.style={rectangle,rounded corners=2pt,draw=thesisred!75!black,fill=softred,align=center,font=\scriptsize,text width=2.25cm},
  lbl/.style={font=\scriptsize\color{thesisink}},
  flow/.style={-{Latex[length=2.0mm]},draw=thesisgray,line width=0.55pt}
]
\coordinate (b1) at (0,1.6);
\coordinate (b2) at (1.7,2.2);
\coordinate (b3) at (3.4,1.55);
\coordinate (b4) at (5.0,2.15);
\coordinate (b5) at (6.5,1.25);
\coordinate (b6) at (1.1,0.25);
\coordinate (b7) at (2.8,-0.35);
\coordinate (b8) at (4.8,0.05);
\coordinate (b9) at (6.2,-0.55);
\draw[line] (b1)--(b2)--(b3)--(b4)--(b5)--(b9)--(b8)--(b7)--(b6)--(b1);
\draw[line] (b2)--(b6) (b2)--(b7) (b3)--(b7) (b3)--(b8) (b4)--(b8) (b5)--(b8);
\foreach \p in {b1,b4,b6,b8} \node[bus] at (\p) {};
\node[gen] (g1) at (b2) {};
\node[conv] (c1) at (b3) {};
\node[gen] (g2) at (b5) {};
\node[conv] (c2) at (b7) {};
\node[load] (l1) at (b9) {};
\node[lbl,above=0.08cm of g1] {SG};
\node[lbl,above=0.08cm of c1] {IBR};
\node[lbl,above=0.08cm of g2] {SG};
\node[lbl,below=0.10cm of c2] {GFM/GFL};
\node[lbl,below=0.08cm of l1] {load};
\node[ctrl] (pll) at (3.4,3.20) {PLL/current loop\\voltage loop\\plant control};
\node[ctrl] (gov) at (0.95,3.20) {governor\\exciter\\stabilizer};
\draw[flow] (gov.south) -- (g1.north);
\draw[flow] (pll.south) -- (c1.north);
\node[lbl,align=center,text width=4.2cm] at (9.1,1.80) {\textbf{Electrical graph}\\
\(Y_{\rm bus}=G+jB\)\\
\(I=Y_{\rm bus}V\)};
\node[lbl,align=center,text width=4.2cm] at (9.1,0.60) {\textbf{Dynamic devices}\\
\(M\ddot\theta+D\dot\theta\)\\
or converter states \(z\)};
\node[lbl,align=center,text width=4.2cm] at (9.1,-0.65) {\textbf{Planning actions}\\
line, damping, inertia,\\
controller retuning};
\draw[flow] (6.95,1.25) -- (7.65,1.62);
\draw[flow] (6.95,0.18) -- (7.65,0.48);
\draw[flow] (6.95,-0.48) -- (7.65,-0.62);
\end{tikzpicture}}
\caption{A converter-rich grid contains a physical graph, dynamic devices, and
hidden controller states. The retained variables may be bus or machine angles,
frequencies, voltages, and selected ports; the eliminated variables reappear as
self-energy when they are condensed.}
\label{fig:framework_networked_grid}
\end{figure}

Figure~\ref{fig:framework_networked_grid} shows the modeling issue. The
electrical graph tells us which buses are connected, but not which dynamic
margin is protected. A synchronous generator contributes inertia, damping,
exciter, governor, and stabilizer dynamics. A grid-following or grid-forming
converter contributes synchronization, current-loop, voltage-loop, filter,
limiter, and plant-level dynamics. If those internal states are removed from a
reduced model, their effect must go somewhere. In this thesis, that place is
the rational self-energy \(\Pi(s)\).

\section{AC Phasors, Power Flow, and Network Stiffness}

The first physical object is the AC phasor voltage. At bus \(i\),
\begin{equation}
  V_i=|V_i|e^{j\theta_i},
  \label{eq:framework_phasor}
\end{equation}
where \(|V_i|\) is voltage magnitude and \(\theta_i\) is phase angle. The
network admittance matrix maps bus voltages to currents,
\begin{equation}
  I=Y_{\rm bus}V,\qquad Y_{\rm bus}=G+jB,
  \label{eq:framework_ybus}
\end{equation}
and the complex power injected at the bus is
\begin{equation}
  S_i=P_i+jQ_i=V_iI_i^* .
  \label{eq:framework_complex_power}
\end{equation}
The operating point is therefore not only a graph. It is a solution of the AC
power-flow equations at a particular dispatch, load, voltage profile, topology,
and controller setting.

The active-power stiffness matrix is the first derivative of active power with
respect to phase angles at the operating point. With voltage magnitudes fixed
inside the small-signal angle block,
\begin{equation}
  \Delta P\approx J_{P\theta}\Delta\theta,
  \qquad
  J_{P\theta}
  =
  \left.
  \frac{\partial P}{\partial \theta}
  \right|_{(V^\star,\theta^\star)} .
  \label{eq:framework_jacobian}
\end{equation}
Depending on the sign convention for exported electrical power, the
synchronizing matrix is \(L=J_{P\theta}\) or \(L=-J_{P\theta}\). The thesis
always states the convention before using a derivative. The important fact is
that \(L\) is an operating-point derivative, not a purely topological matrix.

\begin{figure}[ht]
\centering
\resizebox{0.92\textwidth}{!}{%
\begin{tikzpicture}[
  axis/.style={-{Latex[length=2.0mm]},draw=thesisink!70,line width=0.55pt},
  curve/.style={draw=thesisblue,line width=1.05pt},
  tangent/.style={draw=thesisorange,line width=0.85pt},
  guide/.style={draw=thesisgray!70,densely dashed,line width=0.45pt},
  lbl/.style={font=\small\color{thesisink}}
]
\draw[axis] (-3.55,0) -- (3.65,0) node[right,lbl] {\(\delta=\theta_i-\theta_j\)};
\draw[axis] (0,-1.75) -- (0,1.95) node[above,lbl] {\(P_{ij}\)};
\draw[curve] plot[smooth] coordinates {
  (-3.14,0) (-2.62,-0.50) (-2.09,-0.87) (-1.57,-1.00)
  (-1.05,-0.87) (-0.52,-0.50) (0,0) (0.52,0.50)
  (1.05,0.87) (1.57,1.00) (2.09,0.87) (2.62,0.50) (3.14,0)
};
\coordinate (op) at (0.78,0.70);
\draw[guide] (op) -- (0.78,0) node[below,lbl] {\(\delta^\star\)};
\draw[guide] (op) -- (0,0.70);
\draw[tangent] (-0.35,-0.26) -- (1.85,1.66);
\node[circle,fill=thesisred,inner sep=1.7pt] at (op) {};
\node[lbl,anchor=west] at (1.00,0.44) {operating point};
\node[lbl,align=center,text width=4.8cm] at (-2.05,1.32)
{\(\displaystyle P_{ij}\approx |V_i||V_j|B_{ij}\sin\delta\)};
\node[lbl,align=center,text width=5.3cm] at (2.05,-1.25)
{\(\displaystyle \frac{\partial P_{ij}}{\partial \delta}
=|V_i^\star||V_j^\star|B_{ij}\cos\delta^\star\)};
\end{tikzpicture}}
\caption{A line contributes synchronizing stiffness because active power is
sensitive to angle separation. The slope at the operating point becomes the
edge weight in the linearized stiffness matrix \(L\).}
\label{fig:framework_power_angle_plot}
\end{figure}

Figure~\ref{fig:framework_power_angle_plot} gives the physical reason \(L\)
acts like a stiffness matrix. For a lossless two-bus line,
\begin{equation}
  P_{ij}\approx |V_i||V_j|B_{ij}\sin(\theta_i-\theta_j).
  \label{eq:framework_line_power}
\end{equation}
Linearization gives a slope proportional to
\(|V_i^\star||V_j^\star|B_{ij}\cos(\theta_i^\star-\theta_j^\star)\). This
slope is large when the line is electrically strong and the operating point is
not close to the angular stability limit. It decreases as the operating point
approaches a stressed angle difference. Thus the same physical line can have
different small-signal stiffness under different operating conditions.

\section{From Physical Balance to the Swing Equation}

The second physical object is rotational or emulated energy storage. For a
synchronous machine, the rotor accelerates when mechanical input and
electrical output are imbalanced:
\begin{equation}
  M_i\dot\omega_i=P_{m,i}-P_{e,i}-D_i\omega_i,
  \qquad
  \omega_i=\dot\theta_i .
  \label{eq:framework_single_swing}
\end{equation}
After linearization around a balanced operating point,
\begin{equation}
  M_i\Delta\ddot\theta_i+D_i\Delta\dot\theta_i+\Delta P_{e,i}=p_i .
  \label{eq:framework_single_linear_swing}
\end{equation}
The network supplies \(\Delta P_e\approx L\Delta\theta\). Stacking retained
angle variables gives
\begin{equation}
  M\ddot\theta+D\dot\theta+L\theta=p .
  \label{eq:framework_network_swing}
\end{equation}
The three matrices are not interchangeable. \(M\) multiplies acceleration,
\(D\) multiplies velocity, and \(L\) multiplies angle displacement. A planning
action that changes one term should not be interpreted as if it changed the
others.

The uniform angle direction has no restoring force. Mathematically,
\begin{equation}
  L\ones=0 .
  \label{eq:framework_reference_invariance}
\end{equation}
Physically, adding the same constant angle to every bus does not change any
angle difference or power flow. All spectral statements therefore occur after
fixing an angle reference, using center-of-inertia coordinates, or projecting
onto the nonzero-angle subspace.

\section{Converter States and the DAE Layer}

Converter-rich grids add a hidden dynamic layer between the electrical graph
and the retained swing-like variables. A generic local model has differential
states \(x\), algebraic variables \(y\), and parameters or actions \(\rho\):
\begin{align}
  \dot x &= f(x,y,\rho), \label{eq:framework_dae_f}\\
  0 &= g(x,y,\rho). \label{eq:framework_dae_g}
\end{align}
The algebraic variables include network constraints such as nodal current
balance and voltage equations. The differential states include machine rotor
states, exciter states, governor states, PLL states, current controllers,
voltage controllers, filters, and plant controls.

At a regular equilibrium, \(g_y\) is nonsingular and the algebraic variables
can be eliminated locally. The first-order small-signal matrix is
\begin{equation}
  A_{\rm red}=f_x-f_y g_y^{-1}g_x .
  \label{eq:framework_reduced_A}
\end{equation}
This formula explains why the operating point matters. The Jacobians \(f_x\),
\(f_y\), \(g_x\), and \(g_y\) are all evaluated at
\((x^\star,y^\star,\rho^\star)\). If the dispatch, IBR penetration, line
status, voltage setpoints, or controller gains change, the linearized operator
changes.

\begin{figure}[ht]
\centering
\resizebox{\textwidth}{!}{%
\begin{tikzpicture}[
  layer/.style={rectangle,rounded corners=3pt,draw=#1!75!black,fill=#1!9,
    minimum height=1.0cm,text width=2.85cm,align=center,font=\small},
  arrow/.style={-{Latex[length=2.2mm]},draw=thesisgray,line width=0.7pt},
  note/.style={font=\scriptsize\color{thesisgray},align=center,text width=2.8cm}
]
\node[layer=thesisblue] (net) {Electrical network\\\(Y_{\rm bus},V,\theta\)};
\node[layer=thesisgreen, right=0.55cm of net] (opf) {Operating point\\\(P,Q,V^\star,\theta^\star\)};
\node[layer=thesisorange, right=0.55cm of opf] (jac) {Linearization\\\(L,J_{xy},A_{\rm red}\)};
\node[layer=thesisred, right=0.55cm of jac] (dyn) {Dynamic operator\\\(T(s)\)};
\node[layer=thesisteal, right=0.55cm of dyn] (dec) {Decision\\margin, action,\\certificate};
\draw[arrow] (net) -- (opf);
\draw[arrow] (opf) -- (jac);
\draw[arrow] (jac) -- (dyn);
\draw[arrow] (dyn) -- (dec);
\node[note, below=0.25cm of net] {topology and admittance};
\node[note, below=0.25cm of opf] {dispatch and constraints};
\node[note, below=0.25cm of jac] {small-signal physics};
\node[note, below=0.25cm of dyn] {QEP or Schur NEP};
\node[note, below=0.25cm of dec] {hosting or retuning};
\end{tikzpicture}}
\caption{The thesis follows the physical modeling chain from network data to
an operator-valued decision. A planning result is not reproducible unless the
operating point, retained variables, action family, margin, and certification
gate are all declared.}
\label{fig:framework_modeling_pipeline}
\end{figure}

\section{Why the Spectral Parameter Appears}

Small-signal stability studies the free response after an infinitesimal
disturbance. For the fixed-\(D\) second-order model, set \(p=0\) and try
\(\theta(t)=ve^{st}\). Since
\[
  \dot\theta=sve^{st},
  \qquad
  \ddot\theta=s^2ve^{st},
\]
substitution into \eqref{eq:framework_network_swing} gives
\begin{equation}
  (s^2M+sD+L)v=0.
  \label{eq:framework_qep}
\end{equation}
This quadratic eigenvalue problem is the first spectral object of the thesis.
The term \(s^2M\) is inertia, the term \(sD\) is velocity damping, and the
term \(L\) is synchronizing stiffness.

When controller states are eliminated rather than retained, the operator is no
longer a fixed-coefficient QEP. Partition the linearized dynamics into retained
states \(x_r\) and condensed states \(x_c\). Eliminating \(x_c\) gives a Schur
complement with the resolvent \((sI-A_{cc})^{-1}\). In second-order retained
coordinates this appears as
\begin{equation}
  T(s)=s^2M+sD_0+L+\Pi(s),
  \qquad
  \Pi(s)=\Pi_K(s)+s\Pi_D(s).
  \label{eq:framework_nep}
\end{equation}
The rational term \(\Pi(s)\) is the spectral memory of the eliminated
controller block. It contains poles and residues. A static damping matrix
cannot exactly represent those poles.

\section{What a Pole Means Physically}

A root \(s=\sigma+j\omega\) of the retained operator produces a modal response
of the form
\[
  e^{st}=e^{\sigma t}e^{j\omega t}.
\]
The real part \(\sigma\) controls growth or decay, and the imaginary part
\(\omega\) controls oscillation frequency. The damping ratio is
\begin{equation}
  \zeta=-\frac{\sigma}{\sqrt{\sigma^2+\omega^2}} .
  \label{eq:framework_damping_ratio}
\end{equation}
The same pole can be interpreted in several physical ways: a slow
electromechanical oscillation, a local plant mode, a controller interaction
mode, or a mixed network-control mode. This is why the thesis does not name a
mode only by frequency. It asks which operator term and which port direction
make the mode visible.

\begin{figure}[ht]
\centering
\resizebox{0.86\textwidth}{!}{%
\begin{tikzpicture}[
  axis/.style={-{Latex[length=2.0mm]},draw=thesisink!75,line width=0.55pt},
  stable/.style={fill=softgreen,draw=none,opacity=0.55},
  protect/.style={draw=thesisorange,line width=0.85pt,densely dashed},
  pole/.style={circle,fill=thesisblue,inner sep=1.7pt},
  badpole/.style={circle,fill=thesisred,inner sep=1.8pt},
  arrow/.style={-{Latex[length=2.1mm]},draw=thesisred,line width=0.8pt},
  lbl/.style={font=\small\color{thesisink}}
]
\fill[stable] (-4.0,-2.0) rectangle (0,2.25);
\draw[axis] (-4.05,0) -- (1.45,0) node[right,lbl] {\(\Real(s)\)};
\draw[axis] (0,-2.05) -- (0,2.35) node[above,lbl] {\(\Imag(s)\)};
\draw[protect] (-1.0,-2.0) -- (-1.0,2.15);
\node[lbl,anchor=east] at (-1.05,1.92) {protected boundary};
\node[pole] (p1) at (-2.15,1.18) {};
\node[pole] (p2) at (-2.15,-1.18) {};
\node[pole] (p3) at (-0.58,0.78) {};
\node[badpole] (p4) at (0.18,0.72) {};
\draw[arrow] (p3) -- (p4);
\node[lbl,align=center,text width=3.2cm] at (-2.75,-1.70) {stable and protected};
\node[lbl,align=center,text width=3.4cm] at (0.52,1.38) {action or uncertainty can move a pole};
\end{tikzpicture}}
\caption{The spectral plane makes the margin explicit. A nominally stable pole
can still be fragile if a physically admissible perturbation moves it across a
protected damping, abscissa, resonance, or passivity boundary.}
\label{fig:framework_spectral_plane}
\end{figure}

Figure~\ref{fig:framework_spectral_plane} explains why a nominal eigenvalue
plot is not enough. A system may be stable at the base point but fragile under
small line uncertainty, controller retuning, model-identification error, or
operating-point movement. The thesis therefore uses structured fragility
rather than only nominal roots.

\section{Why Scalar Grid Strength Is Not Enough}

Classical strength metrics are useful because they compress a complicated
network into a small number. Short-circuit ratio describes local electrical
stiffness. Algebraic connectivity describes a global graph bottleneck. A
damping ratio describes one pole. Effective resistance describes a static
graph compliance. None of these scalars is wrong. The problem is that each one
protects only one mechanism.

\begin{table}[ht]
\centering
\caption{Why a scalar metric is not a complete dynamic-strength theory.}
\label{tab:framework_scalar_limits}
\begin{tabular}{L{0.20\textwidth}L{0.28\textwidth}L{0.34\textwidth}}
\toprule
\rowcolor{softblue}
\thead{Scalar} & \thead{What it sees} & \thead{What it can miss}\\
\midrule
SCR & local fault-level or Thevenin stiffness & damping, controller poles,
non-normal amplification, multi-port interactions\\
\(\lambda_2(L)\) & global graph connectivity & heterogeneous damping,
inertia placement, operating-point residues\\
\(\zeta_k\) & decay of one pole & finite-action fragility, nearby controller
poles, hidden input-output gain\\
\(a_e^TL^\dagger a_e\) & static effective resistance & frequency-dependent
compliance, resonant Green-function peaks\\
\bottomrule
\end{tabular}
\end{table}

The dynamic replacement is not another scalar. It is the retained operator
\begin{equation}
  T(s)=s^2M+sD_0+L+\Pi(s)
  \label{eq:framework_core_operator}
\end{equation}
and its Green function
\begin{equation}
  G_T(s)=T(s)^{-1}.
  \label{eq:framework_green}
\end{equation}
The scalar quantities used for planning are projections of this operator onto a
declared port, action family, frequency band, or protected boundary.

\begin{figure}[ht]
\centering
\resizebox{0.86\textwidth}{!}{%
\begin{tikzpicture}[
  axis/.style={-{Latex[length=2.0mm]},draw=thesisink!75,line width=0.55pt},
  curveA/.style={draw=thesisblue,line width=1.0pt},
  curveB/.style={draw=thesisred,line width=1.0pt},
  curveC/.style={draw=thesisorange,line width=1.0pt,densely dashed},
  lbl/.style={font=\small\color{thesisink}},
  leg/.style={font=\scriptsize\color{thesisink},anchor=west}
]
\draw[axis] (0,0) -- (7.55,0) node[right,lbl] {line stiffness \(b_e\)};
\draw[axis] (0,0) -- (0,3.55) node[above,lbl] {normalized margin};
\draw[curveA] plot[smooth] coordinates {(0.45,0.55) (1.15,1.02) (1.95,1.55) (2.85,2.04) (3.75,2.38) (4.70,2.66) (5.55,2.84)};
\draw[curveB] plot[smooth] coordinates {(0.45,2.80) (1.15,2.60) (1.95,2.30) (2.85,1.88) (3.75,1.38) (4.70,0.92) (5.55,0.62)};
\draw[curveC] plot[smooth] coordinates {(0.45,2.45) (1.15,2.42) (1.95,2.26) (2.85,1.90) (3.75,1.20) (4.70,0.78) (5.55,1.02)};
\draw[curveA] (5.95,2.70) -- (6.45,2.70);
\node[leg] at (6.55,2.70) {graph stiffness};
\draw[curveB] (5.95,2.33) -- (6.45,2.33);
\node[leg] at (6.55,2.33) {damping headroom};
\draw[curveC] (5.95,1.96) -- (6.45,1.96);
\node[leg] at (6.55,1.96) {controller-aware margin};
\end{tikzpicture}}
\caption{A stiffness-improving line action need not improve a dynamic margin.
The graph metric, fixed-\(D\) damping screen, and controller-aware margin can
move in different directions.}
\label{fig:framework_scalar_failure_plot}
\end{figure}

Figure~\ref{fig:framework_scalar_failure_plot} is the central planning
motivation. If a line upgrade increases graph stiffness but reduces damping
headroom, a scalar graph score gives the wrong engineering message. If the
static screen accepts the line but the controller-aware model rejects it, the
missing mechanism is not topology. It is controller memory.

\section{A Three-Bus Example Where Classical Screens Fail}

This section gives the smallest possible example that contains the dissertation
problem. The example is not meant to represent a full test system. It is a
teaching model that shows why a line action that is excellent for a classical
graph metric can be bad for a dynamic margin.

\subsection{Step 1: Build the Three-Bus Graph}

Consider the three-bus chain in Figure~\ref{fig:framework_three_bus}. Bus 1 is
a synchronous source with appreciable damping, bus 2 is a passive transfer
bus, and bus 3 is a converter-dominated bus with weak damping. The base system
has two lines, \(1\)--\(2\) and \(2\)--\(3\), both with unit synchronizing
weight. The candidate planning action is a new tie line \(1\)--\(3\) with
incremental stiffness \(p\ge 0\).

\begin{figure}[ht]
\centering
\resizebox{0.90\textwidth}{!}{%
\begin{tikzpicture}[
  bus/.style={circle,draw=thesisink!70,fill=white,line width=0.7pt,
    minimum size=10mm,align=center,font=\small},
  gen/.style={circle,draw=thesisblue!85!black,fill=thesisblue!12,line width=1.0pt,
    minimum size=11mm,align=center,font=\small},
  conv/.style={regular polygon,regular polygon sides=4,draw=thesisorange!85!black,
    fill=thesisorange!13,line width=1.0pt,minimum size=11mm,align=center,font=\small},
  line/.style={draw=thesisgray!78,line width=1.2pt},
  cand/.style={draw=thesisred!82,line width=1.0pt,densely dashed},
  lbl/.style={font=\small\color{thesisink},align=center},
  note/.style={rectangle,rounded corners=2pt,draw=thesisgray!45,fill=softgray,
    align=center,font=\scriptsize,text width=3.3cm}
]
\node[gen] (b1) at (0,0) {1};
\node[bus] (b2) at (3.8,0) {2};
\node[conv,rotate=45] (b3) at (7.6,0) {\rotatebox{-45}{3}};
\draw[line] (b1) -- node[above,lbl] {\(b_{12}=1\)} (b2);
\draw[line] (b2) -- node[above,lbl] {\(b_{23}=1\)} (b3);
\draw[cand] (b1) to[bend right=60] node[pos=0.50,below=0.70cm,lbl] {candidate tie \(p\)} (b3);
\node[lbl,below=0.36cm of b1] {SG\\\(d_1=0.60\)};
\node[lbl,below=0.36cm of b2] {transfer bus\\\(d_2=0.60\)};
\node[lbl,below=0.36cm of b3] {IBR bus\\\(d_3=0.02\)};
\node[note] at (1.45,1.60) {classical graph screen:\\increase connectivity};
\node[note] at (6.05,1.60) {dynamic screen:\\protect damping and resonance};
\end{tikzpicture}}
\caption{Three-bus teaching example. The candidate tie \(1\)--\(3\) improves
the static graph, but it also changes the frequency of the endpoint mode that
uses the weakly damped converter bus.}
\label{fig:framework_three_bus}
\end{figure}

\subsection{Step 2: Derive the Stiffness Matrix}

With only the chain lines present, the Laplacian stiffness is
\begin{equation}
  L_0=
  \begin{bmatrix}
  1 & -1 & 0\\
  -1 & 2 & -1\\
  0 & -1 & 1
  \end{bmatrix}.
  \label{eq:framework_three_bus_L0}
\end{equation}
The nonzero graph modes are
\begin{equation}
  \nu_1=1,\quad
  \phi_1=\frac{1}{\sqrt 2}
  \begin{bmatrix}1\\0\\-1\end{bmatrix},
  \qquad
  \nu_2=3,\quad
  \phi_2=\frac{1}{\sqrt 6}
  \begin{bmatrix}1\\-2\\1\end{bmatrix}.
  \label{eq:framework_three_bus_modes}
\end{equation}
The candidate tie has incidence vector
\begin{equation}
  a_{13}=e_1-e_3=
  \begin{bmatrix}1\\0\\-1\end{bmatrix},
  \qquad
  L(p)=L_0+p\,a_{13}a_{13}^T .
  \label{eq:framework_three_bus_action}
\end{equation}
Because \(a_{13}\) is aligned with \(\phi_1\), the tie line changes the first
nonzero stiffness exactly:
\begin{equation}
  \nu_1(p)=1+2p,\qquad \nu_2(p)=3 .
  \label{eq:framework_three_bus_eigs}
\end{equation}
For \(0\le p\le 1\), the algebraic connectivity therefore increases from
\(1\) to \(3\). The static effective resistance between buses 1 and 3 also
decreases:
\begin{equation}
  R_{13}^{\rm stat}(p)=\frac{2}{1+2p}.
  \label{eq:framework_three_bus_resistance}
\end{equation}
A classical graph-only decision would say that increasing \(p\) is
unambiguously good: the network becomes more connected and the endpoint
resistance becomes smaller.

\subsection{Step 3: Add the First Dynamic Screen}

Now add unit inertia and heterogeneous damping,
\begin{equation}
  M=I,\qquad
  D=\operatorname{diag}(0.60,0.60,0.02).
  \label{eq:framework_three_bus_MD}
\end{equation}
The endpoint mode \(\phi_1\) moves bus 1 against bus 3 while bus 2 stays nearly
still. Its frozen graph-mode damping estimate is
\begin{equation}
  \delta_1=\phi_1^TD\phi_1
  =\frac{d_1+d_3}{2}
  =0.31 .
  \label{eq:framework_three_bus_modal_damping}
\end{equation}
The same line that increases stiffness also raises the modal frequency. A
first dynamic damping-ratio screen gives
\begin{equation}
  \zeta_1(p)\approx
  \frac{\delta_1}{2\sqrt{\nu_1(p)}}
  =
  \frac{0.31}{2\sqrt{1+2p}} .
  \label{eq:framework_three_bus_zeta}
\end{equation}
The sign of the conclusion changes. The graph metric improves with \(p\), but
the damping ratio falls with \(p\). If the planning rule requires
\(\zeta\ge 0.10\), the base case passes and the strong tie \(p=1\) fails.

\begin{figure}[ht]
\centering
\resizebox{0.82\textwidth}{!}{%
\begin{tikzpicture}[
  axis/.style={-{Latex[length=2.0mm]},draw=thesisink!75,line width=0.55pt},
  grid/.style={draw=thesisgray!24,line width=0.35pt},
  stiff/.style={draw=thesisblue,line width=1.1pt},
  damp/.style={draw=thesisred,line width=1.1pt},
  ctrl/.style={draw=thesisorange,line width=1.0pt,densely dashed},
  thresh/.style={draw=thesisink!55,line width=0.6pt,densely dotted},
  lbl/.style={font=\small\color{thesisink}},
  smalllbl/.style={font=\scriptsize\color{thesisink}},
  leg/.style={font=\scriptsize\color{thesisink},anchor=west}
]
\foreach \x in {0.70,1.95,3.20,4.45,5.70} \draw[grid] (\x,0.70) -- (\x,3.35);
\foreach \y in {0.95,1.50,2.05,2.60,3.15} \draw[grid] (0.55,\y) -- (5.95,\y);
\draw[axis] (0.55,0.70) -- (6.25,0.70) node[right,lbl] {tie strength \(p\)};
\draw[axis] (0.55,0.70) -- (0.55,3.55) node[above,lbl] {normalized score};
\node[smalllbl] at (0.70,0.43) {0};
\node[smalllbl] at (3.20,0.43) {0.5};
\node[smalllbl] at (5.70,0.43) {1};
\draw[stiff] plot[smooth] coordinates {(0.70,1.05) (1.95,1.52) (3.20,2.03) (4.45,2.55) (5.70,3.05)};
\draw[damp] plot[smooth] coordinates {(0.70,2.92) (1.95,2.42) (3.20,1.85) (4.45,1.42) (5.70,1.18)};
\draw[ctrl] plot[smooth] coordinates {(0.70,2.65) (1.95,2.06) (3.20,1.33) (4.45,0.88) (5.70,0.80)};
\draw[thresh] (0.55,1.33) -- (5.95,1.33);
\node[smalllbl,anchor=east] at (0.46,1.33) {requirement};
\draw[stiff] (0.90,0.08) -- (1.32,0.08);
\node[leg] at (1.42,0.08) {\(\lambda_2(L(p))\)};
\draw[damp] (2.80,0.08) -- (3.22,0.08);
\node[leg] at (3.32,0.08) {\(\zeta_1(p)\)};
\draw[ctrl] (4.45,0.08) -- (4.87,0.08);
\node[leg] at (4.97,0.08) {controller distance};
\end{tikzpicture}}
\caption{In the three-bus example, the classical graph score and the dynamic
margin move in opposite directions. This is the simplest visible failure of a
connectivity-only planning rule.}
\label{fig:framework_three_bus_plot}
\end{figure}

\subsection{Step 4: Show the Classical Failure}

Table~\ref{tab:framework_three_bus_failure} summarizes the contradiction. The
line action is attractive under two classical static metrics:
\(\lambda_2(L(p))\) rises and \(R_{13}^{\rm stat}(p)\) falls. The same action
is unattractive under the dynamic damping screen because it raises the
frequency of a mode that uses the weakly damped converter bus.

\begin{table}[ht]
\centering
\small
\caption{Three-bus example: static improvement but dynamic degradation.}
\label{tab:framework_three_bus_failure}
\begin{tabular}{c c c c c L{0.23\textwidth}}
\toprule
\rowcolor{softblue}
\thead{\(p\)} & \thead{\(\lambda_2\)} & \thead{\(R_{13}^{\rm stat}\)}
& \thead{\(\omega_1\approx\sqrt{1+2p}\)}
& \thead{\(\zeta_1\)} & \thead{Decision if \(\zeta_{\rm req}=0.10\)}\\
\midrule
0 & 1.00 & 2.00 & 1.00 & 0.155 & passes the damping screen\\
0.5 & 2.00 & 1.00 & 1.41 & 0.110 & barely passes\\
1.0 & 3.00 & 0.67 & 1.73 & 0.089 & fails even though graph strength improved\\
\bottomrule
\end{tabular}
\end{table}

This is the basic reason the thesis cannot stop at graph strength. A graph
method answers, ``Did the network become stiffer?'' The stability question is
different: ``Did the protected dynamic margin improve?'' In this example the
answer to the first question is yes, while the answer to the second is no.

\subsection{Step 5: Add the Hidden Controller Mechanism}

The previous step used only \(M\), \(D\), and \(L(p)\). Converter-rich systems
can fail one level higher. Suppose the eliminated converter/controller block at
bus 3 has a lightly damped internal pole near
\begin{equation}
  \mu_c=-0.05+j1.70 .
  \label{eq:framework_three_bus_controller_pole}
\end{equation}
A static graph screen cannot see this pole. A fixed-\(D\) screen cannot
represent it exactly either, because the reduced effect is frequency
dependent. After Schur condensation, it appears inside \(\Pi(s)\), and the
mode must be tested through
\begin{equation}
  T(s;p)=s^2I+sD+L(p)+\Pi(s).
  \label{eq:framework_three_bus_nep}
\end{equation}
A simple resonance-distance gate is
\begin{equation}
  d_c(p)=
  \frac{|j\sqrt{1+2p}-\mu_c|}
       {|j\sqrt{1+2p}|+|\mu_c|}.
  \label{eq:framework_three_bus_controller_distance}
\end{equation}
For \(p=0\), \(d_c\approx0.260\). For \(p=0.5\), \(d_c\approx0.093\). For
\(p=1\), \(d_c\approx0.017\). Thus the same tie that the classical graph
method ranks best also moves the network mode toward the hidden controller
pole. The correct escalation is not cosmetic; it changes the engineering
decision. The controller-aware analysis asks whether
\(\Pi(s)\), residues, port direction, and dynamic effective resistance create a
fragile interaction near the protected frequency band.

\subsection{What the Example Teaches}

The three-bus example separates three statements that are often conflated:
\begin{enumerate}
\item \textbf{Topology statement:} adding the tie \(1\)--\(3\) increases
\(\lambda_2\) and lowers static resistance.
\item \textbf{Fixed-dynamics statement:} the same tie can lower damping ratio
because it raises frequency without adding proportional damping.
\item \textbf{Controller-aware statement:} if a hidden controller pole sits in
the new frequency band, the reduced operator must include \(\Pi(s)\), and the
planning question becomes a port-resolvent question.
\end{enumerate}
This is why the dissertation develops a hierarchy instead of a single graph
index. Classical methods are valuable first screens, but they fail when the
protected margin depends on heterogeneous damping, hidden controller memory, or
finite line actions. The rest of the thesis formalizes exactly when each level
is enough and when the analysis must escalate.

\section{Dynamic Effective Resistance as the Bridge}

The classical effective resistance of a line \(e=(i,j)\) is
\begin{equation}
  R_e^{\rm stat}=a_e^TL^\dagger a_e,
  \qquad a_e=e_i-e_j .
  \label{eq:framework_static_resistance}
\end{equation}
This quantity measures static graph compliance between the endpoints of the
line. The thesis generalizes it by replacing the static pseudoinverse with the
dynamic Green function:
\begin{equation}
  \mathcal R_e(s)=a_e^TT(s)^{-1}a_e .
  \label{eq:framework_dynamic_resistance}
\end{equation}
At \(s=0\) and \(\Pi(0)=0\), this reduces to
\(a_e^TL^\dagger a_e\). Away from zero, it is a frequency-dependent,
controller-aware port compliance.

This definition matters because a finite line action has an exact scalar pole
condition:
\begin{equation}
  \det\!\left(T(s)+\Delta b_ea_ea_e^T\right)
  =
  \det T(s)\left(1+\Delta b_e\mathcal R_e(s)\right).
  \label{eq:framework_line_determinant}
\end{equation}
Thus the new or shifted poles created by the line action satisfy
\begin{equation}
  1+\Delta b_e\mathcal R_e(s)=0.
  \label{eq:framework_line_condition}
\end{equation}
The usual line sensitivity is the local residue limit of this finite-action
law. That is the reason dynamic effective resistance is the bridge between
spectral graph theory and controller-aware planning.

\section{Operator Hierarchy and Escalation Logic}

The thesis uses a hierarchy because no single model level is always necessary
or always sufficient. The cheapest valid model should be used, but it must pass
the gate attached to the claim.

\begin{figure}[ht]
\centering
\resizebox{\textwidth}{!}{%
\begin{tikzpicture}[
  level/.style={rectangle,rounded corners=3pt,draw=#1!75!black,fill=#1!9,
    minimum height=1.12cm,text width=2.7cm,align=center,font=\small},
  arrow/.style={-{Latex[length=2.2mm]},draw=thesisgray,line width=0.7pt},
  gate/.style={diamond,aspect=2.1,draw=thesisink!55,fill=softyellow,
    align=center,font=\scriptsize,text width=2.15cm,inner sep=1pt}
]
\node[level=thesisblue] (l0) {\(D=\gamma M\)\\uniform decay};
\node[gate, right=0.40cm of l0] (g0) {commuting?};
\node[level=thesisgreen, right=0.40cm of g0] (l1) {fixed \(D\)\\QEP};
\node[gate, right=0.40cm of l1] (g1) {modal coupling small?};
\node[level=thesisorange, right=0.40cm of g1] (l2) {Schur NEP\\\(\Pi(s)\)};
\node[gate, below=0.70cm of l2] (g2) {finite action or uncertainty?};
\node[level=thesisred, left=0.55cm of g2] (l3) {Green function\\\(T(s)^{-1}\)};
\node[level=thesisteal, right=0.55cm of g2] (l4) {structured fragility\\\(r_{\rm frag}\)};
\node[level=thesisgray, below=0.70cm of g2] (l5) {nonlinear memory\\\(K(x,\tau)\)};
\draw[arrow] (l0) -- (g0);
\draw[arrow] (g0) -- (l1);
\draw[arrow] (l1) -- (g1);
\draw[arrow] (g1) -- (l2);
\draw[arrow] (l2) -- (g2);
\draw[arrow] (g2) -- (l3);
\draw[arrow] (g2) -- (l4);
\draw[arrow] (g2) -- (l5);
\end{tikzpicture}}
\caption{Escalation logic of the thesis. The analysis becomes richer only
when the protected claim requires the richer operator: modal coupling,
controller memory, finite actions, structured uncertainty, or nonlinear
amplitude effects.}
\label{fig:framework_escalation}
\end{figure}

Figure~\ref{fig:framework_escalation} is the practical answer to why the full
analysis is necessary. If damping commutes with stiffness, graph modes can
certify more. If damping does not commute, the fixed-\(D\) QEP is needed. If
controller states are condensed, \(\Pi(s)\) is needed. If the question concerns
a finite line, damping, inertia, or controller action, \(T(s)^{-1}\) is needed.
If the question concerns minimum physical perturbation to failure, structured
fragility is needed. If the question concerns finite-amplitude behavior,
nonlinear memory is needed.

\section{Decision Contract for the Rest of the Thesis}

Every claim in the rest of the dissertation should be readable through the
same four-part contract:
\begin{enumerate}
\item \textbf{Operator:} Which model is being used: \(L\), QEP, Schur NEP,
Green function, structured fragility, or nonlinear memory?
\item \textbf{Port or action:} Which line, bus, converter, controller channel,
input-output pair, or uncertainty family is being tested?
\item \textbf{Margin:} Which quantity is protected: damping ratio, spectral
abscissa, protected gain, passivity, resonance distance, or hosting feasibility?
\item \textbf{Certificate:} Which theorem, bound, contour test, finite-action
law, or full-model validation turns the calculation into evidence?
\end{enumerate}
Without these four items, a weak-bus list, line-ranking plot, or hosting number
is not reproducible. With these four items, the thesis can distinguish a graph
screen from a dynamic certificate.

\FloatBarrier

\chapter{From Power Flow to Spectral Dynamics}

This chapter derives the mathematical objects used throughout the thesis from
the electrical model upward. The purpose is pedagogical as much as technical:
the reader should see where the matrix \(L\) comes from, why the swing equation
has the form \(M\ddot\theta+D\dot\theta+L\theta=p\), and why substituting
\(e^{st}\) produces the quadratic operator \(s^2M+sD+L\). Later chapters use
these equations aggressively, so this chapter builds them slowly.

\section{The Physical Picture}

A power grid is a network of buses connected by lines and transformers. A bus
is a node where generators, loads, converters, transformers, or transmission
lines meet. At each AC bus, the voltage is not only a magnitude. It is a
sinusoid with a phase angle. In phasor notation,
\begin{equation}
  V_i=|V_i|e^{j\theta_i},
  \label{eq:guide_voltage_phasor}
\end{equation}
where \(|V_i|\) is the voltage magnitude and \(\theta_i\) is the phase angle.
Power transfer in an AC network depends strongly on differences of angles,
not on absolute angles. If every angle is shifted by the same constant, the
physics is unchanged. This is the reason a Laplacian matrix will later have
one zero eigenvalue.

The network itself is represented by the bus-admittance matrix
\begin{equation}
  Y_{\rm bus}=G+jB.
  \label{eq:guide_ybus}
\end{equation}
This matrix is built from the impedances of lines and transformers. It maps
bus voltages to injected currents:
\begin{equation}
  I=Y_{\rm bus}V.
  \label{eq:guide_current_balance}
\end{equation}
Equation~\eqref{eq:guide_current_balance} is Kirchhoff's current law in matrix
form. It is the electrical network before any dynamic-machine model is added.

\section{Complex Power and Power Flow}

The complex power injected at bus \(i\) is
\begin{equation}
  S_i=P_i+jQ_i=V_i I_i^*,
  \label{eq:guide_complex_power}
\end{equation}
where \(P_i\) is active power and \(Q_i\) is reactive power. Substituting
\(I=Y_{\rm bus}V\) gives the AC power-flow equations:
\begin{align}
  P_i(V,\theta)
  &=
  |V_i|\sum_{j=1}^n |V_j|
  \left(
  G_{ij}\cos(\theta_i-\theta_j)
  +
  B_{ij}\sin(\theta_i-\theta_j)
  \right),
  \label{eq:guide_active_power}\\
  Q_i(V,\theta)
  &=
  |V_i|\sum_{j=1}^n |V_j|
  \left(
  G_{ij}\sin(\theta_i-\theta_j)
  -
  B_{ij}\cos(\theta_i-\theta_j)
  \right).
  \label{eq:guide_reactive_power}
\end{align}
These equations determine the operating point. A power-flow solution provides
the steady voltages \(V^\star\), angles \(\theta^\star\), and injections that
balance generation and load. The dynamic model in this thesis is always a
small-signal model around such an operating point.

\begin{remark}[Where power flow enters the thesis]
The matrix \(L\) is not chosen by hand. It is the linearized active-power
response around a solved operating point. If the loading, generation dispatch,
network topology, line parameters, or converter injections change, the power
flow changes, and therefore the stiffness matrix \(L\) changes.
\end{remark}

\section{From Active-Power Jacobian to the Laplacian \texorpdfstring{\(L\)}{L}}

To obtain a dynamic model, perturb the angles around the operating point:
\begin{equation}
  \theta=\theta^\star+\Delta\theta.
  \label{eq:guide_angle_perturbation}
\end{equation}
The first-order change in active power is
\begin{equation}
  \Delta P
  \approx
  J_{P\theta}\Delta\theta,
  \qquad
  J_{P\theta}
  =
  \left.
  \frac{\partial P}{\partial\theta}
  \right|_{(V^\star,\theta^\star)}.
  \label{eq:guide_active_jacobian}
\end{equation}
In this dissertation \(P\) denotes exported electrical power in the swing
equation. With that convention, the synchronizing stiffness is this incremental
power response:
\begin{equation}
  L=J_{P\theta}
  \quad
  \hbox{on the retained angle coordinates.}
  \label{eq:guide_L_from_jacobian}
\end{equation}
If instead \(P\) is defined as net injected disturbance with the opposite sign,
the formula becomes \(L=-J_{P\theta}\). The sign convention must be fixed before
using any derivative formula.

Under the common lossless approximation \(G_{ij}\approx0\), fixed voltage
magnitudes, and sufficiently small angle differences, the off-diagonal
entries become
\begin{equation}
  L_{ij}
  =
  -|V_i^\star||V_j^\star|B_{ij}
  \cos(\theta_i^\star-\theta_j^\star),
  \qquad i\ne j,
  \label{eq:guide_L_offdiag}
\end{equation}
and the diagonal is chosen so that each row sums to zero:
\begin{equation}
  L_{ii}=-\sum_{j\ne i}L_{ij},
  \qquad
  L\ones=0.
  \label{eq:guide_L_rowsum}
\end{equation}
This is the graph Laplacian structure. The line weight is roughly
\(|V_i^\star||V_j^\star|B_{ij}\cos(\theta_i^\star-\theta_j^\star)\). A line
with high susceptance, high voltage magnitudes, and small angle separation
contributes stronger synchronizing stiffness.

\begin{example}[Two-bus stiffness]
For two buses connected by one lossless line, active power from bus 1 to bus 2
is approximately
\[
  P_{12}=|V_1||V_2|B_{12}\sin(\theta_1-\theta_2).
\]
Linearizing around \(\theta^\star\) gives
\[
  \Delta P_{12}
  \approx
  |V_1^\star||V_2^\star|B_{12}
  \cos(\theta_1^\star-\theta_2^\star)
  (\Delta\theta_1-\Delta\theta_2).
\]
Thus the line resists angle separation like a spring. That spring constant is
the line weight that appears inside \(L\).
\end{example}

\section{The Single-Machine Swing Equation}

The swing equation is Newton's second law for rotor angle. A synchronous
machine stores kinetic energy in a rotating mass. Let \(\theta_i\) be the
electrical angle of machine \(i\) relative to a synchronous reference, and let
\(\omega_i=\dot\theta_i\) be the frequency deviation. The power imbalance
accelerates the rotor:
\begin{equation}
  M_i\dot\omega_i
  =
  P_{m,i}-P_{e,i}-D_i\omega_i.
  \label{eq:guide_single_swing_speed}
\end{equation}
Here \(P_{m,i}\) is mechanical or commanded input power, \(P_{e,i}\) is
electrical air-gap power exported to the network, \(D_i\omega_i\) is a
linearized damping or primary-control term, and \(M_i\) is the inertia
coefficient in the chosen per-unit convention. Since \(\omega_i=\dot\theta_i\),
\eqref{eq:guide_single_swing_speed} becomes
\begin{equation}
  M_i\ddot\theta_i+D_i\dot\theta_i
  =
  P_{m,i}-P_{e,i}.
  \label{eq:guide_single_swing_angle}
\end{equation}

At the operating point, \(P_{m,i}^\star=P_{e,i}^\star\). The dynamic equation
therefore concerns deviations from that balance. Let \(p_i\) be a small input
disturbance and let \(\Delta P_{e,i}\) be the incremental electrical power.
Then
\begin{equation}
  M_i\Delta\ddot\theta_i+D_i\Delta\dot\theta_i+\Delta P_{e,i}=p_i.
  \label{eq:guide_single_linear_swing}
\end{equation}
The previous section showed that, to first order,
\(\Delta P_e\approx L\Delta\theta\). Therefore the electrical network supplies
the restoring term in the swing equation.

\section{The Networked Swing Model}

Stack all retained angles into \(\theta\), all inertia coefficients into
\(M=\diag(M_i)\), and all damping coefficients into \(D=\diag(D_i)\). Dropping
the \(\Delta\) notation for readability gives
\begin{equation}
  M\ddot\theta+D\dot\theta+L\theta=p.
  \label{eq:guide_network_swing}
\end{equation}
This is the model used in the fixed-\(D\) chapters. Each term has a physical
meaning:
\begin{itemize}
\item \(M\ddot\theta\): inertia resists acceleration;
\item \(D\dot\theta\): damping or primary control resists speed deviation;
\item \(L\theta\): network synchronizing stiffness resists angle separation;
\item \(p\): disturbance or control input.
\end{itemize}

Because \(L\ones=0\), the uniform angle direction must be removed by selecting
a reference bus, center-of-inertia coordinates, or a projection onto
\(\ones^\perp\). After that reduction, \(L\) is positive definite on the
retained nonzero-angle subspace if the network is connected and the operating
point is regular.

\section{Why the Exponential Ansatz Produces \texorpdfstring{\(s^2M+sD+L\)}{s2M+sD+L}}

Small-signal stability asks what happens to a tiny perturbation if no further
input is applied. Set \(p=0\) and look for modal solutions of the form
\begin{equation}
  \theta(t)=ve^{st},
  \label{eq:guide_modal_ansatz}
\end{equation}
where \(v\) is the mode shape and \(s\) is the complex growth rate. Then
\begin{equation}
  \dot\theta(t)=sve^{st},\qquad
  \ddot\theta(t)=s^2ve^{st}.
  \label{eq:guide_exponential_derivatives}
\end{equation}
Substituting \eqref{eq:guide_modal_ansatz} into
\eqref{eq:guide_network_swing} gives
\[
  Ms^2ve^{st}+Dsve^{st}+Lve^{st}=0.
\]
Since \(e^{st}\ne0\), it can be divided out:
\begin{equation}
  (s^2M+sD+L)v=0.
  \label{eq:guide_qep_from_swing}
\end{equation}
This is the quadratic eigenvalue problem (QEP). It is not an arbitrary
mathematical object. The factor \(s^2\) comes from acceleration, \(s\) comes
from velocity damping, and the constant term comes from the network restoring
force.

\begin{table}[ht]
\centering
\caption{Physical origin of the QEP terms.}
\label{tab:qep_origin}
\begin{tabular}{lll}
\toprule
\rowcolor{softblue}
\thead{Term} & \thead{Physical source} & \thead{Engineering interpretation}\\
\midrule
\(s^2M\) & rotor or virtual mass & acceleration and inertia\\
\(sD\) & damping, droop, primary response & decay of speed deviations\\
\(L\) & linearized active-power network & synchronizing stiffness\\
\((sI-A_{cc})^{-1}\) & eliminated controller states & dynamic memory of controls\\
\bottomrule
\end{tabular}
\end{table}

\section{How to Read a Complex Pole}

If a root is \(s=\sigma+j\omega\), the perturbation contains
\[
  e^{st}=e^{\sigma t}e^{j\omega t}.
\]
The real part \(\sigma\) controls decay or growth. The imaginary part
\(\omega\) controls oscillation frequency. Stability requires
\(\sigma<0\). The damping ratio is
\begin{equation}
  \zeta=-\frac{\sigma}{\sqrt{\sigma^2+\omega^2}}.
  \label{eq:guide_zeta}
\end{equation}
This number measures how quickly the oscillation decays relative to its
frequency. A line upgrade can increase \(\omega\) by stiffening the network,
but if \(\sigma\) does not become more negative, \(\zeta\) may decrease. This
is the basic reason graph stiffness and damping margin must be kept separate.

\section{Where Inverters Enter the Same Chain}

An inverter-based resource can enter this chain at several levels. In a simple
reduced model, a grid-forming inverter may be represented by effective
inertia, damping, or droop terms added to \(M\) and \(D\). A grid-following
inverter may alter the operating point and the Jacobian through its injection
and controls. In a detailed small-signal model, inverter controllers add
states: PLL states, current-loop states, voltage-loop states, filters, limits,
and plant controls. If those states are retained, they appear in the full
state matrix. If they are mathematically eliminated, they appear through a
Schur complement and produce the rational self-energy \(\Pi(s)\) developed in
the self-energy part of the thesis. The dynamic damping object \(\Sigma(s)\)
is the special case where
the rational correction enters through the velocity channel.

Thus the thesis path is:
\begin{equation}
  \hbox{power flow}
  \rightarrow
  L
  \rightarrow
  M\ddot\theta+D\dot\theta+L\theta=p
  \rightarrow
  s^2M+sD+L
  \rightarrow
  s^2M+sD_0+L+\Pi(s).
  \label{eq:guide_thesis_path}
\end{equation}
Every later result is a refinement of one arrow in this chain.

\begin{figure}[ht]
\centering
\begin{tikzpicture}[
  stage/.style={
    rectangle,
    rounded corners=3pt,
    draw=#1!70!black,
    fill=#1!10,
    text width=0.25\textwidth,
    minimum height=1.15cm,
    align=center,
    font=\small
  },
  flow/.style={-{Latex[length=2.2mm]}, thick, draw=thesisgray},
  note/.style={font=\scriptsize\color{thesisgray}, align=center}
]
\node[stage=thesisblue] (pf) {\textbf{Power-flow solution}\\
\(V^\star,\theta^\star,Y_{\rm bus}\)};
\node[stage=thesisgreen, right=0.65cm of pf] (jac) {\textbf{Active-power Jacobian}\\
\(L=J_{P\theta}\)};
\node[stage=thesisorange, right=0.65cm of jac] (swing) {\textbf{Swing layer}\\
\(M\ddot\theta+D\dot\theta+L\theta=p\)};
\node[stage=thesisblue, below=0.85cm of jac] (qep) {\textbf{Fixed-\(D\) spectral problem}\\
\((s^2M+sD+L)v=0\)};
\node[stage=thesisred, right=0.65cm of qep] (nep) {\textbf{Condensed-controller problem}\\
\(T(s)=s^2M+sD_0+L+\Pi(s)\)};
\draw[flow] (pf) -- (jac);
\draw[flow] (jac) -- (swing);
\draw[flow] (swing.south) -- ++(0,-0.42) -| ([xshift=0.32cm]qep.north);
\draw[flow] (jac.south) -- (qep.north);
\draw[flow] (qep) -- (nep);
\node[note, below=0.18cm of pf] {electrical constraints};
\node[note, below=0.18cm of qep] {modal damping screen};
\node[note, below=0.18cm of nep] {controller memory};
\end{tikzpicture}
\caption{Modeling chain from AC power flow to the rational NEP. The same
network operating point that sets voltages and flows also defines the
stiffness operator \(L\). Controller condensation changes the operator class,
because the damping contribution becomes frequency dependent.}
\label{fig:prelim_modeling_chain}
\end{figure}

\FloatBarrier

\section{Linearization and Spectral Parameter}

Let a power-system dynamic model be written locally as a differential-algebraic
system,
\begin{align}
  \dot x &= f(x,y,\rho),\\
  0 &= g(x,y,\rho),
\end{align}
where \(x\) are dynamic states, \(y\) are algebraic variables, and \(\rho\) is
a parameter or action. At a regular equilibrium
\((x^\star,y^\star)\), algebraic variables can be eliminated locally when
\(g_y\) is nonsingular, giving a reduced state matrix
\begin{equation}
  A_{\rm red}=f_x-f_y g_y^{-1}g_x.
  \label{eq:dae_reduced_A}
\end{equation}
Small-signal stability is determined by the eigenvalues of \(A_{\rm red}\).
The spectral parameter is denoted by \(s\in\C\), and a pole is stable when
\(\Real(s)<0\).

For an oscillatory pole \(s=\sigma+j\omega\), the frequency and damping ratio
are
\begin{equation}
  f=\frac{|\omega|}{2\pi},\qquad
  \zeta=-\frac{\sigma}{\sqrt{\sigma^2+\omega^2}}.
  \label{eq:prelim_zeta}
\end{equation}
The damping-ratio derivative satisfies
\begin{equation}
  d\zeta
  =
  -\frac{\omega^2}{(\sigma^2+\omega^2)^{3/2}}\,d\sigma
  +
  \frac{\sigma\omega}{(\sigma^2+\omega^2)^{3/2}}\,d\omega,
  \label{eq:prelim_dzeta}
\end{equation}
so a frequency shift alone is not a damping certificate.

\section{Graph Laplacians and Synchronizing Stiffness}

Under a lossless active-power approximation, the synchronizing stiffness matrix
has Laplacian structure. For \(i\ne j\),
\begin{equation}
  L_{ij}=-V_iV_j b_{ij}\cos(\theta_i^\star-\theta_j^\star),
  \label{eq:prelim_lap_offdiag}
\end{equation}
with diagonal entries chosen so that
\begin{equation}
  L\ones=0.
\end{equation}
If the network is reciprocal and the operating point is sufficiently regular,
\(L=L^T\succeq0\). The zero eigenvalue corresponds to uniform angle shift and
is removed by fixing a reference or projecting onto \(\ones^\perp\).

Let \(M=\diag(M_i)\succ0\). The inertia-normalized stiffness is
\begin{equation}
  \Lt=M^{-1/2}LM^{-1/2}.
  \label{eq:prelim_ltilde}
\end{equation}
On the nonzero subspace, write
\begin{equation}
  \Lt Q=Q\Lambda,\qquad Q^TQ=I,\qquad
  \Lambda=\diag(\nu_1,\ldots,\nu_n),\quad \nu_k>0.
  \label{eq:prelim_graph_modes}
\end{equation}
The columns of \(Q\) are graph-inertia modes. They diagonalize stiffness. They
do not automatically diagonalize damping.

\section{Second-Order Systems and QEPs}

The retained second-order model is
\begin{equation}
  M\ddot\theta+D\dot\theta+L\theta=p.
  \label{eq:prelim_second_order}
\end{equation}
Substituting \(\theta(t)=ve^{st}\) gives
\begin{equation}
  P(s)v=(s^2M+sD+L)v=0.
  \label{eq:prelim_qep}
\end{equation}
This is a QEP. The corresponding first-order linearization is
\begin{equation}
  \mathcal A=
  \begin{bmatrix}
    0&I\\
    -M^{-1}L&-M^{-1}D
  \end{bmatrix},
  \label{eq:prelim_first_order}
\end{equation}
but the second-order form should be retained conceptually because the matrices
\(M\), \(D\), and \(L\) have distinct physical meanings.

In inertia-normalized variables \(\vartheta=M^{1/2}\theta\),
\begin{equation}
  \ddot\vartheta+\Dt\dot\vartheta+\Lt\vartheta=M^{-1/2}p,\qquad
  \Dt=M^{-1/2}DM^{-1/2}.
  \label{eq:prelim_normalized}
\end{equation}
The modal damping matrix is
\begin{equation}
  \Gamma=Q^T\Dt Q.
  \label{eq:prelim_gamma}
\end{equation}
The equality \(Q^T\Lt Q=\Lambda\) is automatic by construction; the equality
\(Q^T\Dt Q=\diag(\cdot)\) is an additional damping assumption.

\section{Simple Eigenvalue Sensitivity}

For a polynomial or nonlinear eigenvalue problem
\begin{equation}
  T(s,\rho)v=0,\qquad u^HT(s,\rho)=0,
\end{equation}
with a simple eigenvalue \(s_\star\), the first derivative is
\begin{equation}
  \frac{ds_\star}{d\rho}
  =
  -\frac{u^H T_\rho(s_\star,\rho)v}
  {u^H T_s(s_\star,\rho)v}.
  \label{eq:prelim_nep_sens}
\end{equation}
For a QEP \(T(s)=s^2M+sD+L\), this becomes
\begin{equation}
  \frac{ds_\star}{d\rho}
  =
  -\frac{u^H(s_\star D_\rho+L_\rho)v}
  {u^H(2s_\star M+D)v}
  \label{eq:prelim_qep_sens}
\end{equation}
when \(M\) is fixed. Equation~\eqref{eq:prelim_qep_sens} is central because it
shows that line actions and damping actions enter the numerator differently,
while the denominator depends on the mode and the damping operator.

\section{Schur Complements}

Partition a full first-order linearization into retained states \(x_e\) and
condensed states \(x_c\):
\begin{equation}
  A=
  \begin{bmatrix}
    A_{ee} & A_{ec}\\
    A_{ce} & A_{cc}
  \end{bmatrix}.
  \label{eq:prelim_partitioned_A}
\end{equation}
The eigenvalue equation \((sI-A)x=0\) is equivalent, when \(sI-A_{cc}\) is
invertible, to
\begin{equation}
  \left[sI-A_{ee}-A_{ec}(sI-A_{cc})^{-1}A_{ce}\right]x_e=0.
  \label{eq:prelim_schur}
\end{equation}
This is an exact nonlinear eigenvalue problem for the retained states. It does
not approximate the eliminated states; it represents them through their
resolvent.

\section{Rational Matrix Functions and Self-Energy}

A rational matrix function has the form
\begin{equation}
  R(s)=R_0+\sum_{\ell=1}^m \frac{R_\ell}{s-\alpha_\ell}
  \label{eq:prelim_rational}
\end{equation}
after partial-fraction expansion when the poles are simple. The residues
\(R_\ell\) determine how strongly each eliminated pole \(\alpha_\ell\) couples
to the retained coordinates. In the Schur operator, those residues are
generated by \(A_{ec}\), \(A_{ce}\), and the eigenvectors of \(A_{cc}\).

In a second-order retained representation, the rational contribution appears as
a self-energy,
\begin{equation}
  T(s)=s^2M+sD_0+L+\Pi(s),
  \qquad
  \Pi(s)=G(sI-A_{cc})^{-1}(B_\theta+sB_\omega).
  \label{eq:prelim_sigma}
\end{equation}
The term self-energy is useful because it emphasizes that the retained model is
not isolated. It carries the dynamic influence of states that have been
mathematically eliminated. When \(\Pi(s)=s\Pi_D(s)\), the dynamic damping
matrix is \(\Sigma(s)=D_0+\Pi_D(s)\).

\section{Certificates, Screens, and Gates}

The thesis uses three words with specific meanings.

\begin{definition}[Certificate]
A certificate is a mathematical or numerical condition that directly implies a
claimed margin property under stated assumptions.
\end{definition}

\begin{definition}[Screen]
A screen is a lower-cost indicator that may identify candidates for action or
inspection but does not imply the claimed margin property by itself.
\end{definition}

\begin{definition}[Gate]
A gate is an accept/reject test that a proposed certificate or screen must pass
before it is used for a stronger claim.
\end{definition}

The distinction matters throughout the dissertation. A graph eigenvalue can be
a certificate in a proportional or commuting fixed-\(D\) theory, a screen in a
non-commuting fixed-\(D\) theory, and an incomplete descriptor after controller
condensation. The same number changes epistemic status when the operator class
changes.

\chapter{Literature Review and Theoretical Gap}

The theory developed in this dissertation sits at the intersection of five
literatures: classical power-system stability, spectral and network
interpretations of synchronization, second-order and quadratic eigenvalue
analysis, nonlinear eigenvalue problems generated by state condensation, and
converter-dominated stability. The purpose of this chapter is not to catalogue
every contribution in those areas. It is to identify the exact conceptual gap
that motivates a damping-operator hierarchy.

\section{Systematic Review Contract}

In the expanded dissertation this chapter should become a systematic review,
not only a narrative survey. The review protocol should report search strings,
databases, dates, inclusion and exclusion criteria, preprint versus
peer-reviewed status, extraction fields, and a method-assumption-data-certificate
matrix. The central taxonomy is the operator hierarchy
\begin{equation}
  \begin{aligned}
  D=\gamma M
  &\rightarrow
  [\Lt,\Dt]=0
  \rightarrow
  D\ {\rm fixed\ and\ non\!-\!commuting}
  \rightarrow
  \Pi(s)\ {\rm rational} \\
  &\rightarrow
  G_T(s)=T(s)^{-1}
  \rightarrow
  \Lambda_\varepsilon^{\mathcal A}(T)
  \rightarrow
  {\rm nonlinear\ memory}.
  \end{aligned}
  \label{eq:lit_operator_taxonomy}
\end{equation}
The first arrow marks the move from uniform damping-to-inertia ratio to the
larger decouplable class. The second marks modal coupling within a constant
QEP. The third marks controller or auxiliary-state condensation. The fourth
marks the Green function that carries finite actions. The fifth marks
structured fragility under physical perturbation families. The last arrow
marks the move from linearized spectral memory to genuinely nonlinear memory.
This taxonomy is the review's organizing principle.

\section{Classical Small-Signal Stability}

Classical power-system stability theory starts from an operating point and
studies the linearized differential-algebraic model around that point. The
standard engineering workflow is well established: solve the power flow,
initialize dynamic models, linearize, compute eigenvalues, inspect damping
ratios, and use eigenvectors or participation factors to interpret the critical
modes \cite{kundur1994power,rogers2000oscillations}. The IEEE/CIGRE stability
definition paper formalized the classification of rotor-angle, frequency, and
voltage stability and placed small-signal stability within a broader taxonomy
of disturbances, time scales, and mechanisms \cite{kundur2004definition}.

For the present thesis, the central fact is that classical modal analysis is
already an operator theory. A mode is not merely an oscillation observed in a
time-domain trace. It is a spectral object of a linearized system. Its damping
ratio is computed from the real and imaginary parts of a pole. Its engineering
interpretation depends on how the corresponding eigenvectors project onto
angles, speeds, exciters, governors, loads, and controls. Thus, any reduced
model that claims to preserve damping must preserve the relevant spectral
mechanism, not just the visual shape of a response.

\section{Graph and Synchronization Perspectives}

Power networks are also graphs. Under lossless or nearly lossless assumptions,
the active-power Jacobian has a Laplacian structure: off-diagonal terms encode
coupling between connected buses and diagonal terms enforce row-sum balance.
This structure links power-system dynamics to synchronization theory. The work
of Doerfler, Chertkov, and Bullo gives a concise synchronization condition for
complex oscillator networks and smart grids, clarifying how coupling strength,
phase cohesiveness, and network structure interact \cite{dorfler2013synchronization}.

The graph viewpoint is indispensable because it turns a high-dimensional
network into interpretable objects: algebraic connectivity, effective
resistance, cut structure, mode shapes, Kron-reduced Laplacians, and weighted
network strength. Algebraic connectivity traces back to Fiedler's second
Laplacian eigenvalue and its relation to graph connectivity
\cite{fiedler1973algebraic}. Effective-resistance design provides another
view of robustness by penalizing large pairwise electrical distances and by
linking network reinforcement to convex edge-weight allocation
\cite{ghosh2008effective}. Spectral sparsification shows why effective
resistance is also a structural sampling quantity: preserving selected
quadratic forms preserves the Laplacian action that matters to many slow-mode
metrics \cite{spielman2011sparsification}. Kron reduction gives the matching
electrical-network operation for eliminating passive buses through a Schur
complement while preserving a reduced Laplacian
\cite{dorfler2013kron}.

These tools motivate planning intuitions: reinforce a weak line, improve a
weak cut, place a grid-forming converter at a strategic bus, or raise a
critical graph eigenvalue. Those intuitions are valuable. The problem is that
damping is not always a graph eigenvalue. Damping may be local, heterogeneous,
non-proportional, controller-mediated, or frequency dependent. The graph
literature gives the stiffness skeleton; the present thesis asks which damping
operator must be attached to that skeleton.

This boundary is important for the originality claim. For example, Guo, Zhao,
and Low already connect the scaled graph Laplacian spectrum with primary
frequency regulation under uniform damping-to-inertia ratios
\cite{guo2018graph}. The present thesis uses that kind of result as a known
floor. Its contribution is not the existence of Laplacian modes, but the
operator hierarchy that decides when those modes remain damping coordinates,
when rough damping-to-inertia fields mix them, and when controller self-energy
must replace a static graph screen.

\section{Beyond Scalar Laplacians: Cellular Sheaves}

An important structural extension is to look beyond scalar graph
Laplacians. The closest mature mathematical object is the cellular sheaf
Laplacian. Hansen and Ghrist develop spectral sheaf theory by lifting the
ordinary graph Laplacian to a Hodge Laplacian acting on a sheaf of vector
spaces over a cell complex \cite{hansen2019sheaf}. In that setting, vertices
and edges may carry vector spaces, and restriction maps encode how local data
are compared across cells. The theory connects spectral data to sheaf
cohomology and includes analogues of interlacing, sparsification, effective
resistance, and synchronization.

This dissertation does not claim a completed sheaf theory of converter-rich
power grids. The point is more precise: a scalar Laplacian assigns one
number-valued state to each bus, whereas a converter-rich bus carries internal
controller coordinates. A sheaf model is a natural candidate for a future
operator in which the bus stalk contains retained electrical variables,
converter-control states, or reduced controller coordinates, and each line
restriction map describes which combinations are electrically compatible. Such
an operator could turn the present scalar mode-family labels and
self-energy-dominance diagnostics into a structured vector-bundle-like object.

The risk is equally clear. A sheaf Laplacian is useful only if it detects a
weak node, weak link, or hidden margin that the scalar \(L\), the fixed-\(D\)
QEP, and the Schur self-energy screen do not already detect. In this thesis it
is therefore placed as a speculative theoretical extension, not as evidence
for any current planning claim.

\section{Quadratic Eigenvalue Problems and Second-Order Models}

The swing-equation model is naturally second order. In compact form it leads to
a quadratic eigenvalue problem,
\begin{equation}
  P(s)v=(s^2M+sD+L)v=0.
  \label{eq:lit_qep}
\end{equation}
Tisseur and Meerbergen's survey establishes the QEP as a mature mathematical
object, with well-developed linearizations, spectral properties, perturbation
formulas, and numerical methods \cite{tisseur2001qep}. This literature matters
because the power-system damping problem should not be reduced too quickly to a
first-order matrix. The second-order structure carries physical information:
mass or inertia, damping, and stiffness are distinct operators.

The distinction between proportional and non-proportional damping is also
well known in structural dynamics. If damping and stiffness share a modal
basis, modal decoupling is exact. If they do not, modal analysis becomes
coupled, and non-proportionality indices are needed to measure the error of
decoupled approximations \cite{conde2021nonproportionality}. This thesis uses
the same mathematical distinction in a power-system setting. The first
departure from graph-frequency intuition is not necessarily converter
dynamics; it can already occur when \(D\) is fixed but heterogeneous relative
to \(M\).

\section{Nonlinear Eigenvalue Problems and Schur Condensation}

The next mathematical layer is the nonlinear eigenvalue problem (NEP). Guettel
and Tisseur survey matrix-valued functions \(T(s)\) whose dependence on the
spectral parameter is nonlinear and show that such problems arise naturally
when delayed, damped, or condensed dynamics are represented without expanding
all states \cite{guettel2017nep}. A Schur complement is a basic route from a
larger linear eigenvalue problem to a smaller nonlinear one. If the eliminated
block is \(A_{cc}\), the retained operator contains the resolvent
\((sI-A_{cc})^{-1}\).

For computation, contour integral methods and rational-interpolation
viewpoints are particularly relevant. Brennan, Embree, and Gugercin recast
contour integral methods for NEPs in a systems-theoretic framework using
Loewner matrix pencils, allowing interpolation away from infinity and linking
eigenvalue extraction to system realization and modal truncation
\cite{brennan2023contour}. This is important for a Schur self-energy thesis
because \(\Pi(s)\) and its damping-channel projection \(\Sigma(s)\) are
transfer-like rational objects. If the
controller block is not fully known, a Loewner or data-driven contour
formulation is a plausible route for identifying the visible part of
\(\Pi(s)\) from frequency-domain samples rather than from explicit
state-space matrices.

Recent power-system work also strengthens the numerical side of the proposed
planning workflow. Bouterakos, McKeon, and Tzounas develop continuation-based
eigenvalue tracking for large-scale power-system models, supporting sparse
representations and both explicit and semi-implicit DAE forms
\cite{bouterakos2025tracking}. That result is directly aligned with the local
pole-motion map used later in this dissertation: a planning screen that changes
penetration, line weights, inertia, damping, or controller gains must track
selected eigenvalue trajectories rather than repeatedly rediscovering the entire
spectrum from scratch.

This observation is the mathematical hinge of the dissertation. Converter and
auxiliary controller states can be retained explicitly, producing a large
first-order model. They can also be condensed, producing a smaller nonlinear
operator. The condensation is exact away from poles of the eliminated block,
but the price of exactness is frequency dependence. A constant damping matrix
cannot store the poles, zeros, and residues introduced by the eliminated
controller states. The reduced controller object is therefore a rational
self-energy; its damping-channel projection is the rational object denoted
\(\Sigma(s)\).

\section{Pseudospectra, Fragility, and Model Adequacy}

Nonlinear eigenvalue theory also raises a robustness question. A nominal root
set may look safe while nearby structured perturbations can move a pole to a
protected boundary, especially when the retained operator is non-normal. The
classical pseudospectral viewpoint already warns that eigenvalues alone can
miss transient amplification and perturbation sensitivity
\cite{trefethen2005spectra}. Recent work on nonlinear spectral problems
extends the computational discussion of spectra and pseudospectra to broad NEP
settings and emphasizes spectral pollution and invisibility risks in
approximations \cite{colbrook2025universal}.

For this dissertation, the distinction is that power-system perturbations must
be structured. A generic matrix perturbation is not the same as uncertainty in
line susceptance, PLL gains, converter delays, virtual inertia, damping, or an
identified self-energy. This motivates a physical-action pseudospectrum
\(\Lambda_\varepsilon^{\mathcal A}(T)\) and a fragility radius: the smallest
admissible perturbation that drives a retained root to a protected boundary.
The related model-adequacy problem is contour based. If a reduced operator
\(\widehat T\) is used for a decision, it should preserve the protected
eigenvalue count on a contour before its derivative signs are trusted.

\section{Sign Characteristics and Krein-Type Mechanisms}

The mathematical language closest to the proposed controller-resolvent
collision is not a new, free-floating ``Krein signature for NEPs.'' The citable
foundation is the sign characteristic of Hermitian matrix polynomials and
self-adjoint meromorphic matrix functions. Gohberg, Lancaster, and Rodman give
the classical matrix-polynomial framework, including sign characteristics of
spectral data \cite{gohberg1982matrixpolynomials}. Their related work on
self-adjoint meromorphic matrix functions is closer to the rational structure
created by Schur condensation \cite{gohberg1983meromorphic}. Mehrmann,
Noferini, Tisseur, and Xu develop sign characteristics for Hermitian matrix
polynomials and clarify the role of structure-preserving transformations
\cite{mehrmann2016sign}. Kollár and Miller's Evans--Krein survey provides a
systems-oriented view in which rational Krein functions track spectral
collisions \cite{kollar2014graphical}.

Those references do not already solve the converter problem. They mostly live
in Hamiltonian, self-adjoint, or Hermitian-polynomial settings. Converter
self-energies are rational, dissipative, and generally non-normal. The thesis
therefore treats sign characteristic as a foundation to be adapted, not as a
ready-made theorem. The open bridge is to show when a rational
controller-condensed operator has an operational sign collision that coincides
with weak-grid PLL or controller-resonance instability.

\section{Converter-Dominated Stability}

The increased penetration of converter-interfaced generation motivates a
reconsideration of stability definitions and modeling practice. The revisited
IEEE/CIGRE classification explicitly adds converter-driven and electric
resonance phenomena to the stability vocabulary
\cite{hatziargyriou2021stability}. IEEE Std 2800-2022 reflects the same
industrial shift by specifying performance and interconnection requirements for
IBRs connected to transmission systems \cite{ieee2800_2022}.

Within the research literature, grid-following and grid-forming controls are
not interchangeable dynamic objects. Grid-following converters commonly depend
on a PLL for synchronization, so PLL dynamics can contribute to small-signal
instability under weak-grid conditions. Yang, Huang, Xin, and Ju analyze how
placing grid-forming converters can improve the small-signal stability of
PLL-integrated systems by modifying grid strength and Laplacian-related
quantities \cite{yang2021gfm}. Comparative studies of grid-forming and
grid-following inverters further show that control architecture changes the
state-space model and its eigenvalues \cite{ding2021comparative}.

Inertia and primary-control placement studies sharpen the same point. Optimal
virtual inertia allocation can be formulated as an \(H_2\)-type network
performance problem \cite{poolla2017virtual}; disturbance propagation depends
strongly on where inertia is located relative to slow network modes
\cite{pagnier2019inertia_location}; and perturbation theory shows that inertia
and primary-control placement are not interchangeable under heterogeneous
parameters \cite{pagnier2019optimal}. These results support the thesis
distinction between adding inertia, adding damping, and modifying the graph.

Passivity and port-Hamiltonian viewpoints provide a second, complementary
control language. Passivity-based grid-forming work formulates converter
controllers so that local subsystems have dissipativity properties that can be
interconnected modularly. He and D{\"o}rfler derive decentralized stability
conditions for grid-forming converters using passivity indices
\cite{he2024passivity}. Huang et al. combine small-gain and small-phase
arguments to obtain decentralized stability checks for heterogeneous
power-electronics-dominated systems \cite{huang2024gainphase}, while Gui,
Subedi, and Xue use a port-controlled Hamiltonian form for passivity-based
grid-forming control of distributed energy resources \cite{gui2024passivity}.
Recent matrix-valued passivity indices make the same issue sharper for MIMO
ports: scalar indices can discard the direction and coupling structure of a
passivity deficit \cite{ru2026matrix}. Dynamic passivity multipliers further
show why frequency-dependent certificates can be less conservative than strict
static passivity tests for converter-dominated grids \cite{gorbunov2026dynamic}.
In the notation of this thesis, this literature says that the Hermitian,
weighted directional, or multiplier-based behavior of the condensed
self-energy \(\Sigma(s)\) should not be a decorative diagnostic. It is a
candidate certificate for controller retuning, provided that the passivity
screen is tied back to the retained pole motion and to full-model validation.

These studies justify the engineering relevance of the present theory, but
they do not remove the need for a general operator language. A placement rule
derived for a particular control model or grid-strength metric may not
transfer to a different controller, parameter regime, or retained-state
partition. The common mathematical issue is not the label GFL or GFM by
itself; it is whether the reduced model preserves the spectral contribution of
the eliminated control states.

\section{Hosting Capacity, Network Reinforcement, and Braess Effects}

Hosting capacity is usually introduced as a planning quantity: how much
distributed generation, load, or converter-interfaced equipment can be added
before a registered performance limit is violated. Bollen and Roennberg frame
hosting capacity as a distribution-planning tool under uncertainty
\cite{bollen2017hosting}. Later dynamic hosting-capacity studies emphasize
that snapshot limits can miss the time duration, severity, and controller
response of violations \cite{oliveira2019dynamicHosting,jain2020dynamicHosting}.
Those works motivate the same methodological principle used in this thesis:
capacity is not a single number until the limiting metric has been named.

The present dissertation uses the phrase dynamic hosting capacity in a
different but compatible sense. Instead of asking only how long a voltage or
thermal violation lasts, it asks how much IBR conversion can be admitted before
a small-signal spectral margin is lost. Static hosting constraints remain
necessary, but they are insufficient in low-inertia systems because a
voltage-safe operating point can still have an unacceptable damping ratio,
resolvent gain, or controller-resonance exposure.

Recent stability-manifold work for converter-dominated systems reinforces the
need to map controller-parameter regions under multiple operating scenarios
\cite{conte2026stabilitymanifolds}. The distinction here is mechanism: the
present thesis seeks not only a sampled stability boundary, but an analytic
normal, a physical self-energy label, an error certificate, and a minimum
admissible action associated with that boundary.

Network reinforcement has a similar ambiguity. In a purely stiffness-oriented
graph view, increasing a line weight or adding a line seems beneficial.
Braess' original traffic result already showed that adding capacity can worsen
system performance under decentralized routing \cite{braess2005traffic}.
Power-system analogues have been studied for oscillator synchronization and
electric grids. Witthaut and Timme showed that adding links in oscillator
networks can destroy synchronization \cite{witthaut2012braess}, while Coletta
and Jacquod analyzed how added couplings can enhance or reduce linear
stability in coupled-oscillator and power-grid models \cite{coletta2016braess}.
Schaefer and coauthors further demonstrated Braess' paradox in AC power-grid
experiments and large-grid studies, emphasizing that line upgrades can
increase loading elsewhere in the network \cite{schafer2022braess}.

For this thesis, the Braess literature gives a warning rather than a license
to overclaim. A line reinforcement that increases modal frequency but reduces
damping ratio is a Braess-like damping-margin reversal, not automatically a
classical flow Braess paradox. A controller-condensed case in which the line
move pulls a network-family pole toward a nearby control pole is a stronger
mechanistic claim, but it still requires a registered NEP or full-model
validation gate. This calibrated vocabulary is used in the planning chapters.

\section{Computational Modeling and Reproducible Dynamics}

Detailed simulation environments are essential for testing the theory. ANDES,
for example, implements a hybrid symbolic-numeric framework for power-system
DAE modeling and analysis, including eigenvalue analysis and time-domain
simulation \cite{cui2021andes}. Such tools make it possible to build
converter-rich cases, extract Jacobians, compare full and reduced spectra, and
test whether a proposed self-energy approximation is valid.

The next modeling frontier is to identify the port self-energy when the
internal converter model is unavailable. Recent gray-box converter
identification work estimates voltage-source converter parameters from terminal
time-series data and validates the resulting model with frequency-domain
metrics \cite{darii2026graybox}. Structure-preserving identification of
port-Hamiltonian DAEs points in a complementary direction: learn from noisy
input-output data while preserving interconnection and passivity properties
\cite{hagelaars2026identification}. These lines support the thesis requirement
that an identified \(\widehat\Pi(s)\) must carry an uncertainty bound, not just
a fitted curve.

This dissertation, however, intentionally separates theory from validation.
The present manuscript develops the operator framework. A later empirical
chapter or companion article may use ANDES or another DAE platform to test the
framework on specific systems. That separation is methodological: without a
clear theoretical object, simulation can produce plausible plots without
explaining which damping mechanism has been preserved.

\section{Theoretical Gap}

The reviewed literatures leave a precise gap. Classical small-signal stability
provides modal damping ratios and participation factors. Spectral graph theory
provides an interpretable stiffness skeleton. QEP theory explains
second-order fixed-coefficient dynamics. NEP theory explains frequency-dependent
operators created by condensation. Converter-dominated stability shows why
controller dynamics matter. What is missing is a unified damping-operator
hierarchy that tells the engineer which level is mathematically sufficient for
a given claim.

The gap can be stated as four questions.

\begin{enumerate}
\item When does a graph mode carry a true damping certificate rather than only
a frequency screen?
\item When does fixed non-commuting damping create enough modal coupling to
invalidate decoupled modal reasoning?
\item When do controller states require a rational self-energy \(\Pi(s)\),
and when can its damping-channel projection \(\Sigma(s)\) replace a constant
reduced damping matrix?
\item When does a dynamic Green function or structured fragility radius change
the decision that would be made from a scalar graph or damping metric?
\end{enumerate}

The remainder of the dissertation answers these questions by developing the
hierarchy from uniform damping \(D=\gamma M\), to the broader commuting
constant-\(D\) class, to fixed non-commuting damping, to rational
self-energy \(\Pi(s)\), to Green-function action calculus, to structured
fragility, and finally to nonlinear memory.


\part{Operational Graph Dynamics}
\chapter{Operational Graph Dynamics and Graph-Time Resonance}

The graph should not appear only as a static stiffness screen. In a
converter-rich power system, the graph is the spatial support on which power
flow, inertia, damping, controller memory, resonance, fragility, and design
act together. This chapter builds the missing dynamic graph layer. Its central
operator is
\begin{equation}
  T(s;z^\star)
  =
  s^2M+sD+B W(z^\star)B^T+\Pi(s;z^\star),
  \label{eq:graph_dynamic_master_operator}
\end{equation}
where \(B\) is an oriented incidence matrix, \(W(z^\star)\) contains
incremental line stiffnesses at the AC equilibrium \(z^\star\), \(M\) is a
nodal inertia measure, \(D\) is a nodal or device-level dissipative field, and
\(\Pi(s;z^\star)\) is the condensed controller self-energy. This turns the
thesis into a theory of damped, controlled waves on electrical graphs.

\thesisbox{thesisblue}{Graph-dynamic thesis}{
Dynamic stability is a joint property of topology, operating point, spatial
inertia and damping fields, and controller memory. A resonance occurs when
temporal proximity, spatial alignment, port visibility, low damping, and
resolvent amplification coincide.}

\section{Power-Flow Geometry of the Operational Graph}

Consider a connected lossless graph with \(n\) buses, \(m\) lines, and an
arbitrarily oriented incidence matrix \(B\in\R^{n\times m}\). Let
\begin{equation}
  \eta=B^T\theta
  \label{eq:graph_eta}
\end{equation}
be the vector of line angle differences. For line \(e=(i,j)\), define
\begin{equation}
  \kappa_e=b_eV_iV_j>0.
  \label{eq:graph_kappa}
\end{equation}
The oriented active line flow is
\begin{equation}
  f_e(\theta,V)=\kappa_e\sin\eta_e,
  \label{eq:graph_line_flow}
\end{equation}
and the nodal active-power injection map is
\begin{equation}
  p(\theta,V)=B\operatorname{diag}(\kappa)\sin(B^T\theta).
  \label{eq:graph_power_flow_compact}
\end{equation}
At an equilibrium \(z^\star=(\theta^\star,V^\star,\ldots)\),
\begin{equation}
  f_e^\star=\kappa_e^\star\sin\eta_e^\star,
  \qquad
  w_e^\star=\kappa_e^\star\cos\eta_e^\star.
  \label{eq:graph_flow_stiffness_pair}
\end{equation}
Thus the operating-point stiffness is
\begin{equation}
  L^\star=B W^\star B^T,\qquad
  W^\star=\operatorname{diag}(w_e^\star).
  \label{eq:graph_operational_laplacian}
\end{equation}
If \(|\eta_e^\star|<\pi/2\) for all lines, then \(W^\star\succ0\). For a
connected graph, \(L^\star\succeq0\) and
\(\ker L^\star=\operatorname{span}\{\ones\}\). The uniform angle direction is
removed by projection onto \(\ones^\perp\).

The flow-stiffness pair in \eqref{eq:graph_flow_stiffness_pair} is the first
dynamic graph fact:
\begin{equation}
  \frac{\partial w_e}{\partial \eta_e}
  =
  -\kappa_e\sin\eta_e
  =
  -f_e.
  \label{eq:graph_stiffness_angle_derivative}
\end{equation}
A heavily loaded line can lose synchronizing stiffness rapidly under an
additional angular perturbation. This is the direct bridge from power flow to
modal fragility.

\section{Total Derivative of the Operational Laplacian}

Let \(\rho\) be a planning parameter that may change injections,
susceptances, voltages, topology, or the operating point. The reduced
equilibrium equation is
\begin{equation}
  F(\theta,\rho)
  =
  B\operatorname{diag}(\kappa(\rho))\sin(B^T\theta)-p(\rho)=0.
  \label{eq:graph_equilibrium_F}
\end{equation}
Assume the reduced Laplacian \(L\) is invertible on \(\ones^\perp\). Then
\begin{equation}
  \theta_\rho
  =
  L^\dagger
  \left[
  p_\rho-B\operatorname{diag}(\sin\eta)\kappa_\rho
  \right],
  \qquad
  \eta_\rho=B^T\theta_\rho.
  \label{eq:graph_theta_rho}
\end{equation}
Since \(w=\kappa\odot\cos\eta\),
\begin{equation}
  w_\rho
  =
  \cos\eta\odot\kappa_\rho
  -
  f\odot\eta_\rho,
  \label{eq:graph_w_rho}
\end{equation}
and
\begin{equation}
  L_\rho=B\operatorname{diag}(w_\rho)B^T.
  \label{eq:graph_L_rho}
\end{equation}

Equation~\eqref{eq:graph_w_rho} separates the direct line-parameter channel
from the operating-point channel:
\begin{equation}
  L_\rho
  =
  B\operatorname{diag}(\cos\eta\odot\kappa_\rho)B^T
  -
  B\operatorname{diag}(f\odot\eta_\rho)B^T.
  \label{eq:graph_L_rho_split}
\end{equation}
The second term is often the one missed by frozen-graph screens. A physical
line change modifies flows and angles; those changes modify many incremental
line stiffnesses, not only the line being touched.

\section{Mode Sensitivity to Power-Flow Geometry}

For generalized graph modes
\begin{equation}
  L\phi_k=\nu_kM\phi_k,
  \qquad
  \phi_k^TM\phi_\ell=\delta_{k\ell},
  \label{eq:graph_generalized_modes}
\end{equation}
a simple eigenvalue satisfies
\begin{equation}
  (\nu_k)_\rho
  =
  \phi_k^T(L_\rho-\nu_kM_\rho)\phi_k.
  \label{eq:graph_nu_rho_general}
\end{equation}
Let \(a_e\) be the incidence vector of line \(e\), and define the modal line
deformation
\begin{equation}
  g_{e,k}=a_e^T\phi_k.
  \label{eq:graph_modal_line_deformation}
\end{equation}
Then
\begin{equation}
  (\nu_k)_\rho
  =
  \sum_{e=1}^m (w_e)_\rho g_{e,k}^2
  -
  \nu_k\sum_{i=1}^n (M_i)_\rho \phi_k(i)^2.
  \label{eq:graph_nu_rho_edges}
\end{equation}

If only injections change while \(\kappa\) and \(M\) are fixed, then
\(\theta_\rho=L^\dagger p_\rho\), and
\begin{equation}
  (\nu_k)_\rho
  =
  -
  \sum_e f_e g_{e,k}^2 a_e^TL^\dagger p_\rho.
  \label{eq:graph_injection_nu_rho}
\end{equation}
Equivalently,
\begin{equation}
  \nabla_p\nu_k
  =
  -L^\dagger B\operatorname{diag}(f)c_k,
  \qquad
  c_k=(g_{1,k}^2,\ldots,g_{m,k}^2)^T.
  \label{eq:graph_grad_p_nu}
\end{equation}
The flow-weighted modal centrality
\begin{equation}
  \chi_{e,k}^{\rm flow}
  =
  |f_e|\,|g_{e,k}|^2
  \label{eq:graph_flow_modal_centrality}
\end{equation}
flags lines that simultaneously carry high flow and sit across a large jump of
the critical mode. This is more physical than degree, betweenness, or
\(\partial\lambda_2/\partial b_e\) alone.

\section{Cycle and Cut Spaces}

The DC flow induced by \(p\in\ones^\perp\) is
\begin{equation}
  f^\star=W B^T L^\dagger p.
  \label{eq:graph_dc_flow}
\end{equation}
It is the unique solution of
\begin{equation}
  \min_f\ \frac12 f^TW^{-1}f
  \quad
  \text{subject to}\quad Bf=p.
  \label{eq:graph_min_energy_flow}
\end{equation}
The Lagrange multiplier proof is immediate: stationarity gives
\(f=WB^T\lambda\), and the constraint gives \(L\lambda=p\).

If \(C\) spans the cycle space, \(BC=0\), every feasible flow can be written
as \(f=f_0+Cq\). Changing a line weight changes the metric \(W^{-1}\), modifies
cycle redistribution, changes the operating point, and then changes \(L\).
This is the formal bridge between topology control, Braess-like effects, and
modal dynamics.

For a cut \(S\subset\mathcal V\), define
\begin{equation}
  m_S=\sum_{i\in S}M_i,\qquad
  m_{S^c}=\sum_{i\notin S}M_i,\qquad
  w(\partial S)=\sum_{e\in\partial S}w_e.
  \label{eq:graph_cut_quantities}
\end{equation}
The first nonzero generalized eigenvalue satisfies
\begin{equation}
  \nu_2
  \le
  w(\partial S)
  \left(
  \frac1{m_S}+\frac1{m_{S^c}}
  \right)
  \label{eq:graph_inertia_cut_bound}
\end{equation}
for every nontrivial cut. The proof is the Rayleigh quotient evaluated on the
two-area vector that is constant on \(S\) and \(S^c\) and has zero
center-of-inertia mean.

\section{Graph Roughness of Damping-to-Inertia Fields}

Assume \(M=\diag(M_i)\) and \(D=\diag(D_i)\), and define the nodal field
\begin{equation}
  \gamma_i=\frac{D_i}{M_i}.
  \label{eq:graph_gamma_field}
\end{equation}
Then
\begin{equation}
  \Dt=M^{-1/2}DM^{-1/2}=\diag(\gamma_i).
  \label{eq:graph_Dt_gamma}
\end{equation}
For \(i\ne j\),
\begin{equation}
  [\Lt,\Dt]_{ij}
  =
  \Lt_{ij}(\gamma_j-\gamma_i).
  \label{eq:graph_commutator_entry}
\end{equation}
Therefore
\begin{equation}
  \norm{[\Lt,\Dt]}_F^2
  =
  2\sum_{e=(i,j)}
  \frac{w_e^2}{M_iM_j}
  (\gamma_i-\gamma_j)^2.
  \label{eq:graph_roughness_identity}
\end{equation}
This is a Dirichlet energy over an auxiliary graph with weights
\(\widehat w_e=w_e^2/(M_iM_j)\). Within the diagonal nodal damping class, and
for a connected graph,
\begin{equation}
  [\Lt,\Dt]=0
  \quad\Longleftrightarrow\quad
  \gamma_i=\gamma\ \text{for all }i
  \quad\Longleftrightarrow\quad
  D=\gamma M.
  \label{eq:graph_diagonal_commuting_iff}
\end{equation}
Thus non-proportionality is spatial roughness of the damping-to-inertia field.

Let \(\Lt Q=Q\Lambda\) and \(\Gamma=Q^T\Dt Q\). Since
\begin{equation}
  Q^T[\Lt,\Dt]Q=[\Lambda,\Gamma],
  \label{eq:graph_commutator_modal}
\end{equation}
for \(k\ne\ell\),
\begin{equation}
  (\nu_k-\nu_\ell)\Gamma_{k\ell}
  =
  q_k^T[\Lt,\Dt]q_\ell.
  \label{eq:graph_modal_coupling_identity}
\end{equation}
Therefore
\begin{equation}
  |\Gamma_{k\ell}|
  \le
  \frac{\norm{[\Lt,\Dt]}_2}{|\nu_k-\nu_\ell|}
  \le
  \frac{
  \left(
  2\sum_{e=(i,j)}
  \frac{w_e^2}{M_iM_j}(\gamma_i-\gamma_j)^2
  \right)^{1/2}}
  {|\nu_k-\nu_\ell|}.
  \label{eq:graph_modal_mixing_bound}
\end{equation}
Large graph roughness and small spectral separation are the two ingredients of
strong modal mixing. The design reading is direct: smooth the field
\(D_i/M_i\), avoid abrupt local damping-to-inertia jumps across high-stiffness
edges, or escalate from a graph-mode damping screen to the full QEP whenever
\eqref{eq:graph_modal_mixing_bound} is not small in the protected mode band.

\section{Dynamic Roughness of Local Control Fields}

If the damping-channel self-energy is approximately local,
\begin{equation}
  \Sigma(s)=\diag(\sigma_i(s)),
  \qquad
  \gamma_i(s)=\frac{\sigma_i(s)}{M_i},
  \label{eq:graph_dynamic_gamma}
\end{equation}
then, frequency by frequency,
\begin{equation}
  \norm{[\Lt,\widetilde\Sigma(s)]}_F^2
  =
  2\sum_{e=(i,j)}
  \frac{w_e^2}{M_iM_j}
  |\gamma_i(s)-\gamma_j(s)|^2.
  \label{eq:graph_dynamic_roughness}
\end{equation}
The band function
\begin{equation}
  \mathcal E_\gamma(\omega)
  =
  \sum_{e=(i,j)}
  \frac{w_e^2}{M_iM_j}
  |\gamma_i(j\omega)-\gamma_j(j\omega)|^2
  \label{eq:graph_dynamic_roughness_band}
\end{equation}
identifies frequencies where the spatial distribution of control is
incompatible with the electrical topology. If \(\Pi(s)\) is dense, it should
be decomposed into edge-like, shunt-like, and irreducibly nonlocal parts:
\begin{equation}
  \Pi(j\omega)
  =
  B K_\Pi(j\omega)B^T
  +
  \diag(k_\Pi(j\omega))
  +
  \Pi_{\rm nonlocal}(j\omega).
  \label{eq:graph_pi_edge_shunt_nonlocal}
\end{equation}
This is the dynamic extension of the line-shunt decomposition used later for
static perturbations.

\section{Joint Graph-Time Frequency Response}

Normalize the retained operator:
\begin{equation}
  \widetilde T(s)=M^{-1/2}T(s)M^{-1/2}.
  \label{eq:graph_normalized_T}
\end{equation}
Let \(\Lt=Q\Lambda Q^T\). In graph-frequency coordinates,
\begin{equation}
  T_G(s)
  =
  Q^T\widetilde T(s)Q
  =
  s^2I+s\Gamma+\Lambda+\Pi_G(s),
  \label{eq:graph_time_operator}
\end{equation}
where \(\Pi_G(s)=Q^TM^{-1/2}\Pi(s)M^{-1/2}Q\). The graph-time transfer
function is
\begin{equation}
  H_G(\omega)=T_G(j\omega)^{-1}.
  \label{eq:graph_time_transfer}
\end{equation}
The index \(k\) is a spatial graph frequency, while \(\omega\) is temporal
frequency. \(H_{G,kk}(\omega)\) measures amplification of spatial mode \(k\),
and \(H_{G,k\ell}(\omega)\) measures conversion from spatial mode \(\ell\) to
spatial mode \(k\).

Let
\begin{equation}
  d_k(\omega)
  =
  \nu_k-\omega^2
  +j\omega\Gamma_{kk}
  +[\Pi_G(j\omega)]_{kk},
  \label{eq:graph_dk}
\end{equation}
and write
\begin{equation}
  T_G(j\omega)=D_0(\omega)+E(\omega),
  \qquad
  D_0(\omega)=\diag(d_1,\ldots,d_n).
  \label{eq:graph_TG_split}
\end{equation}
If \(\alpha(\omega)=\norm{D_0(\omega)^{-1}E(\omega)}_2<1\), then
\begin{equation}
  \norm{H_G(\omega)-D_0(\omega)^{-1}}_2
  \le
  \frac{\norm{D_0^{-1}}_2^2\norm{E}_2}{1-\alpha(\omega)}.
  \label{eq:graph_cross_resonance_bound}
\end{equation}
To first order,
\begin{equation}
  [H_G(\omega)]_{k\ell}
  \approx
  -\frac{E_{k\ell}(\omega)}{d_k(\omega)d_\ell(\omega)},
  \qquad k\ne\ell.
  \label{eq:graph_cross_resonance_first_order}
\end{equation}
A cross-modal resonance therefore requires temporal amplification in both
modes and a strong off-diagonal channel from damping, self-energy, nonnormal
AC stiffness, or nonlocal control.

\section{Resonance Taxonomy}

The graph-time response separates six mechanisms:
\begin{enumerate}
\item \textbf{Natural resonance:} \(\det T(s)=0\).
\item \textbf{Forced resonance:} \(j\Omega\) lies near a visible pole and the
input shape has strong projection on the corresponding left eigenvector.
\item \textbf{Cross-modal graph resonance:} \(|H_{G,k\ell}(\omega)|\) is large
for \(k\ne\ell\).
\item \textbf{Network-control resonance:} a pole of \(\Pi(s)\) is close to a
network-family pole and has a visible residue.
\item \textbf{Parametric topology resonance:} a time-varying line weight
modulates \(L(t)=L_0+\varepsilon L_1\cos\Omega t\).
\item \textbf{Internal nonlinear resonance:} finite-amplitude dynamics satisfy
\(m\omega_k\approx n\omega_\ell\) and require normal forms, continuation, or
harmonic balance.
\end{enumerate}

For parametric line modulation,
\begin{equation}
  L_1=B\widehat W B^T,\qquad
  c_{k\ell}=\phi_k^TL_1\phi_\ell
  =
  \sum_e \widehat w_e(a_e^T\phi_k)(a_e^T\phi_\ell).
  \label{eq:graph_parametric_coupling}
\end{equation}
The coefficient \(c_{k\ell}\) is the exact spatial bridge from a modulated
edge to the two interacting graph modes.

\section{Dynamic Effective Resistance and Low-Rank Updates}

Let
\begin{equation}
  G(s)=T(s)^{-1}
  \label{eq:graph_green_function}
\end{equation}
on the projected nonzero-angle subspace. For line \(e\), define the dynamic
effective resistance
\begin{equation}
  \mathcal R_e(s)=a_e^TG(s)a_e.
  \label{eq:graph_dynamic_effective_resistance}
\end{equation}
If \(s=0\), \(\Pi(0)=0\), and \(T(0)=L\), this reduces to the static effective
resistance \(a_e^TL^\dagger a_e\). For \(s\ne0\), it contains topology,
operating point, inertia, damping, controller poles, memory, and resolvent
amplification.

For a frozen operating point and a rank-one stiffness update,
\begin{equation}
  T_\Delta(s)=T(s)+\Delta w_ea_ea_e^T,
  \label{eq:graph_rank_one_update}
\end{equation}
the determinant lemma gives
\begin{equation}
  \det T_\Delta(s)
  =
  \det T(s)\left[1+\Delta w_e\mathcal R_e(s)\right].
  \label{eq:graph_det_rank_one}
\end{equation}
New or shifted poles caused by this finite frozen update satisfy
\begin{equation}
  1+\Delta w_e\mathcal R_e(s)=0.
  \label{eq:graph_pole_update_equation}
\end{equation}
The corresponding Green function update is
\begin{equation}
  G_\Delta(s)
  =
  G(s)
  -
  \frac{\Delta w_e\,G(s)a_ea_e^TG(s)}
  {1+\Delta w_e\mathcal R_e(s)}.
  \label{eq:graph_green_update}
\end{equation}

For multiple candidate lines with incidence matrix \(A_C\),
\begin{equation}
  K_C(s)=A_C^TG(s)A_C,
  \label{eq:graph_candidate_kernel}
\end{equation}
and
\begin{equation}
  \det T_\Delta(s)
  =
  \det T(s)\det\!\left[I+\Delta W K_C(s)\right].
  \label{eq:graph_multi_det_update}
\end{equation}
This reduces a large spectral update to a matrix whose dimension is the number
of candidate actions. It is exact for frozen low-rank stiffness updates and a
predictor for physical line changes that also move the operating point.

\section{Dynamic Edge Centrality}

For an input shape \(b\), an output shape \(c\), and
\begin{equation}
  h(s)=c^TG(s)b,
  \label{eq:graph_transfer_scalar}
\end{equation}
the frozen line-weight derivative is
\begin{equation}
  \frac{\partial h(s)}{\partial w_e}
  =
  -
  \left[c^TG(s)a_e\right]
  \left[a_e^TG(s)b\right].
  \label{eq:graph_dynamic_edge_derivative}
\end{equation}
Thus
\begin{equation}
  \mathcal C_e(\omega;b,c)
  =
  |c^TG(j\omega)a_e|\,
  |a_e^TG(j\omega)b|
  \label{eq:graph_dynamic_edge_centrality}
\end{equation}
is a frequency-specific dynamic centrality. It identifies a line that is
reachable from the disturbance, transmits energy toward the protected output,
and matters in the resonant band. It need not have high static degree,
betweenness, or Fiedler sensitivity.

\section{Structured Fragility Radii}

There is no single maximum perturbation. At least four different radii are
needed.

For topology or parameter changes, define a protected spectral region
\(\mathcal S\). With candidate line matrix \(A_C\), the structured topological
radius is
\begin{equation}
  r_{\rm top}
  =
  \inf\left\{
  \norm{\Delta W}:
  \exists s\in\partial\mathcal S,\ 
  \det[T(s)+A_C\Delta W A_C^T]=0
  \right\}.
  \label{eq:graph_topological_radius}
\end{equation}
Using \eqref{eq:graph_multi_det_update}, the reduced condition is
\begin{equation}
  \det[I+\Delta W K_C(s)]=0.
  \label{eq:graph_reduced_fragility_condition}
\end{equation}
With a declared physical uncertainty structure, this becomes the structured
radius
\begin{equation}
  r_{\rm top}
  =
  \left[
  \sup_{s\in\partial\mathcal S}
  \mu_\Delta(K_C(s))
  \right]^{-1}.
  \label{eq:graph_mu_radius}
\end{equation}

For a single frozen line update, the complex unstructured line radius along a
protected boundary is especially transparent:
\begin{equation}
  r_{e}^{\C}
  =
  \left[
  \sup_{s\in\partial\mathcal S}
  |\mathcal R_e(s)|
  \right]^{-1}.
  \label{eq:graph_complex_line_radius}
\end{equation}
This is the smallest complex scalar \(\Delta w_e\) that can satisfy
\(1+\Delta w_e\mathcal R_e(s)=0\) for some boundary point. A physical line
weight is real, so the real frozen-line radius is
\begin{equation}
  r_{e}^{\R}
  =
  \min_{\substack{s\in\partial\mathcal S\\
  \operatorname{Im}\mathcal R_e(s)=0}}
  \frac{1}{|\operatorname{Re}\mathcal R_e(s)|},
  \label{eq:graph_real_line_radius}
\end{equation}
with the convention \(r_e^{\R}=\infty\) if the feasibility set is empty. Thus
dynamic effective resistance is not only a ranking score; it gives an exact
single-line fragility radius for the declared spectral boundary.

For a forced input in a stable linear model, increasing amplitude does not move
poles. It can still violate an operational limit:
\begin{equation}
  |\widehat u|_{\max}^{\Omega}
  =
  \frac{y_{\max}}
  {\sup_{\omega\in\Omega}
  \sigma_{\max}\!\left(C_yG(j\omega)B_u\right)}.
  \label{eq:graph_forced_amplitude_radius}
\end{equation}
For nonlinear loss of stability, the relevant object is a region-of-attraction
radius,
\begin{equation}
  r_{\rm ROA}
  =
  \inf_{x\in\partial\mathcal B(x^\star)}
  \norm{W_x(x-x^\star)}.
  \label{eq:graph_roa_radius}
\end{equation}
The first is a parametric fragility radius, the second is an input-output
exposure bound, and the third is a nonlinear stability quantity.

\section{Stability Recovery Design}

The graph-dynamic design problem combines line, damping, inertia, and
controller actions:
\begin{equation}
  u=
  \begin{bmatrix}
  \Delta w & \Delta D & \Delta M & \Delta\kappa
  \end{bmatrix}^T.
  \label{eq:graph_design_vector}
\end{equation}
A unified recovery problem is
\begin{align}
  \min_u\quad&
  c_w^T|\Delta w|
  +c_D^T|\Delta D|
  +c_M^T|\Delta M|
  +c_\kappa^T|\Delta\kappa|
  \label{eq:graph_recovery_problem}\\
  \text{subject to}\quad&
  F(z^\star,u)=0,\nonumber\\
  &|f_e(z^\star,u)|\le \overline f_e,\nonumber\\
  &|B^T\theta^\star|\le \eta_{\max},\nonumber\\
  &\Real s_k(u)\le-\alpha_\star,\qquad
    \zeta_k(u)\ge\zeta_\star,\nonumber\\
  &\sup_{\omega\in\Omega}
    \sigma_{\max}\!\left(C_yT(j\omega,u)^{-1}B_u\right)
    \le g_\star,\nonumber\\
  &r_{\rm top}(u)\ge r_\star.\nonumber
\end{align}
Sequential versions should use pole sensitivities and resolvent-gradient
steps only inside a trust region, followed by continuation of the operating
point and the protected poles.

\section{AC, Directed, and Matrix-Weighted Extensions}

The scalar lossless Laplacian is the symmetric floor, not the universal model.
For the AC Jacobian,
\begin{equation}
  \begin{bmatrix}
  \Delta P\\
  \Delta Q
  \end{bmatrix}
  =
  \begin{bmatrix}
  J_{P\theta}&J_{PV}\\
  J_{Q\theta}&J_{QV}
  \end{bmatrix}
  \begin{bmatrix}
  \Delta\theta\\
  \Delta V
  \end{bmatrix},
  \label{eq:graph_ac_jacobian}
\end{equation}
eliminating voltage increments under \(\Delta Q=0\) gives
\begin{equation}
  L_{\rm AC}
  =
  J_{P\theta}-J_{PV}J_{QV}^{-1}J_{Q\theta}.
  \label{eq:graph_ac_condensed_L}
\end{equation}
In general \(L_{\rm AC}\ne L_{\rm AC}^T\). Its symmetric part is
stiffness-like; its skew part is circulatory and nonconservative. This
nonnormal layer requires left-right eigenvectors and resolvent gains.

For multivariable ports, assign each bus an \(r\)-dimensional vector and set
\(\mathcal B=B\otimes I_r\). With line admittances
\(Y_e(s)\in\C^{r\times r}\),
\begin{equation}
  \mathcal L(s)
  =
  \mathcal B
  \operatorname{blkdiag}(Y_1(s),\ldots,Y_m(s))
  \mathcal B^T.
  \label{eq:graph_matrix_laplacian}
\end{equation}
For a matrix update on line \(e\), with \(\mathcal A_e=a_e\otimes I_r\),
\begin{equation}
  \mathcal R_e^{(r)}(s)
  =
  \mathcal A_e^T\mathcal T(s)^{-1}\mathcal A_e,
  \label{eq:graph_matrix_effective_resistance}
\end{equation}
and
\begin{equation}
  \det\mathcal T_\Delta(s)
  =
  \det\mathcal T(s)
  \det\!\left[I_r+\Delta Y_e(s)\mathcal R_e^{(r)}(s)\right].
  \label{eq:graph_matrix_det_update}
\end{equation}
This is the natural extension to \(dq\) coupling, voltage-frequency channels,
frequency-dependent line models, directed effects, and multivariable converter
ports.

\section{Paper Extraction}

This chapter supports a focused Transactions-style paper:
\begin{quote}
\itshape
Graph-Time Resonance, Spectral Effective Resistance, and Structured Topology
Fragility in Converter-Rich Power Systems.
\end{quote}
The broad dissertation contribution is operator-valued dynamic grid strength.
The narrower article should be sold as controller-resolvent-aware spectral
graph planning. Its minimum publishable package is:
\begin{enumerate}
\item graph roughness of \(D_i/M_i\) and the modal-mixing bound
\eqref{eq:graph_modal_mixing_bound};
\item graph-time transfer \(H_G(\omega)\) and cross-resonance atlas;
\item finite dynamic-effective-resistance updates
\(1+\Delta w_e\mathcal R_e(s)=0\);
\item structured topology-fragility radii
\eqref{eq:graph_complex_line_radius}--\eqref{eq:graph_real_line_radius};
\item a controller-induced reinforcement-reversal gate, using the Schur
self-energy machinery developed later in the thesis.
\end{enumerate}
The novelty claim should not be that Laplacians or effective resistance are
new, that Schur complements are new, or that adding a line can never hurt in
existing literature. The defensible claim is narrower and stronger:
power-flow-induced stiffness, damping-to-inertia roughness, graph-time modal
conversion, finite low-rank topology updates, rational controller self-energy,
and structured stability radii are projections of one operator
\(T(s;z^\star)\), and that operator determines when a topological improvement
really improves a protected dynamic margin.


\part{Constant Damping Theory}
\chapter{Fixed-Damping Spectral Theory}

The first level of the hierarchy assumes that damping is a fixed matrix. This
is the last level at which one can still try to extract a large amount of
engineering information from the graph Laplacian \(L\), the inertia matrix
\(M\), and the damping matrix \(D\) without retaining converter-control
states. The chapter develops that analysis aggressively. It derives the
spectral screens, the exact first-order eigenvalue sensitivities, the
decoupled modal formulas for \(d/dL\), \(d/dD\), and \(d/dM\), and the limits
of graph-only reinforcement.

The conclusion is deliberately sharp. With fixed \(D\), adding inertia is not
automatically stabilizing for oscillatory damping margins. Adding a line is not
automatically stabilizing either. Damping ratio is a quotient: it improves when
modal damping rises faster than modal frequency. If an action raises stiffness
or inertia but does not raise effective modal damping, the damping ratio can
fall even while a graph metric improves.

\section{The Classical Swing Operator}

Consider
\begin{equation}
  M\ddot\theta+D\dot\theta+L\theta=p,
  \label{eq:fixed_swing}
\end{equation}
with \(M\succ0\), \(D\succeq0\), and \(L=L^T\succeq0\) on the projected
subspace orthogonal to uniform angle shift. The homogeneous modal ansatz
\(\theta(t)=ve^{st}\) gives the QEP
\begin{equation}
  P(s)v=(s^2M+sD+L)v=0.
  \label{eq:fixed_qep}
\end{equation}
This is the basic fixed-\(D\) operator. All planning screens in this chapter
are approximations or consequences of \eqref{eq:fixed_qep}.

In inertia-normalized coordinates,
\begin{equation}
  \ddot\vartheta+\Dt\dot\vartheta+\Lt\vartheta=0,
  \label{eq:fixed_normalized}
\end{equation}
where
\begin{equation}
  \Lt=M^{-1/2}LM^{-1/2},\qquad
  \Dt=M^{-1/2}DM^{-1/2}.
\end{equation}
Let \(\Lt Q=Q\Lambda\). In modal coordinates \(z=Q^T\vartheta\),
\begin{equation}
  \ddot z+\Gamma \dot z+\Lambda z=0,\qquad
  \Gamma=Q^T\Dt Q.
  \label{eq:fixed_modal}
\end{equation}
The graph diagonalizes stiffness. It diagonalizes damping only when
\begin{equation}
  \Lt\Dt=\Dt\Lt
  \label{eq:commutation_condition}
\end{equation}
on the nonzero modal subspace. Equation~\eqref{eq:commutation_condition} is the
first hidden assumption behind graph-only damping arguments.

A basis-dependent off-diagonal ratio is useful only when the eigenvalues of
\(\Lt\) are well separated. For clustered eigenvalues, the basis inside the
cluster can rotate without changing the physics. A coordinate-free global gate is
\begin{equation}
  \eta_{\rm comm}
  =
  \frac{\norm{\Lt\Dt-\Dt\Lt}_F}
  {2\norm{\Lt}_F\norm{\Dt}_F+\varepsilon_{\rm reg}}.
  \label{eq:commutator_nonprop_index}
\end{equation}
For two spectral clusters \(I\) and \(J\), with orthogonal projectors
\(P_I\) and \(P_J\), the cross-cluster damping gate is
\begin{equation}
  \eta_{IJ}
  =
  \frac{\norm{P_I\Dt P_J}_2}
  {\operatorname{sep}(\Lambda_I,\Lambda_J)+\varepsilon_{\rm reg}}.
  \label{eq:projector_cluster_nonprop}
\end{equation}
This form is invariant to rotations of the modal basis inside each cluster and
is the safer trigger for adaptive reduced-QEP construction.

\section{Exact QEP Energy Identities}

The scalar decoupled formulas used later as screens have an exact ancestor in
the full QEP. Let \(P(s)x=0\) for
\[
  P(s)=s^2M+sD+L,
  \qquad
  M\succ0,\quad D\succeq0,\quad L\succ0
\]
on the nonzero-angle subspace. Assume \(s=\sigma+j\omega\) with
\(\omega\ne0\). Multiplying by \(x^H\) gives
\begin{equation}
  s^2 m+s d+\ell=0,
  \qquad
  m=x^HMx,\quad d=x^HDx,\quad \ell=x^HLx.
  \label{eq:qep_scalar_energy_identity}
\end{equation}
Since \(M,D,L\) are Hermitian in this fixed-\(D\) layer, \(m,d,\ell\) are real
and \(m>0\). Separating real and imaginary parts yields
\begin{align}
  \sigma
  &=
  -\frac{x^HDx}{2x^HMx}, \label{eq:qep_exact_sigma}\\
  |s|^2
  &=
  \frac{x^HLx}{x^HMx}, \label{eq:qep_exact_modulus}\\
  \zeta
  &=
  \frac{x^HDx}
  {2\sqrt{(x^HMx)(x^HLx)}}. \label{eq:qep_exact_zeta}
\end{align}
These identities do not require \(D=\gamma M\). They require the fixed
Hermitian QEP structure and an oscillatory root.

\begin{proposition}[Absence of undamped oscillatory modes]
For the fixed Hermitian QEP, an oscillatory eigenvalue on the imaginary axis
\(s=j\omega\), \(\omega\ne0\), exists only if there is a nonzero vector \(x\)
such that
\begin{equation}
  Dx=0,\qquad Lx=\omega^2Mx.
  \label{eq:undamped_mode_condition}
\end{equation}
\end{proposition}

\begin{proof}
If \(s=j\omega\), then \eqref{eq:qep_exact_sigma} gives \(x^HDx=0\). Since
\(D\succeq0\), this implies \(Dx=0\). The QEP then reduces to
\((-\,\omega^2M+L)x=0\). The converse follows by substitution.
\end{proof}

\section{Generalized Graph Modes}

For placement calculations it is often cleaner to use the generalized
eigenproblem
\begin{equation}
  L\phi_k=\nu_k M\phi_k,\qquad
  \phi_k^T M\phi_\ell=\delta_{k\ell}.
  \label{eq:generalized_modes}
\end{equation}
The zero mode is removed. The scalar modal stiffness is \(\nu_k\), and the
frozen-mode scalar damping is
\begin{equation}
  \delta_k=\phi_k^TD\phi_k.
  \label{eq:modal_delta_physical}
\end{equation}
If the modes are decoupled, mode \(k\) is governed by
\begin{equation}
  s^2+\delta_k s+\nu_k=0.
  \label{eq:fixed_scalar_generalized}
\end{equation}
The damping ratio is
\begin{equation}
  \zeta_k=\frac{\delta_k}{2\sqrt{\nu_k}}.
  \label{eq:fixed_zeta_generalized}
\end{equation}
This formula is a screen outside the decoupled case. It is still extremely
useful because it exposes the competing effects:
\begin{equation}
  \frac{d\zeta_k}{\zeta_k}
  =
  \frac{d\delta_k}{\delta_k}
  -
  \frac{1}{2}\frac{d\nu_k}{\nu_k}.
  \label{eq:zeta_log_differential}
\end{equation}
An action improves damping ratio only if it raises modal damping faster than
half the relative increase in modal stiffness.

\section{Uniform-Decay Case \texorpdfstring{\(D=\gamma M\)}{D=gamma M}}

The most restrictive decoupled case is
\begin{equation}
  D=\gamma M,\qquad \gamma>0.
  \label{eq:proportional_damping}
\end{equation}
Then \(\delta_k=\gamma\), \(\Gamma=\gamma I\), and every nonzero graph mode
satisfies
\begin{equation}
  s^2+\gamma s+\nu_k=0.
  \label{eq:fixed_scalar_mode}
\end{equation}
The roots are
\begin{equation}
  s_{k,\pm}=\frac{-\gamma\pm\sqrt{\gamma^2-4\nu_k}}{2}.
  \label{eq:fixed_proportional_roots}
\end{equation}
When \(\nu_k>\gamma^2/4\), the roots are underdamped and
\begin{equation}
  \Real(s_{k,\pm})=-\gamma/2.
  \label{eq:fixed_decay_floor}
\end{equation}

\begin{proposition}[Proportional-damping floor]
Under \(D=\gamma M\), every underdamped nonzero mode has the same decay rate
\(-\gamma/2\). Network topology changes \(\nu_k\) and therefore modal
frequency, but it does not change the real decay rate of an underdamped mode
inside the proportional class.
\end{proposition}

\begin{proof}
Equation~\eqref{eq:fixed_scalar_mode} gives the two roots directly. If
\(\nu_k>\gamma^2/4\), the square root is imaginary and the real part is
\(-\gamma/2\). The expression contains no dependence on \(\nu_k\) in its real
part.
\end{proof}

The proposition is sometimes counterintuitive because it does not say topology
is irrelevant. Topology can move a mode between overdamped and underdamped
regions, and it changes frequency. It states only that, after a mode is
underdamped, its decay rate is fixed by \(\gamma\) in this uniform
damping-to-inertia model.

\section{Commuting and Generalized Proportional Damping}

Uniform damping is not the whole decouplable class. The exact condition is
\begin{equation}
  [\Lt,\Dt]=0.
  \label{eq:fixed_commuting_decoupling}
\end{equation}
For real symmetric matrices, \eqref{eq:fixed_commuting_decoupling} is
equivalent to simultaneous orthogonal diagonalization. Therefore there exists
an orthogonal \(Q\) such that
\begin{equation}
  Q^T\Lt Q=\diag(\nu_k),
  \qquad
  Q^T\Dt Q=\diag(\delta_k).
  \label{eq:fixed_simultaneous_diagonalization}
\end{equation}
Each mode then satisfies
\begin{equation}
  s_k^2+\delta_k s_k+\nu_k=0,
  \qquad
  \zeta_k=\frac{\delta_k}{2\sqrt{\nu_k}}.
  \label{eq:fixed_commuting_scalar_modes}
\end{equation}

Rayleigh damping is the most important generalized proportional example:
\begin{equation}
  D=\alpha M+\beta L.
  \label{eq:rayleigh_damping}
\end{equation}
In normalized coordinates,
\begin{equation}
  \Dt=\alpha I+\beta\Lt,
  \label{eq:rayleigh_normalized}
\end{equation}
so \([\Lt,\Dt]=0\). Mode \(k\) satisfies
\begin{equation}
  s_k^2+(\alpha+\beta\nu_k)s_k+\nu_k=0,
  \label{eq:rayleigh_mode_polynomial}
\end{equation}
and its underdamped real part and damping ratio are
\begin{align}
  \Real(s_k)&=-\frac{\alpha+\beta\nu_k}{2},
  \label{eq:rayleigh_real_part}\\
  \zeta_k&=
  \frac{\alpha}{2\sqrt{\nu_k}}
  +
  \frac{\beta\sqrt{\nu_k}}{2}.
  \label{eq:rayleigh_zeta}
\end{align}
Unlike the \(D=\gamma M\) floor, Rayleigh damping lets network stiffness change
both frequency and decay. Depending on the modal range and on the signs and
magnitudes of \(\alpha\) and \(\beta\), reinforcement can improve one set of
modes while degrading another damping-ratio target. This is why the thesis
uses the commuting condition as the decoupling test and reserves
\(D=\gamma M\) for the uniform-decay case.

\section{Exact QEP Sensitivities}

Let a scalar action \(\rho\) perturb all three operators:
\begin{equation}
  P(s,\rho)=s^2M(\rho)+sD(\rho)+L(\rho).
\end{equation}
For a simple eigenvalue \(s_\star\), right vector \(v\), and left vector \(w\),
the QEP sensitivity is
\begin{equation}
  \frac{ds_\star}{d\rho}
  =
  -\frac{
  w^H\left(s_\star^2M_\rho+s_\star D_\rho+L_\rho\right)v
  }{
  w^H\left(2s_\star M+D\right)v
  }.
  \label{eq:fixed_full_qep_sens}
\end{equation}
This formula, standard in QEP perturbation theory \cite{tisseur2001qep}, is
the most compact answer to how \(d/dL\), \(d/dD\), and \(d/dM\) mix.

The three numerator terms have different phase factors:
\begin{equation}
  L_\rho,\qquad s_\star D_\rho,\qquad s_\star^2M_\rho.
  \label{eq:phase_weighted_terms}
\end{equation}
For an oscillatory pole, those terms are not collinear in the complex plane.
Thus, line reinforcement, damping addition, and inertia addition do not point
in the same eigenvalue direction. A scalar ranking that ignores these phase
weights is not a damping certificate.

For \(s=\sigma+j\omega\), the damping-ratio derivative is
\begin{equation}
  \frac{d\zeta}{d\rho}
  =
  -\frac{\omega^2}{(\sigma^2+\omega^2)^{3/2}}
  \Real\left(\frac{ds}{d\rho}\right)
  +
  \frac{\sigma\omega}{(\sigma^2+\omega^2)^{3/2}}
  \Imag\left(\frac{ds}{d\rho}\right).
  \label{eq:fixed_zeta_sensitivity}
\end{equation}
The real-part movement alone is insufficient when the frequency movement is
large.

\section{Derivative with Respect to the Graph \texorpdfstring{\(L\)}{L}}

Let a candidate line \(e=(i,j)\) have incidence vector \(a_e=e_i-e_j\) and
susceptance-like weight \(b_e\). In the Laplacian approximation,
\begin{equation}
  L=\sum_{e} b_e a_ea_e^T,\qquad
  \frac{\partial L}{\partial b_e}=a_ea_e^T.
  \label{eq:laplacian_edge_sum}
\end{equation}
For the generalized graph mode \eqref{eq:generalized_modes},
\begin{equation}
  \frac{\partial \nu_k}{\partial b_e}
  =
  (\phi_k(i)-\phi_k(j))^2.
  \label{eq:edge_spectral_derivative}
\end{equation}
Equation~\eqref{eq:edge_spectral_derivative} is the core weak-link screen:
edges across large modal angle differences have the strongest effect on modal
frequency.

For the scalar decoupled damping-ratio screen with fixed \(D\),
\begin{equation}
  \frac{\partial \zeta_k}{\partial b_e}
  =
  -\frac{\zeta_k}{2\nu_k}
  (\phi_k(i)-\phi_k(j))^2
  +
  \frac{\zeta_k}{\delta_k}
  \frac{\partial \delta_k}{\partial b_e}.
  \label{eq:zeta_edge_derivative_general}
\end{equation}
If modal damping is frozen under the line action, the second term is zero and
line reinforcement reduces \(\zeta_k\) for that mode. It may still improve
angle coherence, synchronizing torque, or protected resolvent gain, but it does
not improve scalar damping ratio by itself.

\begin{example}[A reinforced cut can hide a damping loss]
Suppose a critical mode has \(\delta_k\) fixed and a line across its modal cut
has \((\phi_k(i)-\phi_k(j))^2\) large. Reinforcing that line raises
\(\nu_k\), which is good for graph stiffness. Yet
\eqref{eq:zeta_edge_derivative_general} gives
\(\partial\zeta_k/\partial b_e<0\) if \(\partial\delta_k/\partial b_e=0\).
The graph has become stiffer, but the damping ratio has decreased. This is a
hidden margin failure.
\end{example}

\section{Derivative with Respect to Damping \texorpdfstring{\(D\)}{D}}

For a local damping addition at bus \(i\),
\begin{equation}
  D_\rho=e_ie_i^T.
\end{equation}
The frozen-mode damping derivative is
\begin{equation}
  \frac{\partial\delta_k}{\partial D_i}=\phi_k(i)^2.
  \label{eq:damping_node_derivative}
\end{equation}
Therefore
\begin{equation}
  \frac{\partial\zeta_k}{\partial D_i}
  =
  \frac{\phi_k(i)^2}{2\sqrt{\nu_k}}
  \label{eq:zeta_damping_derivative}
\end{equation}
in the decoupled scalar screen. The best damping location for a single mode is
an antinode of that mode, not necessarily the largest generator, the largest
load, or the highest-degree bus.

For multiple modes, a risk-weighted damping placement score is
\begin{equation}
  \mathcal D_i
  =
  \sum_{k\in\mathcal K}
  \omega_k
  \frac{\phi_k(i)^2}{2\sqrt{\nu_k}},
  \label{eq:damping_placement_score}
\end{equation}
where \(\mathcal K\) is the set of monitored modes and \(\omega_k\ge0\) weights
criticality. This is a screen. The exact acceptance test remains
\eqref{eq:fixed_full_qep_sens} or full eigenvalue recomputation.

\begin{example}[A low-degree bus can be the best damping site]
A bus with modest graph degree can have large \(\phi_k(i)^2\) for the critical
mode. Equation~\eqref{eq:zeta_damping_derivative} ranks it highly for damping
placement even if graph centrality ranks it poorly. This is a weak-node
mechanism: the bus is weak because the critical mode lives there.
\end{example}

\section{Derivative with Respect to Inertia \texorpdfstring{\(M\)}{M}}

For a local inertia addition at bus \(i\),
\begin{equation}
  M_\rho=e_ie_i^T.
\end{equation}
The generalized eigenvalue derivative is
\begin{equation}
  \frac{\partial\nu_k}{\partial M_i}
  =
  -\nu_k\phi_k(i)^2.
  \label{eq:inertia_frequency_derivative}
\end{equation}
Local inertia lowers the modal frequency most strongly where the mode has large
amplitude. That fact supports inertia placement for frequency nadir or RoCoF
objectives, and it is consistent with studies showing that virtual-inertia
placement can be posed as a network performance problem
\cite{poolla2017virtual}, that inertia location and slow modes shape
disturbance propagation \cite{pagnier2019inertia_location}, and that inertia
and primary-control placement separate under heterogeneous parameters
\cite{pagnier2019optimal}.

For oscillatory damping ratio, however, the conclusion is more delicate. The
exact answer is still the QEP derivative \eqref{eq:fixed_full_qep_sens}. As a
fixed-mode planning screen, if physical damping \(D\) is not co-added, the
frozen-shape normalization gives
\begin{equation}
  \frac{\partial\delta_k}{\partial M_i}
  \approx
  -\delta_k\phi_k(i)^2.
  \label{eq:inertia_damping_dilution}
\end{equation}
Combining \eqref{eq:zeta_log_differential},
\eqref{eq:inertia_frequency_derivative}, and
\eqref{eq:inertia_damping_dilution} yields
\begin{equation}
  \frac{\partial\zeta_k}{\partial M_i}
  \approx
  -\frac{1}{2}\zeta_k\phi_k(i)^2.
  \label{eq:inertia_zeta_dilution}
\end{equation}
Thus, in a fixed-\(D\) model, inertia-only placement at a modal antinode can
decrease the scalar damping-ratio screen. It may still improve RoCoF or reduce
frequency peaks, but it is not a free damping improvement. Any candidate that
matters operationally should therefore be rechecked with
\eqref{eq:fixed_full_qep_sens} or with a full eigenvalue recomputation.

If inertia is co-designed with proportional damping,
\begin{equation}
  D_i=\gamma M_i,
  \label{eq:co_designed_proportional}
\end{equation}
then the damping-to-inertia ratio is preserved. The real decay floor remains
\(-\gamma/2\), and the damping ratio increases as \(\nu_k\) decreases. This is
the engineering difference between adding inertia alone and adding a
properly-damped inertial device.

\begin{example}[Inertia placement has two different objectives]
Placing inertia where \(\phi_k(i)^2\) is large is attractive if the objective
is to slow a disturbance. Under fixed \(D\), the same placement can reduce
\(\zeta_k\) through damping dilution. A correct planning statement must say
whether the target is RoCoF, frequency nadir, decay rate, damping ratio, or
input-output amplification.
\end{example}

\section{Graph Robustness Without Adding Inertia}

If inertia is unavailable or undesirable, the graph can be strengthened through
line weights, switching, topology control, or series compensation. The oldest
spectral screen is algebraic connectivity: maximize \(\nu_2\), the first
nonzero Laplacian eigenvalue introduced by Fiedler's graph connectivity theory
\cite{fiedler1973algebraic}. The local derivative is
\begin{equation}
  \frac{\partial\nu_2}{\partial b_e}
  =
  (\phi_2(i)-\phi_2(j))^2.
  \label{eq:fiedler_edge_score}
\end{equation}
This favors lines across the Fiedler cut.

A different robustness screen is total effective resistance, also called the
Kirchhoff index:
\begin{equation}
  K_f=n\,\mathrm{trace}(L^+)
  =
  n\sum_{k=2}^n\frac{1}{\lambda_k}.
  \label{eq:kirchhoff_index}
\end{equation}
Minimizing effective resistance is a convex edge-weight allocation problem in
standard graph settings \cite{ghosh2008effective}. Its derivative is
\begin{equation}
  \frac{\partial\,\mathrm{trace}(L^+)}{\partial b_e}
  =
  -a_e^T(L^+)^2a_e
  =
  -\sum_{k=2}^n
  \frac{(q_k(i)-q_k(j))^2}{\lambda_k^2}.
  \label{eq:effective_resistance_derivative}
\end{equation}
This favors lines that reduce many slow modes, not only the Fiedler mode.
Effective resistance also underlies spectral sparsification, where edges are
sampled or weighted according to their resistance contribution while preserving
quadratic forms \cite{spielman2011sparsification}.

The hidden damping question remains:
\begin{equation}
  \hbox{Does the graph action also improve }
  \zeta_k,\ \Real(s_k),\hbox{ or }g_{\rm prot}?
  \label{eq:graph_hidden_question}
\end{equation}
If not, the action is graph-robust but damping-ambiguous. It must be paired
with damping, controller retuning, or a certificate that the chosen margin is
not damping ratio.

\section{Weak Nodes, Weak Links, and Hidden Margins}

The fixed-\(D\) theory gives precise first-order definitions.

\begin{definition}[Weak node for damping]
For a monitored mode set \(\mathcal K\), bus \(i\) is weak for damping if its
modal damping placement score \(\mathcal D_i\) in
\eqref{eq:damping_placement_score} is large while the current local damping
contribution \(D_i\phi_k(i)^2\) is small for critical modes.
\end{definition}

\begin{definition}[Weak link for graph stiffness]
Line \(e\) is weak for a graph metric if the derivative of that metric with
respect to \(b_e\) is large. Examples are
\eqref{eq:fiedler_edge_score} for algebraic connectivity and
\eqref{eq:effective_resistance_derivative} for total effective resistance.
\end{definition}

\begin{definition}[Hidden damping margin]
A graph action has hidden damping risk for mode \(k\) if
\begin{equation}
  \frac{\partial \nu_k}{\partial\rho}>0
  \quad\hbox{but}\quad
  \frac{\partial \zeta_k}{\partial\rho}<0
  \label{eq:hidden_margin_definition}
\end{equation}
under the selected damping model. It improves a stiffness screen while
degrading a damping-ratio screen.
\end{definition}

This definition prevents a common overclaim: a stronger graph is not the same
as a better damping margin.

\section{A Fixed-\texorpdfstring{\(D\)}{D} Spectral Planning Screen}

The following screen extracts the maximum information available before moving
to the full QEP or Schur NEP.

\begin{enumerate}
\item Compute generalized modes \(L\phi_k=\nu_kM\phi_k\).
\item Compute modal damping \(\delta_k=\phi_k^TD\phi_k\) and
\(\zeta_k=\delta_k/(2\sqrt{\nu_k})\).
\item Compute node damping scores \(\mathcal D_i\).
\item Compute inertia-dilution scores
\begin{equation}
  \mathcal M_i
  =
  \sum_{k\in\mathcal K}
  \omega_k\zeta_k\phi_k(i)^2.
  \label{eq:inertia_dilution_score}
\end{equation}
\item Compute line graph scores using \(\partial\nu_2/\partial b_e\) and
\(\partial\mathrm{trace}(L^+)/\partial b_e\).
\item Flag hidden-margin conflicts where graph scores are positive but
\(\partial\zeta_k/\partial b_e<0\).
\item Verify accepted candidates with the exact QEP derivative
\eqref{eq:fixed_full_qep_sens}.
\end{enumerate}

This screen is intentionally not the final thesis method. It is the strongest
fixed-\(D\) layer before controller self-energy is introduced.

\chapter{Fixed Non-Commuting Damping and Modal Coupling}

The second level of the hierarchy keeps the damping matrix fixed but drops the
commutation condition \([\Lt,\Dt]=0\). The eigenvalue problem remains a QEP,
but graph modes are no longer independent. This chapter develops the resulting
modal coupling and shows that it already has the algebraic form of a rational
correction. The term non-proportional is used only in this precise sense:
damping is non-proportional for the present modal analysis when it is not
simultaneously diagonalizable with stiffness in inertia-normalized
coordinates.

\section{Modal Damping Matrix}

In graph modal coordinates,
\begin{equation}
  \ddot z+\Gamma\dot z+\Lambda z=0,\qquad
  \Gamma=Q^T\Dt Q.
  \label{eq:nonprop_modal_system}
\end{equation}
Write
\begin{equation}
  \Gamma=\Gamma_0+E,\qquad
  \Gamma_0=\diag(\Gamma_{11},\ldots,\Gamma_{nn}),
  \label{eq:nonprop_split}
\end{equation}
where \(E\) contains the off-diagonal entries. If \(E=0\), modes are decoupled.
If \(E\ne0\), the velocity of one graph mode enters the equation of another.
Equivalently, for a simple spectrum of \(\Lt\), nonzero off-diagonal entries
of \(\Gamma\) are the modal expression of the commutator \([\Lt,\Dt]\).

The off-diagonal entries are not numerical noise. They are the representation
of damping heterogeneity in the stiffness eigenbasis. A bus with large local
damping can still be ineffective for a critical mode if that mode has small
shape at the bus. Conversely, a small local damping change can matter if it
projects strongly onto the left-right structure of a sensitive eigenpair.

\section{Two-Mode Illustration}

For two graph modes, the QEP is
\begin{equation}
  \begin{bmatrix}
    s^2+s\gamma_1+\nu_1 & s\epsilon\\
    s\epsilon & s^2+s\gamma_2+\nu_2
  \end{bmatrix}
  \begin{bmatrix} z_1\\z_2\end{bmatrix}=0.
  \label{eq:two_mode_nonprop}
\end{equation}
The characteristic equation is
\begin{equation}
  (s^2+s\gamma_1+\nu_1)(s^2+s\gamma_2+\nu_2)-s^2\epsilon^2=0.
  \label{eq:two_mode_char}
\end{equation}
Eliminating \(z_2\) yields an exact scalar nonlinear equation for \(z_1\):
\begin{equation}
  s^2+s\gamma_1+\nu_1
  -
  \frac{s^2\epsilon^2}{s^2+s\gamma_2+\nu_2}
  =0.
  \label{eq:two_mode_rational}
\end{equation}
The correction is rational in \(s\). Even before converter controls are
introduced, eliminating coupled modal coordinates converts a fixed-coefficient
QEP into a frequency-dependent scalar equation for the retained mode.

\section{General Modal Correction}

Partition the modal coordinates into one target mode \(k\) and the remaining
modes \(\bar k\). The modal QEP can be written as
\begin{equation}
  \begin{bmatrix}
    p_k(s) & sE_{k\bar k}\\
    sE_{\bar k k} &
    s^2I+s\Gamma_{\bar k\bar k}+\Lambda_{\bar k\bar k}
  \end{bmatrix}
  \begin{bmatrix}
    z_k\\z_{\bar k}
  \end{bmatrix}=0,
  \label{eq:general_modal_partition}
\end{equation}
where
\begin{equation}
  p_k(s)=s^2+s\Gamma_{kk}+\nu_k.
\end{equation}
Eliminating \(z_{\bar k}\) gives
\begin{equation}
  p_k(s)
  -
  s^2E_{k\bar k}
  \left(s^2I+s\Gamma_{\bar k\bar k}+\Lambda_{\bar k\bar k}\right)^{-1}
  E_{\bar k k}
  =0.
  \label{eq:general_modal_correction}
\end{equation}
This expression shows three facts.

\begin{enumerate}
\item Off-diagonal damping affects a mode through products of coupling terms,
not through diagonal damping alone.
\item Nearby modal denominators matter more than distant denominators.
\item The correction has the same resolvent structure that later appears in
controller condensation.
\end{enumerate}

\section{Perturbation View}

For small \(E\), one may evaluate the correction near an unperturbed root
\(s_{k,0}\) of \(p_k(s)\). The leading eigenvalue shift is governed by
\begin{equation}
  \Delta s_k
  \approx
  \frac{
  s_{k,0}^2E_{k\bar k}
  \left(s_{k,0}^2I+s_{k,0}\Gamma_{\bar k\bar k}
  +\Lambda_{\bar k\bar k}\right)^{-1}
  E_{\bar k k}
  }{
  2s_{k,0}+\Gamma_{kk}
  }.
  \label{eq:nonprop_shift}
\end{equation}
The denominator \(2s_{k,0}+\Gamma_{kk}\) is the derivative of the decoupled
scalar polynomial. This is the scalar version of the general QEP sensitivity
formula. The modal correction can move both real and imaginary parts, so its
effect on damping ratio must be computed rather than inferred from graph
frequency alone.

\section{Non-Proportionality as a Gate}

The previous equations suggest a theoretical gate before using graph metrics as
damping screens. A fixed-\(D\) graph screen is plausible when
\begin{equation}
  \norm{E}\ll\norm{\Gamma_0}
  \label{eq:small_nonprop_gate_1}
\end{equation}
and when the resolvent denominators in
\eqref{eq:general_modal_correction} are not near resonance over the frequency
band of interest. This does not prove that the screen is accurate; it states
the conditions under which the modal coupling neglected by the screen is small.

Structural-dynamics non-proportionality indices formalize similar ideas:
measure how far the damping operator is from the modal subspace in which the
undamped system is diagonal \cite{conde2021nonproportionality}. The power-grid
interpretation is more specific. The off-diagonal damping matrix identifies
which graph modes exchange damping through heterogeneous device locations and
damping-to-inertia ratios.

\section{Connection to the Next Level}

Fixed non-commuting damping is not the same as converter self-energy. In
\eqref{eq:general_modal_correction}, the eliminated coordinates are other
graph modes of the same fixed-QEP system. In controller self-energy, the
eliminated coordinates are physical control and auxiliary states. The algebraic
similarity is nevertheless important. It shows that rational dependence is not
an exotic artifact; it is the natural result of eliminating coupled dynamic
coordinates. The next chapter applies that logic to controller-state
condensation.


\part{Rational Self-Energy}
\chapter{Rational Self-Energy and Schur Nonlinear Eigenproblems}

The third level of the hierarchy begins when converter and auxiliary controller
states are condensed. The reduced problem is no longer a fixed-coefficient QEP.
It is a nonlinear eigenvalue problem whose retained correction is rational in
the spectral parameter. This chapter derives that object and explains why
\(\Pi(s)\), not a static damping matrix, is the general theoretical replacement
after controller or algebraic states are eliminated.

\section{Full and Retained State Partitions}

Let the linearized state matrix be partitioned as
\begin{equation}
  A=
  \begin{bmatrix}
    A_{ee} & A_{ec}\\
    A_{ce} & A_{cc}
  \end{bmatrix}.
  \label{eq:self_full_partition}
\end{equation}
The retained states \(x_e\) may include angles, speeds, or selected interface
variables. The condensed states \(x_c\) may include exciters, governors,
stabilizers, PLLs, voltage and current controllers, plant controllers, filters,
and other auxiliary dynamics. The eigenvalue equation is
\begin{align}
  (sI-A_{ee})x_e-A_{ec}x_c &=0,\\
  -A_{ce}x_e+(sI-A_{cc})x_c &=0.
\end{align}
When \(sI-A_{cc}\) is invertible,
\begin{equation}
  x_c=(sI-A_{cc})^{-1}A_{ce}x_e.
\end{equation}
Substitution gives
\begin{equation}
  S(s)x_e=0,\qquad
  S(s)=sI-A_{ee}-A_{ec}(sI-A_{cc})^{-1}A_{ce}.
  \label{eq:self_schur_operator}
\end{equation}

\begin{proposition}[Exact Schur equivalence]
If \(sI-A_{cc}\) is invertible, then \(s\) is an eigenvalue of the full matrix
\(A\) with retained component \(x_e\ne0\) if and only if \(S(s)x_e=0\).
\end{proposition}

\begin{proof}
The proof follows by the two substitutions above. Starting from a full
eigenpair gives \(S(s)x_e=0\). Conversely, if \(S(s)x_e=0\), define
\(x_c=(sI-A_{cc})^{-1}A_{ce}x_e\). Then the block eigenvalue equations are
satisfied. \qed
\end{proof}

The result is exact, but the retained operator is nonlinear in \(s\). Any
constant reduced matrix that tries to replace \(A_{ec}(sI-A_{cc})^{-1}A_{ce}\)
must be an approximation.

\section{Descriptor Schur Complement}

Power-system linearizations are often descriptor systems rather than explicit
ODEs. The more general eigenvalue equation is
\begin{equation}
  (sE-A)x=0.
  \label{eq:self_descriptor_pencil}
\end{equation}
Partition retained and condensed variables as
\begin{equation}
  sE-A
  =
  \begin{bmatrix}
    sE_{rr}-A_{rr} & sE_{rc}-A_{rc}\\
    sE_{cr}-A_{cr} & sE_{cc}-A_{cc}
  \end{bmatrix}.
  \label{eq:self_descriptor_partition}
\end{equation}
When \(sE_{cc}-A_{cc}\) is invertible, the exact retained pencil is
\begin{equation}
  S(s)
  =
  sE_{rr}-A_{rr}
  -
  (sE_{rc}-A_{rc})
  (sE_{cc}-A_{cc})^{-1}
  (sE_{cr}-A_{cr}).
  \label{eq:self_descriptor_schur}
\end{equation}
This form is the correct starting point for DAEs, semi-explicit models,
algebraic variables that are not eliminated first, and identified port models.
The explicit-state formula \eqref{eq:self_schur_operator} is the special case
\(E=I\) with compatible zero off-diagonal descriptor blocks. The descriptor
form also makes visible which eliminated poles are genuine dynamic poles and
which are algebraic or infinite-eigenvalue artifacts of the chosen
representation.

\section{Port Equivalence of Self-Energy}

The self-energy must be treated as a port map, not as a property of one hidden
state realization. Let two condensed subsystems have hidden coordinates
\(z_1\) and \(z_2\), possibly of different dimensions, and suppose they produce
the same retained correction on a spectral domain \(\mathcal D\):
\begin{equation}
  \Pi_1(s)=\Pi_2(s),\qquad s\in\mathcal D .
  \label{eq:self_port_equiv_pi}
\end{equation}
Then every retained-port statement that uses only
\begin{equation}
  T(s)=s^2M+sD_0+L+\Pi(s)
  \label{eq:self_port_equiv_T}
\end{equation}
is invariant under the change of hidden realization.

\begin{theorem}[Spectral sufficiency of the port self-energy]
\label{thm:self_port_sufficiency}
Assume that the retained matrices \(M,D_0,L\) are the same and that two hidden
subsystems satisfy \eqref{eq:self_port_equiv_pi} on a domain \(\mathcal D\)
that contains the protected roots and the contours used for certification.
Then the two reduced operators have the same retained roots in
\(\mathcal D\), the same retained Green function \(T(s)^{-1}\) wherever the
inverse exists, and the same low-rank action determinants for any action that
is expressed only through the retained ports.
\end{theorem}

\begin{proof}
Under \eqref{eq:self_port_equiv_pi}, the two retained operators
\eqref{eq:self_port_equiv_T} are identical on \(\mathcal D\). Equality of
roots, inverses on the resolvent set, and any determinant formula built from
that inverse follows directly. The full hidden-state eigenvectors may differ,
but those coordinates are not observables of a retained-port decision.
\end{proof}

This theorem is deliberately modest. It does not say that hidden
participation factors, Euclidean controller-state norms, or internal energy
coordinates are invariant. It says that a retained-port planning claim should
depend on the map \(\Pi(s)\), the port normalization, and the protected
spectral region. Similarity transformations \(z=S\widehat z\), changes of
hidden units, or a different order of Schur elimination are harmless only to
the extent that they preserve the same port self-energy on the domain being
used.

\section{General Second-Order Self-Energy Form}

The exact Schur operator \eqref{eq:self_schur_operator} does not by itself imply
that the reduced second-order model contains only a dynamic damping term. To see
the required structure, write a retained mechanical model with eliminated
controller states \(z\):
\begin{align}
  M\ddot\theta+D_0\dot\theta+\Lt\theta+Gz &=0,\\
  \dot z &= A_{cc}z+B_\theta\theta+B_\omega\dot\theta .
  \label{eq:self_structured_controller}
\end{align}
After Laplace transformation and elimination of \(z\),
\begin{equation}
  T(s)
  =
  s^2M+sD_0+\Lt+\Pi(s),
  \qquad
  \Pi(s)=G(sI-A_{cc})^{-1}(B_\theta+sB_\omega).
  \label{eq:self_general_pi}
\end{equation}
The sign of \(G,B_\theta,B_\omega\) depends on the port convention; the
important point is the channel separation. In general, \(\Pi(s)\) contains a
rational stiffness contribution and a rational damping contribution:
\begin{equation}
  \Pi(s)=\Pi_K(s)+s\Pi_D(s).
  \label{eq:self_pi_decomposition}
\end{equation}
Only under the additional assumption \(\Pi_K(s)=0\), or after a defensible
absorption of \(\Pi_K\) into a stated stiffness correction, does the operator
reduce to the dynamic-damping form
\begin{equation}
  T(s)=s^2M+s\Sigma(s)+\Lt,
  \label{eq:self_second_order}
\end{equation}
where
\begin{equation}
  \Sigma(s)=D_0+\Pi_D(s)
  =
  D_0+C(sI-A_{cc})^{-1}B
  \quad\hbox{in the damping-channel special case.}
  \label{eq:self_sigma_def}
\end{equation}
This is the precise transition from a fixed \(D\) to a rational
self-energy \(\Pi(s)\). The broader self-energy is \(\Pi(s)\);
\(\Sigma(s)\) names the dynamic damping component.

The notation \(\Sigma(s)\) is therefore deliberately restricted. It is not
simply damping in the classical viscous sense. It can include dissipative,
anti-dissipative, phase-lag, resonant, and cross-channel effects, depending on
the poles and residues of the condensed block. Its Hermitian part on
\(s=j\omega\) may suggest dissipative behavior over a frequency band, but the
full eigenvalue effect also depends on eigenvectors and on \(T_s(s)\).

\section{Dynamic Green Function and Effective Resistance}

The inverse of the retained operator is the dynamic Green function of the
condensed network:
\begin{equation}
  G_T(s)=T(s)^{-1},
  \qquad s\in\rho(T),
  \label{eq:self_green_function}
\end{equation}
where \(\rho(T)\) denotes the resolvent set. This object is the map from a
retained port injection to the retained response at spectral parameter \(s\).
It is therefore the natural carrier of dynamic grid strength: it contains the
graph stiffness, inertia, fixed damping, controller self-energy, and
non-normal amplification in one operator.

\begin{definition}[Dynamic effective resistance]
For a line or candidate line \(e=(i,j)\), let \(a_e=e_i-e_j\). The dynamic
effective resistance of the line is
\begin{equation}
  \mathcal R_e(s)=a_e^TG_T(s)a_e
  =
  a_e^T
  \left[
    s^2M+sD_0+L+\Pi(s)
  \right]^{-1}
  a_e .
  \label{eq:self_dynamic_resistance}
\end{equation}
\end{definition}

When \(s=0\), \(\Pi(0)=0\), and the uniform angle reference is removed, this
quantity reduces to the classical graph effective resistance
\begin{equation}
  \mathcal R_e(0)=a_e^TL^\dagger a_e .
  \label{eq:self_static_resistance_limit}
\end{equation}
For \(s\ne0\), \(\mathcal R_e(s)\) is not a physical resistor. It is a
meromorphic port compliance. Its poles are the retained roots of \(T(s)\), and
its residues measure how strongly each mode is seen across the endpoint
difference of line \(e\).

\begin{theorem}[Exact finite line-action law]
\label{thm:self_finite_line_action}
Let a finite line-weight change enter the retained operator as
\begin{equation}
  T_{\rm new}(s)=T(s)+\Delta b_ea_ea_e^T .
  \label{eq:self_line_action_update}
\end{equation}
For all \(s\) where \(T(s)\) is invertible,
\begin{equation}
  \det T_{\rm new}(s)
  =
  \det T(s)\left[1+\Delta b_e\mathcal R_e(s)\right].
  \label{eq:self_line_action_determinant}
\end{equation}
Thus any new or shifted root created by this rank-one line action and not
already a root of \(T\) satisfies
\begin{equation}
  1+\Delta b_e\mathcal R_e(s)=0 .
  \label{eq:self_line_action_scalar}
\end{equation}
\end{theorem}

\begin{proof}
Apply the matrix determinant lemma to
\(T(s)+\Delta b_ea_ea_e^T\). This gives
\[
  \det(T+\Delta b_ea_ea_e^T)
  =
  \det T\left(1+\Delta b_ea_e^TT^{-1}a_e\right),
\]
which is \eqref{eq:self_line_action_determinant}.
\end{proof}

This result is stronger than a local derivative. It gives an exact scalar
equation for a finite rank-one line modification even when \(T(s)\) is a large
rational NEP. The same equation also explains why classical effective
resistance is a static shadow of a richer dynamic object: the graph quantity
is the zero-frequency limit of the port Green function.

\section{Low-Rank Dynamic Action Calculus}

Line reinforcement is only one member of a broader action family. Suppose a
planning or retuning action enters through a low-rank retained-port update
\begin{equation}
  T_u(s)=T(s)+U(s)K_u(s)V(s)^* .
  \label{eq:self_low_rank_action}
\end{equation}
Define the action interaction matrix
\begin{equation}
  \mathcal G_u(s)=V(s)^*T(s)^{-1}U(s).
  \label{eq:self_action_green_matrix}
\end{equation}
Then the determinant lemma gives
\begin{equation}
  \det T_u(s)
  =
  \det T(s)\det\!\left[I+K_u(s)\mathcal G_u(s)\right].
  \label{eq:self_low_rank_determinant}
\end{equation}
The large NEP has been reduced, for the purpose of this action, to a
determinant whose dimension is the number of action ports.

Several planning levers become special cases of
\eqref{eq:self_low_rank_determinant}:
\begin{align}
  \hbox{line } e:\quad
  &U=V=a_e,\qquad K_u(s)=\Delta b_e, \label{eq:self_action_line}\\
  \hbox{local damping at bus } i:\quad
  &U=V=e_i,\qquad K_u(s)=s\Delta d_i, \label{eq:self_action_damping}\\
  \hbox{local inertia at bus } i:\quad
  &U=V=e_i,\qquad K_u(s)=s^2\Delta m_i, \label{eq:self_action_inertia}\\
  \hbox{controller retuning through ports}:\quad
  &U=U_c,\quad V=V_c,\quad K_u(s)=\Delta K_c(s).
  \label{eq:self_action_controller}
\end{align}
Thus line, damping, inertia, and controller actions are not four unrelated
sensitivity formulas. They are four kernels projected through the same Green
function \(T(s)^{-1}\).

\section{Modal Sensitivity as a Local Limit}

The usual NEP sensitivity is recovered by expanding the Green function near a
simple root. Let \(T(s_k)x_k=0\), \(y_k^*T(s_k)=0\), and
\(y_k^*T_s(s_k)x_k\ne0\). Then
\begin{equation}
  T(s)^{-1}
  \sim
  \frac{x_ky_k^*}
  {(s-s_k)y_k^*T_s(s_k)x_k}.
  \label{eq:self_green_pole_expansion}
\end{equation}
For a line \(e\),
\begin{equation}
  \mathcal R_e(s)
  \sim
  \frac{(a_e^Tx_k)(y_k^*a_e)}
  {(s-s_k)y_k^*T_s(s_k)x_k}.
  \label{eq:self_resistance_residue}
\end{equation}
Substituting this expansion into the finite-action equation
\eqref{eq:self_line_action_scalar} gives
\begin{equation}
  s-s_k
  \approx
  -\Delta b_e
  \frac{y_k^*a_ea_e^Tx_k}
  {y_k^*T_s(s_k)x_k},
  \label{eq:self_line_local_shift}
\end{equation}
which is exactly the first-order line sensitivity
\begin{equation}
  \frac{ds_k}{db_e}
  =
  -
  \frac{y_k^*a_ea_e^Tx_k}
  {y_k^*T_s(s_k)x_k}.
  \label{eq:self_line_sensitivity_limit}
\end{equation}
The sensitivity is therefore not an isolated formula. It is the local residue
limit of an exact finite-action law.

\section{Pole-Residue Structure}

Assume for exposition that \(A_{cc}\) is diagonalizable:
\begin{equation}
  A_{cc}=R\Lambda_c R^{-1},\qquad
  \Lambda_c=\diag(\alpha_1,\ldots,\alpha_m).
\end{equation}
Then
\begin{equation}
  C(sI-A_{cc})^{-1}B
  =
  \sum_{\ell=1}^m
  \frac{C r_\ell \ell_\ell^T B}{s-\alpha_\ell},
  \label{eq:self_partial_fraction}
\end{equation}
where \(r_\ell\) and \(\ell_\ell\) are right and left eigenvectors of
\(A_{cc}\), normalized so that \(\ell_\ell^Tr_\ell=1\). The residue
\(C r_\ell \ell_\ell^T B\) measures how strongly the eliminated controller pole
\(\alpha_\ell\) is seen by retained coordinates.

This expression explains why a nearby controller pole can dominate a damping
margin. If \(s\) is close to \(\alpha_\ell\), the denominator is small. If the
residue projects strongly onto the retained eigenvectors, the self-energy term
can dominate the eigenvalue sensitivity.

\section{NEP Sensitivity with Self-Energy}

For \(T(s,\rho)v=0\), the simple-root sensitivity is
\begin{equation}
  \frac{ds_\star}{d\rho}
  =
  -\frac{u^H T_\rho(s_\star,\rho)v}
  {u^H T_s(s_\star,\rho)v}.
  \label{eq:self_nep_sens}
\end{equation}
For the general operator \eqref{eq:self_general_pi},
\begin{equation}
  T_s(s)=2sM+D_0+\Pi_s(s).
  \label{eq:self_Ts_general}
\end{equation}
In the damping-channel specialization \eqref{eq:self_second_order},
\begin{equation}
  T_s(s)=2sM+\Sigma(s)+s\Sigma'(s),
  \label{eq:self_Ts}
\end{equation}
with
\begin{equation}
  \Sigma'(s)=-C(sI-A_{cc})^{-2}B.
  \label{eq:self_sigma_prime}
\end{equation}
Thus even the denominator of the sensitivity contains controller-resolvent
terms. A fixed scalar damping bridge cannot reproduce this dependence unless
the rational term is negligible over the critical spectral region.

The additional denominator term \(s\Sigma'(s)\) is the precise mathematical
place where the condensed controller block retains spectral memory. It is not
another static damping coefficient; it is a slope of the self-energy with
respect to the eigenvalue being solved for. Using
\eqref{eq:self_sigma_prime}, this slope contains the squared resolvent
\((sI-A_{cc})^{-2}\). Hence it becomes large when the retained pole is close to
a visible condensed-controller pole. In that regime, replacing
\(\Sigma(s)\) by a constant matrix may preserve one fitted value of the
self-energy while losing the local pole-motion geometry.

\section{Why Static Damping Bridges Fail}

A static bridge assumes that the reduced damping mechanism can be represented
by a constant matrix \(D_{\rm red}\). In the self-energy theory, that assumption
requires
\begin{equation}
  C(sI-A_{cc})^{-1}B\approx D_{\rm red}-D_0
  \label{eq:self_static_approx}
\end{equation}
over all poles and frequency bands relevant to the claim. This is a strong
condition. It can hold locally if the controller poles are far away and the
resolvent varies slowly. It fails when the critical pole is near a controller
pole, when residues are large, or when the relevant mode is dominated by
condensed states.

The failure is structural rather than cosmetic. A constant matrix has no poles.
It cannot represent the denominator \(s-\alpha_\ell\). It can approximate the
rational term at a selected point, but it cannot preserve the full spectral
mechanism over a region unless additional assumptions are imposed.

\begin{remark}[Static models are local models]
The conclusion is not that every static reduced model is useless. The
conclusion is that a static bridge has a validity region. It is defensible
when controller poles are far from the protected region, when residues are
small in the retained eigenvector directions, and when \(\Sigma'(s)\) is small
enough that the denominator of \eqref{eq:self_nep_sens} is not materially
changed. It is not defensible as a universal damping-margin surrogate near a
controller resonance.
\end{remark}

\section{Certified Adequacy of Reduced Operators}

The hierarchy is useful only if it contains an escalation rule. Let
\(\widehat T(s)\) be a reduced or approximate operator, such as a static bridge,
a fitted self-energy, or a lower-order rational approximation of \(\Pi(s)\).
On a closed contour \(\Gamma\) that bounds a protected region, define
\begin{equation}
  \eta_\Gamma
  =
  \sup_{s\in\Gamma}
  \left\|
  \widehat T(s)^{-1}\left[T(s)-\widehat T(s)\right]
  \right\|_2 .
  \label{eq:self_eta_gamma}
\end{equation}

\begin{theorem}[Contour adequacy certificate]
\label{thm:self_contour_adequacy}
Assume that \(T\) and \(\widehat T\) are regular analytic matrix functions on
and inside \(\Gamma\), that \(\widehat T(s)\) is nonsingular on \(\Gamma\), and
that
\begin{equation}
  \eta_\Gamma<1 .
  \label{eq:self_eta_condition}
\end{equation}
Then \(T\) and \(\widehat T\) have the same number of eigenvalues inside
\(\Gamma\), counted with algebraic multiplicity.
\end{theorem}

\begin{proof}
The statement is the matrix-function analogue of Rouche's theorem. Since
\(\|\widehat T^{-1}(T-\widehat T)\|_2<1\) on \(\Gamma\), the homotopy
\(\widehat T+\tau(T-\widehat T)\) remains nonsingular on \(\Gamma\) for
\(0\le\tau\le1\). The argument principle for regular analytic matrix
functions then preserves the eigenvalue count inside the contour
\cite{kato1995perturbation,guettel2017nep}.
\end{proof}

This certificate turns the operator hierarchy into a computational policy:
\begin{equation}
  \eta_\Gamma<1
  \quad\Longrightarrow\quad
  \hbox{the reduced operator preserves the protected eigenvalue count,}
  \label{eq:self_eta_policy_accept}
\end{equation}
whereas
\begin{equation}
  \eta_\Gamma\ge1
  \quad\Longrightarrow\quad
  \hbox{escalate to a richer NEP, identified self-energy, or full DAE.}
  \label{eq:self_eta_policy_escalate}
\end{equation}

A local decision also needs a sign certificate. If an approximate model
predicts a margin derivative \(\widehat g\), and an error analysis gives
\begin{equation}
  |g-\widehat g|\le\varepsilon_g ,
  \label{eq:self_sign_error_bound}
\end{equation}
then
\begin{equation}
  |\widehat g|>\varepsilon_g
  \quad\Longrightarrow\quad
  \operatorname{sign}(g)=\operatorname{sign}(\widehat g).
  \label{eq:self_sign_certificate}
\end{equation}
Without a bound of this type, a ranked action is a screen rather than a
certified recommendation.

\section{Data-Identified Self-Energy}

The self-energy theory should not require access to proprietary hidden
matrices. If a manufacturer, EMT model, or field experiment provides only
terminal data, the retained-port object can be estimated as
\begin{equation}
  \widehat\Pi(s)
  \approx
  \Pi(s)
  \label{eq:self_identified_pi}
\end{equation}
from frequency scans, injected perturbations, PMU records, or black-box
simulation. The estimation method may be vector fitting, Loewner realization,
gray-box calibration, or a structure-preserving DAE identification method.
Recent converter-identification work is moving in this direction by combining
terminal time-series data with physics-informed converter structure and
frequency-domain validation \cite{darii2026graybox}; recent
port-Hamiltonian DAE identification similarly emphasizes preserving
interconnection and passivity structure under noisy input-output data
\cite{hagelaars2026identification}.

An identified self-energy is useful for planning only if it comes with an
error model. On a protected frequency or contour set \(\Omega\), require
\begin{equation}
  \|\Pi(s)-\widehat\Pi(s)\|_2\le\varepsilon_\Pi(s),
  \qquad s\in\Omega .
  \label{eq:self_pi_identification_error}
\end{equation}
This bound propagates through the same adequacy certificate:
\begin{equation}
  \widehat T(s)=s^2M+sD_0+L+\widehat\Pi(s),
  \qquad
  \eta_\Gamma
  =
  \sup_{s\in\Gamma}
  \left\|
  \widehat T(s)^{-1}\left[\Pi(s)-\widehat\Pi(s)\right]
  \right\|_2 .
  \label{eq:self_identified_eta}
\end{equation}
The practical chain is therefore
\begin{equation}
  \hbox{data}
  \rightarrow
  \widehat\Pi(s)
  \rightarrow
  \varepsilon_\Pi
  \rightarrow
  \eta_\Gamma
  \rightarrow
  \hbox{certified or rejected action}.
  \label{eq:self_data_to_certification}
\end{equation}
The thesis contribution is not merely to identify a converter. It is to decide
how much identification error can be tolerated before a planning recommendation
changes sign or loses its protected eigenvalue count.

\section{Mode Families}

The Schur partition motivates a mode-family diagnostic, but a raw Euclidean
ratio of retained and condensed state norms is not invariant under a change of
units or controller realization. If physically meaningful weights are available,
use the weighted index
\begin{equation}
  \pi_c^{(W)}
  =
  \frac{x_c^H W_c x_c}
  {x_e^H W_e x_e+x_c^H W_c x_c},
  \label{eq:self_control_participation}
\end{equation}
where \(W_e,W_c\succeq0\) represent compatible energy, power, or per-unit
normalization choices. Small \(\pi_c^{(W)}\) indicates a retained-state or
network-family mode. Large \(\pi_c^{(W)}\) indicates a control-family mode.
Intermediate values indicate mixed modes.

When no defensible \(W_c\) exists, the more invariant diagnostic is not internal
state size but port visibility:
\begin{equation}
  \mathcal V_c(s)
  =
  \norm{C(sI-A_{cc})^{-1}B}_2 .
  \label{eq:self_controller_visibility}
\end{equation}
This quantity asks whether the condensed block is visible from retained
coordinates at the frequency being tested. It changes under input-output
scaling, so the ports must still be normalized, but it is not changed by a
similarity transformation of the hidden controller state.

\section{Structured Pseudospectrum and Fragility Radius}

Nominal eigenvalues do not fully describe a non-normal converter-rich model.
The unstructured NEP pseudospectrum is
\begin{equation}
  \Lambda_\varepsilon(T)
  =
  \left\{
  s:\sigma_{\min}(T(s))\le\varepsilon
  \right\}.
  \label{eq:self_unstructured_pseudospectrum}
\end{equation}
This set is useful as a numerical warning, but it permits arbitrary matrix
perturbations. Planning perturbations are not arbitrary: they come from line
uncertainty, damping placement, inertia placement, controller gains, delays,
operating-point errors, or identified-model errors. The relevant object is
therefore the action-structured pseudospectrum
\begin{equation}
  \Lambda_\varepsilon^{\mathcal A}(T)
  =
  \left\{
  s:\exists u\in\mathcal A,\ 
  \|W_uu\|\le\varepsilon,\ 
  \det T(s,u)=0
  \right\},
  \label{eq:self_structured_pseudospectrum}
\end{equation}
where \(\mathcal A\) is the declared physical action or uncertainty family and
\(W_u\) normalizes units and costs. Recent nonlinear spectral algorithms make
the broader computation of spectra and pseudospectra increasingly relevant for
NEPs, but the power-system decision must still restrict perturbations to
physically admissible directions \cite{colbrook2025universal}.

For a protected spectral boundary \(\partial\mathcal S\), define the dynamic
fragility radius
\begin{equation}
  r_{\rm frag}
  =
  \inf
  \left\{
  \|W_uu\|:
  \exists s\in\partial\mathcal S,\ 
  \det T(s,u)=0
  \right\}.
  \label{eq:self_fragility_radius}
\end{equation}
This number answers the planning question directly: what is the smallest
physical perturbation that can violate the protected margin? A damping-ratio
margin, an abscissa margin, a distance-to-resonance margin, a passivity margin,
or a protected-gain margin can each be used to define
\(\partial\mathcal S\). Local pole-motion cones and hidden-margin derivatives
are then first-order approximations to \eqref{eq:self_fragility_radius}, not
substitutes for it.

\section{Passivity and Resolvent Gates}

A stable \(A_{cc}\) block does not guarantee that the condensed self-energy is
benign. A passivity statement requires a declared input-output pairing and a
supply rate in the sense of dissipativity theory \cite{willems1972dissipative}.
For a MIMO retained port, reducing passivity to one scalar index can hide the
direction that causes the deficiency. With input and output normalizers
\(W_u\) and \(W_y\), define the directional passivity matrix
\begin{equation}
  \mathcal P_\Sigma(\omega)
  =
  \frac12
  \left[
  W_y\Sigma(j\omega)W_u^{-1}
  +
  W_u^{-*}\Sigma(j\omega)^*W_y^*
  \right].
  \label{eq:self_directional_passivity}
\end{equation}
This matrix records both intensity and direction of passivity excess or
deficit. Recent matrix-valued passivity work is useful here precisely because
scalar indices can lose MIMO channel structure \cite{ru2026matrix}. The
unweighted Hermitian screen used below is the special case
\begin{equation}
  H_\Sigma(\omega)=
  \frac{\Sigma(j\omega)+\Sigma(j\omega)^H}{2}.
  \label{eq:self_hermitian}
\end{equation}
Negative eigenvalues of \(H_\Sigma(\omega)\), or of the weighted matrix
\(\mathcal P_\Sigma(\omega)\), indicate frequencies where the condensed
contribution acts like negative damping in a retained port direction. Large
singular values of the resolvent identify frequencies where the eliminated
block strongly amplifies retained-state motion. Dynamic passivity multipliers
provide a broader certificate family when strict passivity in the trivial
metric is too conservative \cite{gorbunov2026dynamic}.

Recent passivity-based grid-forming work suggests a useful interpretation of
this gate \cite{he2024passivity,gui2024passivity}. If a controller parameter
\(\kappa\) changes the condensed block, write
\(A_{cc}=A_{cc}(\kappa)\) and define the band-limited self-energy passivity
screen
\begin{equation}
  \mu_\Sigma(\Omega;\kappa)
  =
  \inf_{\omega\in\Omega}
  \lambda_{\min}\!\left(H_\Sigma(\omega;\kappa)\right),
  \label{eq:self_passivity_index}
\end{equation}
over the frequency band \(\Omega\) that contains the protected electromechanical
or controller-interaction modes. If \(\mu_\Sigma(\Omega;\kappa)>0\), the condensed
controller contribution is dissipative in the retained coordinates over that
band. If \(\mu_\Sigma(\Omega;\kappa)<0\), there exists a retained direction and a
frequency where the condensed block supplies negative damping.

\begin{theorem}[Local self-energy passivity-retuning certificate]
\label{thm:self_passivity_retuning}
Let \(\Omega\subset\R\) be compact. Assume \(A_{cc}(\kappa)\), \(B(\kappa)\),
and \(C(\kappa)\) are continuously differentiable on an interval
\(\mathcal K\), and that \(j\omega I-A_{cc}(\kappa)\) is invertible for all
\((\omega,\kappa)\in\Omega\times\mathcal K\). At a fixed \(\kappa\), suppose
the infimum in \eqref{eq:self_passivity_index} is attained at a unique
\(\omega_\star\), and that the smallest eigenvalue of
\(H_\Sigma(\omega_\star;\kappa)\) is simple with unit eigenvector \(q_\star\).
Then \(\mu_\Sigma\) is differentiable at \(\kappa\), and
\begin{equation}
  \frac{d\mu_\Sigma}{d\kappa}
  =
  q_\star^H
  \frac{\partial H_\Sigma(\omega_\star;\kappa)}{\partial\kappa}
  q_\star .
  \label{eq:self_passivity_envelope}
\end{equation}
If this quantity is nonnegative throughout an interval, then
\(\mu_\Sigma\) is nondecreasing on that interval. If it is positive before a
single zero and negative afterward, the band-limited passivity margin is locally
unimodal and the zero gives a candidate retuning optimum.
\end{theorem}

\begin{proof}
For fixed \(\omega\), \(H_\Sigma(\omega;\kappa)\) is Hermitian and differentiable
under the stated invertibility assumptions. The standard simple-eigenvalue
perturbation formula gives
\(\partial_\kappa\lambda_{\min}=q_\star^H(\partial_\kappa H_\Sigma)q_\star\)
at the active frequency. Since the margin is an infimum over the compact
frequency set, Danskin's envelope theorem applies to the active minimizer
\cite{danskin1967maxmin,rockafellar1998variational}. Under uniqueness and
simplicity, the directional formula collapses to
\eqref{eq:self_passivity_envelope}. Monotonicity follows by integrating the
derivative over the interval. The unimodal statement is the same first-order
condition with one sign change.
\end{proof}

The derivative in \eqref{eq:self_passivity_envelope} is computable from the same
resolvent that defines the self-energy. With \(R(s,\kappa)=(sI-A_{cc}(\kappa))^{-1}\),
\begin{equation}
  \frac{\partial\Sigma}{\partial\kappa}
  =
  C_\kappa R B
  +
  C R_\kappa B
  +
  C R B_\kappa,\qquad
  R_\kappa
  =
  R\,\frac{\partial A_{cc}}{\partial\kappa}\,R.
  \label{eq:self_sigma_kappa_derivative}
\end{equation}
If only the controller matrix \(A_{cc}\) depends on \(\kappa\), the first and
third terms vanish. The sign of \(d\mu_\Sigma/d\kappa\) is therefore not a
generic algebraic fact; it is a physical property of the controller family,
operating point, and frequency band.

For a grid-following PLL, this distinction matters. Increasing PLL bandwidth can
improve phase tracking over one range and degrade weak-grid damping over another.
The thesis therefore treats monotonicity as a local certificate, not as a global
law. The actionable design problem is
\begin{equation}
  \kappa_\star\in
  \operatorname*{arg\,max}_{\kappa\in\mathcal K}
  \mu_\Sigma(\Omega;\kappa),
  \qquad
  \text{subject to a protected-pole gate.}
  \label{eq:self_passivity_optimal_kappa}
\end{equation}
This formulation also explains what must be validated after retuning:
\(\mu_\Sigma\) should increase in the relevant band, and the protected pole
should move in the desired direction.

This is not a complete stability theorem for the retained NEP. It is a
certificate candidate. The missing companion step is the pole sensitivity:
\begin{equation}
  \frac{d}{d\rho}\mu_\Sigma(\Omega;\rho),
  \qquad
  \frac{ds_\star}{d\rho}
  =
  -\frac{u^H T_\rho(s_\star,\rho)v}
  {u^H T_s(s_\star,\rho)v}.
  \label{eq:self_passivity_retuning}
\end{equation}
A controller-retuning action is credible only if it increases the passivity
screen in the relevant band and moves the protected pole in the correct
direction. This gives a principled version of ``PLL retuning'': the action is
not merely moving an eigenvalue numerically; it is attempting to restore a
band-limited dissipativity margin of the Schur self-energy.

These diagnostics are gates, not standalone stability theorems. They must be
connected to eigenvalue sensitivities and full-model validation before they can
support planning decisions. Their theoretical value is that they point to the
operator that must be inspected: \(\Sigma(s)\), not a scalar damping constant.

\section{Sign Characteristic and Controller-Resolvent Collisions}

The language of Krein signature is mature for Hamiltonian and self-adjoint
spectral problems, and the matrix-polynomial literature uses the more precise
term \emph{sign characteristic} \cite{gohberg1982matrixpolynomials,
gohberg1983meromorphic,mehrmann2016sign}. Kollár and Miller give a useful
graphical interpretation through Evans--Krein functions for polynomial operator
pencils \cite{kollar2014graphical}. That literature should not be cited as if it
already proved the converter result needed here. The thesis uses it as the
correct mathematical foundation for a narrower question: can a rational,
controller-condensed self-energy produce an operational sign collision that
predicts weak-grid synchronization loss?

For a Hermitian matrix polynomial, the sign characteristic is invariant under
structure-preserving transformations and is tied to the perturbation behavior of
real or imaginary-axis eigenvalues. A converter-condensed operator
\[
  T(s)=s^2M+sD_0+\Lt+\Pi(s)
\]
is usually not Hermitian as a function of \(s\), and its rational term may be
dissipative and non-normal. Therefore the following object is only an
operational signature gate, not a completed Krein theorem:
\begin{equation}
  \chi_k^{\rm op}
  =
  \operatorname{sign}
  \Real\!\left(v_k^H T_s(s_k)v_k\right),
  \qquad
  T(s_k)v_k=0.
  \label{eq:self_operational_signature}
\end{equation}
If left and right eigenvectors differ strongly, the bilinear variant
\(u_k^H T_s(s_k)v_k\) should be tracked together with eigenvalue conditioning.
If \(s_k=j\omega_k\) lies on a symmetry line where a Hermitian reduction is
available, one may instead inspect the inertia of
\((T_s(j\omega_k)+T_s(j\omega_k)^H)/2\) on the active eigenspace.

The reason this is relevant to controller-resonant instability is the derivative
term
\[
  T_s(s)=2sM+D_0+\Pi_s(s)
  \quad
  \hbox{or, in the damping-channel case,}\quad
  T_s(s)=2sM+\Sigma(s)+s\Sigma'(s).
\]
Near a pole of \(A_{cc}\), the resolvent derivative contains squared-resolvent
growth. A mode can then become dominated by the controller self-energy rather
than by the mechanical part of the QEP. The proposed mechanism is a
controller-resolvent sign collision: two tracked modes approach while their
operational signatures differ, and the denominator
\(u^H T_s(s)v\) in the NEP sensitivity becomes small or ill-conditioned.

This bridge is a research claim, not a settled theorem. It is credible only if a
benchmark shows the same event in three views: a sign change or singularity in
\eqref{eq:self_operational_signature}, growth of the absolute self-energy
dominance index, and a full-model eigenvalue or time-domain loss of damping. The
general theory of dissipation-induced instabilities explains why such collisions
are plausible \cite{krechetnikov2007dissipation}; the new power-system claim
would be that the Schur self-energy identifies the controller channel that
causes them.


\part{Nonlinear Models and Memory}
\chapter{Nonlinear Models, Memory, and Stability Beyond the NEP}

The rational nonlinear eigenvalue problem developed in the previous part is
nonlinear in the spectral parameter \(s\). It is still a small-signal object:
it comes from a linearization in time and from exact or approximate
elimination of linearized states. This chapter separates that NEP theory from
the genuinely nonlinear power-system problem. The latter concerns equilibria,
local stability, bifurcations, regions of attraction, amplitude-dependent
damping, and memory generated by eliminating nonlinear controllers.

\thesisbox{thesisorange}{Separation of roles}{
The Schur NEP answers: what rational operator preserves the linearized
controller memory around an operating point? The nonlinear chapter answers:
which equilibria exist, how they lose stability, and whether a spectral margin
survives finite perturbations.}

\section{Nonlinear DAE Layer}

Let the full converter-rich model be written as a semi-explicit DAE,
\begin{align}
  \dot x &= f(x,y,p), \label{eq:nonlin_dae_f}\\
  0 &= g(x,y,p), \label{eq:nonlin_dae_g}
\end{align}
where \(x\) collects dynamic states, \(y\) collects algebraic network
variables, and \(p\) is an operating or planning parameter. A regular
equilibrium \((x^\star,y^\star)\) satisfies
\begin{equation}
  f(x^\star,y^\star,p)=0,\qquad
  g(x^\star,y^\star,p)=0,
  \qquad
  g_y(x^\star,y^\star,p)\ \hbox{nonsingular}.
  \label{eq:nonlin_equilibrium_regular}
\end{equation}
When the algebraic Jacobian is nonsingular, the implicit-function theorem gives
a local algebraic manifold \(y=h(x,p)\). The reduced ODE is
\begin{equation}
  \dot x=F(x,p)=f(x,h(x,p),p),
  \label{eq:nonlin_reduced_ode}
\end{equation}
and the small-signal state matrix is
\begin{equation}
  A_{\rm red}
  =
  f_x-f_y g_y^{-1}g_x
  \label{eq:nonlin_reduced_jacobian}
\end{equation}
at the equilibrium. The NEP chapters use this Jacobian, or a structured
second-order projection of it. The nonlinear theory must also study what
happens when the equilibrium moves, becomes singular, or ceases to attract
finite disturbances.

\section{Canonical Converter-Rich Family}

A useful hierarchical model for later chapters is
\begin{align}
  M\ddot\theta+D_0\dot\theta+P_e(\theta,V)+Gz &= p, \label{eq:nonlin_angle_layer}\\
  \dot z &= f_c(z,\theta,\dot\theta,V,p), \label{eq:nonlin_control_layer}\\
  0 &= g(\theta,V,z,p). \label{eq:nonlin_network_layer}
\end{align}
The state \(z\) may represent a synchronous-machine exciter or governor, a
grid-following PLL and current loop, a grid-forming droop or virtual
synchronous-machine controller, a voltage loop, a filter, or a plant-level
control state. Limiters and protections should be treated as smooth models
only inside a declared operating region; outside it they create hybrid
dynamics and require a different analysis.

Linearizing \eqref{eq:nonlin_angle_layer}--\eqref{eq:nonlin_network_layer}
recovers the QEP or Schur NEP used earlier. The nonlinear chapter keeps the
prelinearized equations visible so that a small-signal margin is not mistaken
for a large-signal certificate.

\section{Linear Controller Elimination as Memory}

For a linear controller block
\begin{equation}
  \dot z=A_cz+B_\theta\theta+B_\omega\dot\theta,
  \label{eq:nonlin_linear_controller}
\end{equation}
the exact variation-of-constants formula is
\begin{equation}
  z(t)
  =
  e^{A_ct}z(0)
  +
  \int_0^t
  e^{A_c(t-\tau)}
  \left[
  B_\theta\theta(\tau)+B_\omega\dot\theta(\tau)
  \right]d\tau .
  \label{eq:nonlin_linear_memory}
\end{equation}
Substitution into the retained equation produces an integro-differential
model. Its temporal kernel is
\begin{equation}
  K(t)=Ge^{A_ct}B,
  \label{eq:nonlin_memory_kernel}
\end{equation}
and its Laplace transform is the Schur-resolvent term
\begin{equation}
  G(sI-A_c)^{-1}B.
  \label{eq:nonlin_memory_transform}
\end{equation}
Thus the rational self-energy \(\Pi(s)\) is the spectral representation of
linear memory. This is the conceptual bridge between the NEP part and the
nonlinear-memory part of the dissertation.

\begin{theorem}[Self-energy as the tangent of memory]
\label{thm:nonlin_self_energy_tangent}
Consider a smooth nonlinear retained-hidden system near a regular equilibrium,
and suppose the hidden variables can be eliminated locally as a causal memory
functional acting on the retained trajectory. If this memory functional is
Fr\'echet differentiable at the equilibrium trajectory, then the Laplace
transform of its first variation is the Schur self-energy \(\Pi(s)\) of the
linearized hidden block.
\end{theorem}

\begin{proof}
Linearize the full nonlinear DAE or ODE around the equilibrium and partition
the linearization into retained and hidden coordinates. The hidden linear
subsystem has the variation-of-constants representation
\eqref{eq:nonlin_linear_memory}. Its contribution to the retained equation is
a convolution with kernel \(K(t)=Ge^{A_ct}B\). Taking the Laplace transform
gives \(G(sI-A_c)^{-1}B\), with the appropriate stiffness and velocity-channel
blocks collected as \(\Pi(s)\). This is exactly the first variation of the
nonlinear memory functional.
\end{proof}

\section{Nonlinear Elimination}

If the controller block is nonlinear,
\begin{equation}
  \dot z=f_c(z,\theta,\dot\theta,V,p),
  \label{eq:nonlin_nonlinear_controller}
\end{equation}
then eliminating \(z\) no longer gives a single matrix-valued rational function.
The reduced dynamics may contain history-dependent nonlinear functionals,
amplitude-dependent damping, and higher-order kernels. Locally, one can
organize those effects through:
\begin{itemize}
\item Volterra-series kernels when the nonlinear response is analytic and
small enough for a convergent expansion;
\item center-manifold or invariant-manifold reductions near critical
eigenvalues;
\item Mori-Zwanzig-type memory terms when unresolved variables are projected
out of the dynamics;
\item amplitude-dependent damping and stiffness coefficients when only the
dominant oscillatory envelope is retained.
\end{itemize}
The thesis will use this list as a research contract. A nonlinear memory claim
must state which reduction is being used and which finite-amplitude validation
gate supports it.

The same idea can be expressed as an amplitude-dependent self-energy when a
dominant oscillatory envelope is meaningful:
\begin{equation}
  \Pi(s;A)
  =
  \Pi_1(s)
  +
  A^2\Pi_3(s)
  +
  A^4\Pi_5(s)
  +\cdots .
  \label{eq:nonlin_amplitude_dependent_pi}
\end{equation}
Here \(\Pi_1(s)\) is the small-signal self-energy from Chapter~6, while higher
terms represent finite-amplitude stiffness and damping corrections. This
expansion is only legitimate under a declared Volterra, harmonic-balance, or
normal-form approximation. Its value is conceptual and testable: it asks
whether a small-signal damping reversal survives finite amplitude, disappears
through saturation, or becomes a Hopf mechanism. Recent non-Markovian closure
work reinforces the same principle that hidden variables generally reappear as
memory terms when only partial state information is retained
\cite{monsel2024neuraldde}.

\section{Bifurcation and Large-Signal Stability Contracts}

Small-signal stability asks whether the eigenvalues of the linearization have
negative real part. Nonlinear stability asks additional questions:
\begin{enumerate}
\item Does a regular equilibrium exist for the planned operating point?
\item Does the equilibrium lose stability through a Hopf, saddle-node, or
other bifurcation as \(p\) changes?
\item If a Hopf bifurcation occurs, what is the first Lyapunov coefficient and
does the emerging oscillation remain bounded?
\item What region of attraction is certified by an energy, Lyapunov,
contraction, or incremental-passivity argument?
\item Does a controller-induced damping reversal seen in the NEP persist under
finite perturbations, or is it only a local pole-motion effect?
\end{enumerate}

For a simple Hopf pair \(s(p)=\sigma(p)\pm j\omega(p)\), the spectral crossing
condition is
\begin{equation}
  \sigma(p_\star)=0,\qquad
  \omega(p_\star)\ne0,\qquad
  \frac{d\sigma}{dp}(p_\star)\ne0.
  \label{eq:nonlin_hopf_crossing}
\end{equation}
This condition locates a candidate bifurcation. It does not state the
amplitude, stability, or basin of the post-bifurcation motion. Those require
the nonlinear terms of \(F(x,p)\), not only the Schur NEP.

\section{Chapter Deliverables}

The nonlinear module is complete only when it contains the following
deliverables:
\begin{enumerate}
\item a regular DAE equilibrium theorem or a clearly stated regularity
assumption for each benchmark family;
\item a local-stability result tied to the QEP or Schur NEP;
\item at least one bifurcation calculation with finite-difference or complex
step verification of the derivative;
\item an energy, passivity, contraction, or Lyapunov-region argument for a
restricted model;
\item a finite-perturbation simulation showing whether the small-signal margin
predicts the observed nonlinear response;
\item a paper extraction section that states which nonlinear-memory result is
new, which result is known, and which result is only numerical evidence.
\end{enumerate}

This chapter is therefore not a replacement for the Schur NEP. It is the next
layer above it. The NEP preserves linearized memory; nonlinear analysis decides
which parts of that memory survive beyond infinitesimal perturbations.


\part{Applications and Reproducible Results}
\chapter{Risk Vectors, Weak Elements, and Planning Metrics}

The damping-operator hierarchy changes how weak elements and planning actions
should be defined. A weak bus, weak line, or useful converter placement is not
an intrinsic property of a graph alone. It is a property relative to a dynamic
margin, an operator class, an action family, and a validation gate. This
chapter turns the spectral derivatives into a planning language.

\section{Margins Must Name the Mechanism}

The most common small-signal margin is the critical damping ratio,
\begin{equation}
  S_\zeta=\log\frac{\zetaReq}{\zeta_{\rm crit}}.
  \label{eq:imp_zeta_margin}
\end{equation}
The case satisfies the requirement when \(S_\zeta\le0\). This margin is natural
for oscillatory eigenvalues, but it does not measure non-normal amplification,
input-output gain, or protection-weighted response.

For a state matrix \(A\), selected inputs \(B\), selected outputs \(C\), and
weights \(\Gamma_u,\Gamma_y\), a protected resolvent gain can be defined as
\begin{equation}
  g_{\rm prot}=
  \sup_{\omega\in\Omega}
  \sigma_1\left(\Gamma_yC(j\omega I-A)^{-1}B\Gamma_u\right).
  \label{eq:imp_protected_gain}
\end{equation}
The corresponding log margin is
\begin{equation}
  S_g=\log g_{\rm prot}.
  \label{eq:imp_Sg}
\end{equation}
The two margins can rank actions differently. That disagreement is not an
error; it means the mechanisms being measured are different.

\section{Action Coordinates}

Let a planning action be parameterized as
\begin{equation}
  \rho=
  \begin{bmatrix}
  \rho_L\\
  \rho_D\\
  \rho_M
  \end{bmatrix},
  \label{eq:action_vector}
\end{equation}
where \(\rho_L\) changes line weights or topology, \(\rho_D\) changes damping
or primary-control gain, and \(\rho_M\) changes physical or virtual inertia.
For a simple pole \(s_\star\), the total first-order movement is
\begin{equation}
  ds_\star
  =
  \nabla_Ls_\star^T d\rho_L
  +
  \nabla_Ds_\star^T d\rho_D
  +
  \nabla_Ms_\star^T d\rho_M.
  \label{eq:mixed_gradient}
\end{equation}
The three gradients are computed from the exact QEP sensitivity
\eqref{eq:fixed_full_qep_sens} or, after condensation, from the NEP
sensitivity \eqref{eq:self_nep_sens}. The key point is that the gradients are
not interchangeable. \(L\)-moves are stiffness moves, \(D\)-moves are
dissipation moves, and \(M\)-moves are inertia and normalization moves.

\section{Finite-Action Coordinates}

The action coordinates above are local, but the rational self-energy chapter
gives a finite-action object. For any low-rank action
\[
  T_u(s)=T(s)+U(s)K_u(s)V(s)^*,
\]
the action interaction matrix is
\begin{equation}
  \mathcal G_u(s)=V(s)^*T(s)^{-1}U(s),
  \label{eq:imp_action_green_matrix}
\end{equation}
and the protected roots created or shifted by the action satisfy
\begin{equation}
  \det\!\left[I+K_u(s)\mathcal G_u(s)\right]=0.
  \label{eq:imp_action_root_condition}
\end{equation}
This formula changes how a planning action should be reported. A derivative
\(ds_\star/du\) is the tangent of \eqref{eq:imp_action_root_condition}; it is
not the full action law. For a finite line increment, the scalar interaction is
the dynamic effective resistance
\(\mathcal R_e(s)=a_e^TT(s)^{-1}a_e\). For finite damping, inertia, and
controller-retuning actions, the relevant object is the corresponding block of
\(\mathcal G_u(s)\).

A local risk vector should therefore be interpreted as the first-order
linearization of three finite objects:
\begin{equation}
  \hbox{action determinant}
  \quad
  \hbox{protected margin boundary}
  \quad
  \hbox{structured fragility radius}.
  \label{eq:imp_finite_action_objects}
\end{equation}
This interpretation prevents a candidate from being accepted only because its
infinitesimal derivative is favorable at the base point. The finite action must
remain inside the protected margin and outside the structured pseudospectral
boundary over the declared budget.

\section{Risk Vector for a Candidate Action}

For any candidate action \(a\), define a local risk vector after normalizing the
action coordinate and each margin. A gain component must be dimensionless, e.g.
\(S_g=\log(g_{\rm prot}/g_{\rm req})\), not the logarithm of a dimensional
quantity. With a declared action scale \(\bar a\), derivatives below are taken
with respect to \(\hat a=a/\bar a\):
\begin{equation}
  r(a)=
  \begin{bmatrix}
  r_\sigma(a)\\
  r_\zeta(a)\\
  r_g(a)\\
  r_h(a)\\
  r_f(a)
  \end{bmatrix}.
  \label{eq:risk_vector}
\end{equation}
The components are:
\begin{align}
  r_\sigma(a)
  &=
  \left[
  \frac{d\Real(s_{\rm crit})}{da}
  \right]_+,
  \label{eq:risk_sigma}\\
  r_\zeta(a)
  &=
  \left[
  \frac{dS_\zeta}{da}
  \right]_+,
  \label{eq:risk_zeta}\\
  r_g(a)
  &=
  \left[
  \frac{dS_g}{da}
  \right]_+,
  \label{eq:risk_gain}\\
  r_h(a)
  &=
  \mathbf 1\left\{
  \frac{d\nu_k}{da}>0,\ 
  \frac{d\zeta_k}{da}<0
  \right\},
  \label{eq:risk_hidden}\\
  r_f(a)
  &=
  \hbox{normalized feasibility penalty}.
  \label{eq:risk_feasibility}
\end{align}
Here \([x]_+=\max(x,0)\). The first entry detects rightward pole movement. The
second detects damping-ratio margin loss. The third detects protected
input-output amplification. The fourth detects the hidden-margin conflict from
Chapter~4. The fifth records voltage, flow, current, protection, or dispatch
constraints.

All entries in \eqref{eq:risk_vector} are dimensionless local rates after the
chosen normalization. Without that convention, the angle between two risk
vectors can change simply because the same action is expressed in MW, per-unit,
or percent penetration.

A scalar screening score may be formed as
\begin{equation}
  \mathcal R(a)=w^Tr(a),\qquad w\ge0,
  \label{eq:scalar_risk_score}
\end{equation}
but the vector should be reported before scalarization. Otherwise, a strong
graph improvement can hide a damping loss.

\section{Bus Risk Vector}

For bus \(i\), define a fixed-\(D\) spectral bus vector
\begin{equation}
  r_i^{\rm bus}=
  \begin{bmatrix}
  \mathcal D_i\\
  \mathcal M_i\\
  \Pi_i\\
  H_i
  \end{bmatrix},
  \label{eq:bus_risk_vector}
\end{equation}
where
\begin{equation}
  \mathcal D_i=
  \sum_{k\in\mathcal K}
  \omega_k
  \frac{\phi_k(i)^2}{2\sqrt{\nu_k}}
  \label{eq:bus_D_score}
\end{equation}
is damping placement leverage,
\begin{equation}
  \mathcal M_i=
  \sum_{k\in\mathcal K}
  \omega_k\zeta_k\phi_k(i)^2
  \label{eq:bus_M_score}
\end{equation}
is inertia-dilution exposure under fixed \(D\),
\begin{equation}
  \Pi_i=
  \sum_{k\in\mathcal K}\omega_k\phi_k(i)^2
  \label{eq:bus_participation_score}
\end{equation}
is modal participation exposure, and \(H_i\) records whether incident graph
actions produce hidden damping conflicts. A high \(\mathcal D_i\) says
``put damping here.'' A high \(\mathcal M_i\) says ``do not put inertia here
without co-designed damping if the objective is damping ratio.'' A high
\(\Pi_i\) says the bus is dynamically visible.

\section{Line Risk Vector}

For a line \(e=(i,j)\), define
\begin{equation}
  r_e^{\rm line}=
  \begin{bmatrix}
  G_e^{(2)}\\
  G_e^{(K)}\\
  Z_e\\
  H_e\\
  F_e
  \end{bmatrix}.
  \label{eq:line_risk_vector}
\end{equation}
The algebraic-connectivity gain is
\begin{equation}
  G_e^{(2)}=(\phi_2(i)-\phi_2(j))^2.
  \label{eq:line_fiedler_gain}
\end{equation}
The effective-resistance gain is
\begin{equation}
  G_e^{(K)}
  =
  a_e^T(L^+)^2a_e.
  \label{eq:line_kirchhoff_gain}
\end{equation}
The damping-ratio risk for a monitored mode \(k\) is
\begin{equation}
  Z_e^{(k)}
  =
  -\frac{\partial\zeta_k}{\partial b_e},
  \qquad
  Z_e=\max_{k\in\mathcal K} Z_e^{(k)}.
  \label{eq:line_zeta_risk}
\end{equation}
The hidden-margin flag is
\begin{equation}
  H_e=\mathbf 1\left\{
  G_e^{(2)}>0,\ Z_e>0
  \right\},
  \label{eq:line_hidden_flag}
\end{equation}
and \(F_e\) is a feasibility penalty for thermal, voltage, protection, or
right-of-way constraints. A high \(G_e^{(2)}\) or \(G_e^{(K)}\) is not enough:
the line is attractive only if \(Z_e\) and \(F_e\) are acceptable.

\section{Weak Nodes Are Relative Objects}

A bus is weak only after the following items are specified:
\begin{enumerate}
\item the base operating point and dynamic model;
\item the operator class being used: uniform \(D=\gamma M\), commuting fixed
\(D\), non-commuting fixed \(D\), rational \(\Pi(s)\), or damping-channel
\(\Sigma(s)\);
\item the margin: damping ratio, decay rate, protected resolvent gain, or
another registered metric;
\item the action family: line change, damping addition, inertia change,
converter placement, controller retuning, or protection setting;
\item the accept/reject gate that turns a screen into evidence.
\end{enumerate}

Without those items, a weak-node list is not reproducible. A bus may be weak
for damping-ratio improvement under controller retuning and irrelevant for
protected-gain reduction under line reinforcement.

\section{Improving the Graph Without Adding Inertia}

When the objective is to improve graph robustness without adding inertia, the
candidate action should be selected from graph derivatives and then passed
through damping gates. A conservative procedure is:
\begin{enumerate}
\item rank candidate lines by \(G_e^{(2)}\) if the target is algebraic
connectivity;
\item rank candidate lines by \(G_e^{(K)}\) if the target is total effective
resistance or broad slow-mode coherence;
\item remove candidates with unacceptable \(F_e\);
\item remove or pair candidates with damping support if \(H_e=1\);
\item verify the remaining candidates with QEP or NEP sensitivity.
\end{enumerate}

This gives a mathematically defensible answer to the question ``how do we
improve \(L\) without adding inertia?'' Improve \(L\) where it moves the slow
graph modes most, but do not confuse that with damping unless the damping gate
passes.

\section{Substitutability of Actions}

Two actions are substitutable only if they improve the same margin through
compatible mechanisms. Adding physical damping, changing virtual damping,
placing a grid-forming converter, retuning a PLL, reinforcing a line, or
placing inertia can all move the critical eigenvalue, but they do not generally
act through the same operator term.

Let \(a\) and \(b\) be candidate actions. A local substitutability test is
\begin{equation}
  \cos\angle(r(a),r(b))
  =
  \frac{r(a)^Tr(b)}{\norm{r(a)}\,\norm{r(b)}}.
  \label{eq:risk_vector_angle}
\end{equation}
High alignment means the actions stress or relieve the same risks. Low
alignment means they are not substitutes even if both improve one scalar
metric. If one action improves \(S_\zeta\) by adding damping and another
improves \(S_g\) by reducing resolvent amplification, their risk vectors will
not generally align.

\section{Theory-to-Validation Protocol}

The theoretical hierarchy implies a protocol for future empirical work.

\begin{enumerate}
\item Start with the proportional model only as a baseline and state which
margin it certifies.
\item Measure non-commutation through \(\Gamma=Q^T\Dt Q\) before trusting
graph-mode damping claims.
\item If controller states are eliminated, construct or approximate
\(\Pi(s)\), then inspect its poles, residues, channel decomposition, and any
damping-channel \(\Sigma(s)\) over the frequency band of interest.
\item Use eigenvalue sensitivity to decompose action effects into retained and
condensed contributions.
\item Validate screens against full eigenvalue recomputation before using them
as planning rules.
\end{enumerate}

This protocol prevents a visually plausible graph or time-domain result from
being mistaken for a damping certificate.

\section{Implications for Converter-Rich Systems}

Converter-rich systems make the risk vector more important because converter
controls can supply both damping and destabilizing feedback. A grid-following
PLL, a grid-forming droop loop, a current controller, and a plant controller do
not merely add scalar damping. They add states. If those states are retained,
they appear in the full matrix \(A\). If they are condensed, they appear in
\(\Pi(s)\), or in its declared damping-channel specialization \(\Sigma(s)\).

The planning implication is direct. A reduced model that removes converter
states must show where their spectral effect went. If the answer is a constant
coefficient, the model is making a static approximation. If the answer is a
rational self-energy, the model has preserved the mechanism in the reduced
operator. The practical value of the theory is to make that distinction visible
before numerical ranking begins.

\chapter{Dynamic Hosting Capacity and Spectral Planning Levers}

The preceding chapters define the damping operator hierarchy. This chapter
uses that hierarchy to formulate the planning questions that motivated the
project: how much inverter-based resource (IBR) conversion can be hosted, where
virtual inertia or virtual damping should be placed, and which network
reinforcements can improve robustness before resorting to batteries or other
large storage devices. The common theme is that every planning answer must be
attached to a margin. Voltage feasibility, graph connectivity, damping ratio,
protected resolvent gain, and controller-resonance distance are different
objects.

\section{From Static Hosting to Spectral Hosting}

Let an operating point be described by complex bus voltages
\(V_i=|V_i|e^{j\theta_i}\) and bus admittance matrix
\(Y_{\rm bus}=G+jB\). The exact AC active- and reactive-power injections are
\begin{align}
  P_i(V,\theta)
  &=
  |V_i|\sum_{j=1}^n |V_j|
  \left(
  G_{ij}\cos(\theta_i-\theta_j)
  +
  B_{ij}\sin(\theta_i-\theta_j)
  \right),
  \label{eq:ac_active_power}\\
  Q_i(V,\theta)
  &=
  |V_i|\sum_{j=1}^n |V_j|
  \left(
  G_{ij}\sin(\theta_i-\theta_j)
  -
  B_{ij}\cos(\theta_i-\theta_j)
  \right).
  \label{eq:ac_reactive_power}
\end{align}
The linearized active-power stiffness is the angle block of the AC Jacobian.
After removing the uniform angle reference, write the corresponding positive
semidefinite stiffness as \(L_{\rm op}\). In the lossless small-angle
approximation, \(L_{\rm op}\) reduces to a weighted Laplacian.

A static hosting set for a vector of converter injections or conversions
\(\rho\in\mathcal U\) can be written abstractly as
\begin{equation}
  \mathcal H_{\rm stat}
  =
  \{\rho\in\mathcal U:\ c_{\rm V}(\rho)\le0,\ 
  c_{\rm th}(\rho)\le0,\ c_{\rm q}(\rho)\le0\},
  \label{eq:static_hosting_set}
\end{equation}
where the constraints may represent voltage, thermal, protection, power
quality, or operational limits. This is the classical hosting-capacity
language. It is indispensable, but it does not by itself certify small-signal
damping.

\begin{definition}[Dynamic spectral hosting set]
For a target damping ratio \(\zetaReq\), protected resolvent threshold
\(g_{\rm req}\), and admissible static set \(\mathcal H_{\rm stat}\), define
the dynamic spectral hosting set
\begin{equation}
  \mathcal H_{\rm dyn}
  =
  \left\{
  \rho\in\mathcal H_{\rm stat}:
  \zeta_{\min}(\rho)\ge \zetaReq,\ 
  g_{\rm prot}(\rho)\le g_{\rm req},\
  T(s;\rho)\hbox{ has no forbidden roots}
  \right\}.
  \label{eq:dynamic_hosting_set}
\end{equation}
For a penetration direction \(\alpha\ge0\), the directional dynamic hosting
capacity is
\begin{equation}
  h_{\rm dyn}(\alpha)
  =
  \sup\{\eta\ge0:\ \eta\alpha\in\mathcal H_{\rm dyn}\}.
  \label{eq:directional_dynamic_hosting}
\end{equation}
\end{definition}

This definition deliberately extends time-varying hosting-capacity ideas
\cite{oliveira2019dynamicHosting,jain2020dynamicHosting}. In distribution
planning, dynamic hosting often means that voltage or thermal violations are
counted with their duration and severity. Here it means that the electromechanical
and controller spectra also impose capacity limits. A system can be voltage
safe and still dynamically unsafe if the critical pole violates
\(\zetaReq\), if the protected resolvent gain is too large, or if a stable
network-family pole lies close to a controller pole with large self-energy
sensitivity.

\section{First-Order Spectral Hosting Screen}

In the fixed-\(D\) layer, let
\begin{equation}
  L(\rho)\phi_k(\rho)
  =
  \nu_k(\rho)M(\rho)\phi_k(\rho),
  \qquad
  \phi_k(\rho)^TM(\rho)\phi_\ell(\rho)=\delta_{k\ell}.
  \label{eq:hosting_generalized_modes}
\end{equation}
The diagonal damping-region screen from Chapter~4 is
\begin{equation}
  h_k(\rho)
  =
  \delta_k(\rho)-2\zetaReq\sqrt{\nu_k(\rho)}
  \ge0,
  \qquad
  \delta_k(\rho)=\phi_k(\rho)^TD(\rho)\phi_k(\rho).
  \label{eq:hosting_modal_gap}
\end{equation}
The quantity \(h_k\) is a modal headroom, not an eigenvalue. It is exact only
when the relevant modes decouple and becomes a screen otherwise.

For a simple generalized mode and small perturbation \(\rho_a\),
\begin{equation}
  \frac{\partial \nu_k}{\partial \rho_a}
  =
  \phi_k^T
  \left(
  L_a-\nu_k M_a
  \right)
  \phi_k,
  \label{eq:hosting_nu_derivative}
\end{equation}
where \(L_a=\partial L/\partial\rho_a\) and
\(M_a=\partial M/\partial\rho_a\). The direct damping derivative is
\(\phi_k^TD_a\phi_k\), with modal-rotation corrections when damping is
non-proportional. The first-order hosting-gap derivative is therefore
\begin{equation}
  g_{ka}
  =
  \frac{\partial h_k}{\partial\rho_a}
  =
  \frac{\partial\delta_k}{\partial\rho_a}
  -
  \frac{\zetaReq}{\sqrt{\nu_k}}
  \frac{\partial\nu_k}{\partial\rho_a}.
  \label{eq:hosting_gap_gradient}
\end{equation}

\begin{proposition}[Linear spectral hosting screen]
Assume \(h_k(0)>0\) for every monitored mode \(k\in\mathcal K\), and assume
that the first-order expansion of \eqref{eq:hosting_modal_gap} is used as a
tangent screen. Define
\begin{equation}
  A_{ka}=-g_{ka},
  \qquad
  b_k=h_k(0).
  \label{eq:hosting_polytope_matrices}
\end{equation}
Then the polytope
\begin{equation}
  \mathcal P_{\rm host}^{(1)}
  =
  \{\rho\in\mathcal U:\ A\rho\le b\}
  \label{eq:hosting_polytope}
\end{equation}
is the first-order spectral hosting screen associated with the damping-region
margin. It is not, by itself, a certified inner approximation of
\(\mathcal H_{\rm dyn}\).
\end{proposition}

\begin{proof}
The first-order expansion gives
\[
  h_k(\rho)=h_k(0)+\sum_a g_{ka}\rho_a+O(\norm{\rho}^2).
\]
Dropping the higher-order term and imposing \(h_k(\rho)\ge0\) gives
\(-\sum_ag_{ka}\rho_a\le h_k(0)\), which is \eqref{eq:hosting_polytope}.
\end{proof}

\begin{proposition}[Curvature-certified hosting region]
Suppose that, on a trust region \(\mathcal B_r=\{\rho:\norm{\rho}_2\le r\}\),
each monitored headroom satisfies
\[
  \norm{\nabla^2 h_k(\rho)}_2\le \beta_k .
\]
Then every \(\rho\in\mathcal U\cap\mathcal B_r\) satisfying
\begin{equation}
  h_k(0)+g_k^T\rho-\frac{\beta_k}{2}\norm{\rho}_2^2\ge0,
  \qquad k\in\mathcal K,
  \label{eq:hosting_curvature_certificate}
\end{equation}
also satisfies the second-order lower bound for the monitored modal headrooms.
\end{proposition}

\begin{proof}
Taylor's theorem with the stated Hessian bound gives
\[
  h_k(\rho)
  \ge
  h_k(0)+g_k^T\rho-\frac{\beta_k}{2}\norm{\rho}_2^2
\]
for all \(\rho\in\mathcal B_r\). Imposing
\eqref{eq:hosting_curvature_certificate} therefore certifies nonnegative
headroom for the monitored screen.
\end{proof}

The certified set is generally quadratic or conic, not a polytope. This is the
right object when the claim is an inner dynamic-hosting region rather than a
first-order ranking.

\section{Certified Hosting and Model Escalation}

A dynamic-hosting screen becomes a certified hosting statement only after the
model level is certified over the protected region. Let \(\Gamma\) enclose the
set of eigenvalues whose damping ratio, abscissa, or controller-resonance
distance limits hosting. For a reduced screen \(\widehat T(s;\rho)\), define
\begin{equation}
  \eta_\Gamma(\rho)
  =
  \sup_{s\in\Gamma}
  \left\|
  \widehat T(s;\rho)^{-1}
  \left[T(s;\rho)-\widehat T(s;\rho)\right]
  \right\|_2 .
  \label{eq:hosting_eta_gamma}
\end{equation}
If \(\eta_\Gamma(\rho)<1\), the reduced model preserves the protected
eigenvalue count inside \(\Gamma\). If \(\eta_\Gamma(\rho)\ge1\), the hosting
decision must escalate from the diagonal screen to the QEP, from the QEP to the
Schur NEP, or from the identified NEP to the full DAE, depending on where the
uncertainty enters.

For finite penetration increments, the same principle should be stated as a
fragility budget. Given a protected boundary \(\partial\mathcal S\), define
\begin{equation}
  r_{\rm host}(\rho)
  =
  \inf
  \left\{
  \|W_u u\|:
  \exists s\in\partial\mathcal S,\ 
  \det T(s;\rho,u)=0
  \right\}.
  \label{eq:hosting_fragility_budget}
\end{equation}
A hosting point is not robust merely because the first-order headroom is
positive. It is robust to the declared uncertainty family only if
\(r_{\rm host}\) exceeds the uncertainty or action budget assigned to that
point. This is the dynamic analogue of a voltage or thermal security margin.

\begin{remark}[SG-to-IBR conversion]
If action \(\rho_i\) replaces synchronous inertia and damping at bus \(i\)
without changing the network stiffness, then locally
\[
  M_i(\rho_i)=M_i^0-\rho_i\Delta M_i,\qquad
  D_i(\rho_i)=D_i^0-\rho_i\Delta D_i.
\]
The direct mode-frequency contribution is
\[
  \frac{\partial\nu_k}{\partial\rho_i}
  =
  \nu_k\Delta M_i\,\phi_k(i)^2.
\]
The direct damping contribution is approximately
\[
  \frac{\partial\delta_k}{\partial\rho_i}
  \approx
  -\Delta D_i\,\phi_k(i)^2
  +
  \hbox{modal-rotation terms}.
\]
Therefore SG-to-IBR conversion is most restrictive where the monitored mode
has large \(\phi_k(i)^2\) and the removed damping-to-inertia package is large.
This is the mathematical reason a hosting limit is not only a voltage limit.
\end{remark}

\section{Where to Place Virtual Inertia}

Virtual inertia is not a single planning lever. It may mean inertia-only
emulation, virtual synchronous-machine control with damping, grid-forming
droop with filtering, or a controller whose reduced effect is frequency
dependent. The fixed-\(D\) screen is still useful because it identifies what
would be true if the virtual device were represented locally by added mass and
added damping:
\begin{equation}
  M_i(\mu_i)=M_i^0+\mu_i,\qquad
  D_i(\mu_i)=D_i^0+\gamma_i^v\mu_i.
  \label{eq:virtual_inertia_local_model}
\end{equation}
Using the fixed-mode approximations of Chapter~4,
\begin{equation}
  \frac{\partial h_k}{\partial \mu_i}
  \approx
  \left(\gamma_i^v-\delta_k+\zetaReq\sqrt{\nu_k}\right)\phi_k(i)^2.
  \label{eq:virtual_inertia_hosting_gradient}
\end{equation}
Thus inertia-only support \((\gamma_i^v=0)\) can reduce the damping-region
headroom near the boundary \( \delta_k\approx2\zetaReq\sqrt{\nu_k}\). The same
bus can still be excellent for RoCoF or frequency-nadir objectives because
inertia reduces modal acceleration. It is not automatically excellent for
damping ratio.

For a set of monitored modes, a placement score that respects this distinction
is
\begin{equation}
  J_i^{\rm VI}
  =
  \sum_{k\in\mathcal K}\omega_k
  \left[
  \eta_{\rm f}\nu_k\phi_k(i)^2
  +
  \eta_{\rm h}
  \left(\gamma_i^v-\delta_k+\zetaReq\sqrt{\nu_k}\right)\phi_k(i)^2
  \right],
  \label{eq:virtual_inertia_score}
\end{equation}
where the first term rewards frequency-security leverage and the second term
rewards damping-headroom leverage. The weights \(\eta_{\rm f}\) and
\(\eta_{\rm h}\) must be declared before ranking. Otherwise the ranking hides
a tradeoff between slower frequency motion and better modal damping.

\section{Line-Shunt Taxonomy and Substitutability}

The project also asks whether the network can be improved before installing
batteries or local inertial devices. A structural answer comes from separating
line and shunt directions.

Let \(b_{ij}=e_i-e_j\) and define
\begin{equation}
  \mathcal L_{\rm all}
  =
  \operatorname{span}\{b_{ij}b_{ij}^T:\ 1\le i<j\le n\},
  \qquad
  \mathcal D_{\rm sh}
  =
  \{\operatorname{diag}(g):\ g\in\R^n\}.
  \label{eq:line_shunt_spaces}
\end{equation}

\begin{theorem}[Line-shunt decomposition]
Every real symmetric matrix \(X\in\R^{n\times n}\) has a unique decomposition
\begin{equation}
  X=X_{\mathcal L}+X_{\mathcal D},
  \qquad
  X_{\mathcal L}\in\mathcal L_{\rm all},\quad
  X_{\mathcal D}\in\mathcal D_{\rm sh}.
  \label{eq:line_shunt_decomposition}
\end{equation}
Equivalently,
\[
  \operatorname{Sym}(n)=\mathcal L_{\rm all}\oplus\mathcal D_{\rm sh}.
\]
\end{theorem}

\begin{proof}
Set the off-diagonal entries of \(X_{\mathcal L}\) equal to those of \(X\) and
choose the diagonal of \(X_{\mathcal L}\) so that every row sums to zero. Then
\(X_{\mathcal L}\) is a Laplacian-type matrix and
\(X_{\mathcal D}=X-X_{\mathcal L}\) is diagonal. The intersection of the two
spaces is zero because a diagonal zero-row-sum matrix must be zero.
\end{proof}

For local inertia,
\begin{equation}
  \frac{\partial\Lt}{\partial M_i}
  =
  -\frac{1}{2M_i}
  \left(E_i\Lt+\Lt E_i\right),
  \qquad
  E_i=e_ie_i^T.
  \label{eq:local_inertia_operator_location}
\end{equation}
This perturbation cannot create off-diagonal support on a non-existing edge.
Therefore it lies structurally in existing-line directions plus a shunt
residual:
\begin{equation}
  \frac{\partial\Lt}{\partial M_i}
  \in
  \mathcal L_{\rm ex}\oplus\mathcal D_{\rm sh}.
  \label{eq:inertia_existing_shunt}
\end{equation}
It has no new-line component. This statement does not mean inertia and line
reinforcement are dynamically equivalent. It says where the inertia action
lives in the perturbation space.

Let
\begin{equation}
  -\frac{\partial\Lt}{\partial M_i}
  =
  X_i^{\rm ex}+\operatorname{diag}(r_i),
  \qquad
  X_i^{\rm ex}\in\mathcal L_{\rm ex}.
  \label{eq:inertia_shunt_residual}
\end{equation}
For a critical mode \(q_c\), define the modal substitutability score
\begin{equation}
  s_i
  =
  1-
  \frac{|r_i^T(q_c\odot q_c)|}
  {\norm{r_i}_2\,\norm{q_c\odot q_c}_2+\varepsilon_{\rm reg}}.
  \label{eq:modal_substitutability_score}
\end{equation}
High \(s_i\) means the shunt residual is nearly invisible to the critical
mode; existing-line actions may be tried before local inertia. Low \(s_i\)
means the shunt residual aligns with the critical mode; topology-only actions
are less likely to substitute for a local inertial or synchronous-condenser
lever. Internal random-graph checks in \texttt{temp/poster/} verified the
algebraic residual of this decomposition to approximately \(10^{-14}\), which
is a structural verification rather than a power-system performance proof.
The score is still not a feasibility certificate for construction. A physical
line-only substitute must solve a constrained projection with nonnegative line
increments, corridor limits, and costs; membership in a linear span is only the
algebraic first screen.
\begin{equation}
  \min_{\Delta b\in\mathcal C_{\rm line}}
  \left\|
  -\frac{\partial\Lt}{\partial M_i}
  -
  \sum_{e\in\mathcal E_{\rm cand}}\Delta b_e a_ea_e^T
  \right\|_W,
  \qquad
  \mathcal C_{\rm line}=\{\Delta b\ge0:\ c^T\Delta b\le B,\ 
  \Delta b_e\le \bar b_e\}.
  \label{eq:line_inertia_constrained_projection}
\end{equation}

\section{Reinforcing Lines Without Batteries}

For a line \(e=(i,j)\), the modal-frequency derivative is
\begin{equation}
  \frac{\partial\nu_k}{\partial b_e}
  =
  (\phi_k(i)-\phi_k(j))^2.
  \label{eq:hosting_line_frequency}
\end{equation}
The damping-headroom derivative is
\begin{equation}
  \frac{\partial h_k}{\partial b_e}
  =
  \frac{\partial\delta_k}{\partial b_e}
  -
  \frac{\zetaReq}{\sqrt{\nu_k}}
  (\phi_k(i)-\phi_k(j))^2.
  \label{eq:hosting_line_headroom}
\end{equation}
This equation is the cleanest answer to the line-reinforcement question. A
line that strongly raises modal stiffness can still reduce damping headroom
unless it also raises modal damping through eigenvector rotation, damping
placement, or controller action.

The network-only benefits should still be used. If the objective is to improve
algebraic connectivity, use \(G_e^{(2)}=(\phi_2(i)-\phi_2(j))^2\). If the
objective is to reduce total effective resistance, use
\begin{equation}
  G_e^{(K)}
  =
  a_e^T(L^+)^2a_e,
  \qquad
  \frac{\partial}{\partial b_e}
  n\operatorname{trace}(L^+)
  =
  -nG_e^{(K)}.
  \label{eq:effective_resistance_line_gain}
\end{equation}
But a dynamic-hosting line screen should report the vector
\begin{equation}
  \ell_e
  =
  \begin{bmatrix}
  G_e^{(2)}\\
  G_e^{(K)}\\
  -\min_{k\in\mathcal K}\partial h_k/\partial b_e\\
  \partial S_g/\partial b_e\\
  F_e
  \end{bmatrix},
  \label{eq:line_screen_vector}
\end{equation}
where \(F_e\) is a feasibility penalty. A line should be called a robustness
improvement only after the selected margin improves and feasibility is
acceptable. Otherwise it is only a graph-stiffness improvement.

For controller-aware planning, the line report should be extended to a safe
reinforcement envelope:
\begin{equation}
  \mathcal R_e
  =
  \begin{bmatrix}
  G_e^{(2)}\\
  G_e^{(K)}\\
  \zeta_{p,e}^{\rm static}\\
  \zeta_{p,e}^{\rm NEP}\\
  \rho_{{\rm ctrl},e}\\
  \rho_{{\rm mem},e}\\
  p_{e,\max}^{\rm safe}
  \end{bmatrix}.
  \label{eq:line_safe_reinforcement_envelope}
\end{equation}
Here \(\rho_{{\rm ctrl},e}\) is a magnitude-normalized controller-channel
dominance score, \(\rho_{{\rm mem},e}\) is the memory-dominance score induced
by \(\Pi_s(s)\), and \(p_{e,\max}^{\rm safe}\) is the local action budget before
a protected margin is exhausted. The controller-distance term is
\(d^{\rm ctrl}_{k\ell}=|s_k-\mu_\ell|/(|s_k|+|\mu_\ell|)\), where
\(\mu_\ell\) is a condensed controller pole. A first-order version is
\begin{equation}
  p_{e,\max}^{\rm safe}
  =
  \min_{k,\ell}
  \left\{
  \frac{\zeta_k-\zetaReq}
       {[-\zeta_{p,e,k}^{\rm NEP}]_+},
  \frac{d^{\rm ctrl}_{k\ell}-d_{\rm req}}
       {[-d_{p,e,k\ell}^{\rm ctrl}]_+}
  \right\},
  \label{eq:line_safe_reinforcement_budget}
\end{equation}
with the convention that a zero denominator gives no active local restriction.
Finite upgrades must still be checked by continuation or full eigenvalue
recomputation.

This envelope separates four planning classes:
\begin{enumerate}
\item Safe reinforcement: graph strength and protected damping both improve.
\item Ordinary tradeoff: graph strength improves, but damping ratio falls for
the static stiffness-frequency reason.
\item Control-amplified reversal: the controller channel worsens a degradation
already visible in the static screen.
\item Controller-induced reversal: the static screen accepts the line but the
controller-aware model rejects it.
\end{enumerate}
The fourth class is the most publishable because it says the reduced static
model can choose the wrong reinforcement.

\section{Braess, Braess-Like, and Control-Resonant Reversals}

The literature uses Braess' paradox for cases where adding or increasing
network capacity decreases system performance
\cite{braess2005traffic,schafer2022braess}. In power systems, the performance
loss may be larger flow on an already stressed line, loss of synchrony, or
weaker linear stability \cite{witthaut2012braess,coletta2016braess}. The
present thesis separates three cases.

\begin{definition}[Three reinforcement reversals]
Let \(b_e\) be an upgraded line.
\begin{enumerate}
\item A flow Braess reversal occurs when an upgrade increases loading on a
protected line in the direction of its pre-existing flow.
\item A damping-margin reversal occurs when
\(\partial h_c/\partial b_e<0\) or
\(\partial\zeta_c/\partial b_e<0\) for the protected critical mode.
\item A control-resonant reversal occurs when the line action decreases the
damping of a network-family pole, moves it closer to a condensed controller
pole, and the self-energy term dominates the NEP sensitivity.
\end{enumerate}
\end{definition}

Only the first case is a classical flow Braess statement. The second is a
Braess-like spectral warning. The third is a stronger mechanism that belongs
to the rational self-energy layer:
\begin{equation}
  \frac{ds}{db_e}
  =
  -
  \frac{
  y^*\left(s\,\Sigma_{b_e}(s)+L_{b_e}\right)x
  }{
  y^*T_s(s)x
  }.
  \label{eq:control_resonant_line_sensitivity}
\end{equation}
If the dominant negative contribution comes from the \(s\Sigma_{b_e}(s)\)
channel rather than from \(L_{b_e}\) alone, the reinforcement is not explained
by scalar graph stiffness. It is a network-control resonance.

\section{Internal Validation Gates from the Project}

The project material preserved in \texttt{temp/} contains useful evidence, but
it must be read as gates and diagnostics. The key validated and negative
records are summarized in Table~\ref{tab:internal_validation_gates}.

\begin{table}[ht]
\centering
\caption{Internal validation gates used to calibrate the theory. These are not
universal claims; they define what the later empirical thesis must reproduce
or falsify under registered controls.}
\label{tab:internal_validation_gates}
\begin{tabular}{L{0.27\textwidth}L{0.35\textwidth}L{0.28\textwidth}}
\toprule
\rowcolor{softblue}
\thead{Gate} & \thead{Main record} & \thead{Thesis use}\\
\midrule
Adaptive damping screen &
In the \texttt{ieee39\_rational\_filter} output, 300 surrogate trials gave
median relative error \(4.68\times10^{-3}\) for the diagonal screen,
\(3.21\times10^{-4}\) for the second-order correction, and
\(3.69\times10^{-4}\) for the reduced QEP. &
Use diagonal damping as a fast screen and reduced QEP as a correction, not as
a universal replacement for the full model.\\
\addlinespace
Weak nodes and weak links &
The same surrogate ranked bus 30 as the strongest local damping action
(\(\Delta\zeta=1.12\times10^{-2}\)) and line 25--37 as the strongest positive
line action (\(\Delta\zeta=6.31\times10^{-3}\)); other buses and lines were
negative. &
Weak elements are action- and margin-relative. A centrality list is not enough.\\
\addlinespace
Braess-NEP phase &
The \texttt{phase\_braess\_nep} output tested 460 line-mode perturbations:
348 margin reversals, 232 network-family reversals, and 40 resonant reversals.
The clearest case was line 20--34 with self-energy-dominated
\(\Delta\zeta\) and a reported internal denominator metric of \(0.959\),
which must be recomputed as a magnitude-normalized memory or controller-channel
dominance index. &
Call the effect control-resonant reinforcement risk after NEP validation, not a
generic graph Braess theorem. The missing extraction is the strict subset with
\(d\zeta_{\rm static}/dp\ge0\) and \(d\zeta_{\rm NEP}/dp<0\).\\
\addlinespace
Inertia-control cross term &
The registered full-ANDES gate found 0/13 measurable singleton
GENROU-inertia by PLL-gain cross cases; the strongest cross-dominance index was
\(2.24\times10^{-4}\). &
Keep cross terms in the local theory, but do not use them as the empirical
headline.\\
\bottomrule
\end{tabular}
\end{table}

\FloatBarrier

\section{Measurement Gate for Frequency Security}

Dynamic hosting is also limited by what the protection and estimation layer can
see. The estimator draft in the project compares frequency estimators under
low-inertia IBR conditions and uses the trip-risk exposure
\begin{equation}
  T_{\rm risk}(e_{\rm th})
  =
  \sum_{m=0}^{N-1}\Delta t\,
  \mathbf 1_{\{|e[m]|>e_{\rm th}\}},
  \label{eq:trip_risk_duration}
\end{equation}
where \(e[m]\) is the estimator error relative to the reference frequency.
This metric is included here for one reason: a low RMS estimator can still be
unacceptable if its peak error or latency creates a long protection exposure.
Thus a complete dynamic hosting study should pair spectral margins with
measurement and protection margins. The theoretical hosting set
\eqref{eq:dynamic_hosting_set} can be extended by adding
\(T_{\rm risk}\le T_{\rm max}\) as a protection-aware constraint.

\section{Planning Rule}

The final planning rule is deliberately conservative.

\begin{enumerate}
\item Use static hosting constraints to reject voltage, thermal, and power
quality violations.
\item Use \(\mathcal P_{\rm host}^{(1)}\) to screen SG-to-IBR conversions by
damping headroom.
\item Rank virtual-inertia sites with separate RoCoF and damping-headroom
weights; do not mix the objectives silently.
\item Rank line upgrades by \(G_e^{(2)}\), \(G_e^{(K)}\), and
\(\partial h_k/\partial b_e\), then reject actions that improve graph
stiffness while degrading the protected dynamic margin.
\item If controller states are condensed, verify important candidates with the
NEP sensitivity \eqref{eq:control_resonant_line_sensitivity} or full
eigenvalue recomputation.
\end{enumerate}

This rule gives a defensible version of the practical advice: improve the
graph first when the weak mechanism is graph-mediated, add damping when the
weak mechanism is modal dissipation, add virtual inertia only with the correct
frequency-security objective or co-designed damping, and call a reinforcement
Braess-like only after the measured margin actually worsens.

\chapter{Conclusions}

This dissertation has developed an operator-valued theory of dynamic grid
strength in converter-rich power systems and used it to define dynamic hosting
capacity, virtual-inertia placement, weak elements, damping-aware network
reinforcement, and a traceable source-to-claim ledger for the project
artifacts. The hierarchy begins with the operational graph induced by the
power-flow operating point, continues through commuting fixed damping, isolates
the uniform-decay case \(D=\gamma M\), passes through fixed non-commuting
damping, continues to rational self-energy \(\Pi(s)\) after controller-state
condensation, introduces the Green function \(T(s)^{-1}\) as the carrier of
finite-action behavior, and then separates genuine nonlinear memory from
small-signal NEP theory. The central conclusion is that a damping, hosting, or
grid-strength claim is meaningful only relative to the operator, port/action
family, margin, and validation gate that carry the mechanism.

\section{Theoretical Conclusions}

A first conclusion is that graph topology becomes dynamic only after the
operating point is included. In the lossless retained model, the linearized
stiffness is
\[
  L(z^\star)=B W(z^\star)B^T,\qquad
  w_e=\kappa_e\cos\eta_e .
\]
Thus line loading, angle cohesion, cuts, cycles, inertia placement, and local
control fields enter the same graph-time operator. Static graph metrics are
useful screens, but the protected object is the response of
\(T(s;z^\star)=s^2M+sD+BWB^T+\Pi(s;z^\star)\).

The second conclusion is that \(D=\gamma M\) is a special floor, not the full
decouplable class. The exact fixed-\(D\) decoupling condition is
\([\Lt,\Dt]=0\). Under \(D=\gamma M\), graph modes decouple and every
underdamped mode has real part \(-\gamma/2\). Under Rayleigh or other
commuting damping, graph modes still decouple, but modal decay is not uniform.
This distinction gives a precise foundation for classical graph-frequency
intuition and also limits it.

The third conclusion is that fixed non-commuting damping breaks modal
decoupling. The graph eigenvectors diagonalize \(\Lt\), but the modal damping
matrix \(\Gamma=Q^T\Dt Q\) need not be diagonal. Off-diagonal damping creates
coupled modal dynamics, and eliminating the other modes produces rational
corrections. Thus, even a fixed-\(D\) system can require a frequency-dependent
modal interpretation when one insists on reducing to a single retained mode.

The fourth conclusion is that controller-state condensation changes the reduced
operator class. A Schur complement of the eliminated controller block gives
\[
  T(s)=s^2M+sD_0+L+\Pi(s),
  \qquad
  \Pi(s)=\Pi_K(s)+s\Pi_D(s).
\]
This rational correction has poles and residues inherited from the eliminated
dynamics. A constant damping matrix cannot represent those features exactly.
Static damping bridges are therefore local approximations whose validity must
be tested, not assumed.

The fifth conclusion is that dynamic effective resistance closes the loop
between graph metrics and finite actions. The Green function
\(G_T(s)=T(s)^{-1}\) defines
\[
  \mathcal R_e(s)=a_e^TG_T(s)a_e ,
\]
which reduces to classical effective resistance at the static graph limit and
becomes a meromorphic port compliance away from zero. A finite line action
satisfies \(1+\Delta b_e\mathcal R_e(s)=0\), and general low-rank actions
satisfy a small action determinant. Therefore local sensitivities are tangent
approximations of exact finite-action laws, not standalone planning objects.

The sixth conclusion is that reduced models need certificates. A QEP screen,
static self-energy approximation, data-identified \(\widehat\Pi(s)\), or
condensed NEP should be accepted over a protected contour only when the
relative perturbation bound
\[
  \sup_{\Gamma}\|\widehat T(s)^{-1}(T(s)-\widehat T(s))\|_2<1
\]
or an equivalent certificate preserves the protected eigenvalue count. A
margin derivative should be treated as a recommendation only when its sign is
larger than its error bound.

The seventh conclusion is that fragility is structured. An unstructured
pseudospectrum warns about non-normality, but a planning decision needs a
physical perturbation family: line uncertainty, damping, inertia, controller
gains, delays, operating-point variation, or identified-model error. The
resulting fragility radius measures the smallest admissible perturbation that
can move a pole to a protected boundary.

The eighth conclusion is that NEP theory and nonlinear stability are different
modules. The Schur NEP preserves linearized memory around an operating point.
The nonlinear DAE layer must still address equilibrium regularity,
bifurcations, finite-amplitude stability, and regions of attraction. A
small-signal damping margin is therefore a local gate, not a large-signal
certificate.

The ninth conclusion is that hosting capacity in converter-rich systems is
not only a voltage or thermal number. A static-feasible IBR penetration can
still violate a damping-ratio margin, a protected resolvent-gain threshold, or
a controller-resonance gate. The first-order spectral hosting screen is
therefore a useful tangent screen, but it becomes an inner certificate only with
curvature bounds and must be checked against the QEP or NEP when modal coupling
or controller self-energy is important.

The tenth conclusion is that graph reinforcement and dynamic robustness are
not synonymous. A line upgrade may improve algebraic connectivity or effective
resistance while reducing damping headroom. The correct planning object is a
vector containing graph gain, damping-margin derivative, protected-gain
derivative, feasibility, and possible Braess-like or control-resonant reversal
flags.

\section{Consequences for Research Practice}

The hierarchy imposes discipline on future simulation and planning studies.
Before a graph metric is used as evidence, the study must state whether the
operator is an operational graph at a fixed operating point, uniform
proportional damping, commuting fixed damping, non-commuting fixed damping,
rational self-energy, or nonlinear memory. Before a weak bus or weak line is
named, the margin and action family must be named. Before a reduced converter
model is trusted, the spectral effect of removed states must be visible in
either the retained matrix or the self-energy.

This thesis therefore separates three epistemic levels:

\begin{enumerate}
\item a certificate, which implies a margin under explicit assumptions;
\item a screen, which identifies candidates but does not imply the margin;
\item a gate, which tests whether the screen is reliable enough for the next
claim.
\end{enumerate}

The distinction is not editorial. It is necessary because the same graph
quantity can be a certificate under proportional damping, a screen under fixed
non-commuting damping, and an incomplete descriptor after controller
condensation.

\section{Open Theoretical Work}

Several theoretical problems remain open.

\begin{enumerate}
\item Develop graph-time resonance atlases that relate operating-point
geometry, graph-mode conversion, controller self-energy, and protected
frequency-response peaks.
\item Turn the total derivative of \(L=B W(z^\star)B^T\) into a continuation
tool for planning actions, so that line, dispatch, and converter changes are
distinguished from frozen-stiffness perturbations.
\item Develop computable error bounds between the full Schur NEP and static or
identified approximations of \(\Pi(s)\) or its damping-channel projection
\(\Sigma(s)\) over specified spectral regions.
\item Turn the dynamic effective resistance law
\(1+\Delta b_e\mathcal R_e(s)=0\) into scalable algorithms for finite line
screening and certified reinforcement budgets.
\item Define non-commutation measures tailored to inertia-normalized power
networks, with explicit links to damping-ratio error.
\item Connect passivity or positive-real properties of \(\Sigma(s)\) to
small-signal damping margins in non-normal retained systems.
\item Replace scalar passivity summaries with directional passivity matrices
or dynamic multipliers when MIMO channel geometry is important.
\item Develop a passivity-retuning protocol in which converter-controller
gains are adjusted to improve a band-limited self-energy passivity index while
also moving the protected Schur NEP pole in a damping-improving direction.
\item Compute structured fragility radii for topology, line loading, damping,
inertia, controller, delay, and operating-point uncertainty families.
\item Characterize when grid-strength or Laplacian placement metrics are
certificates, screens, or invalid descriptors for a given converter-control
family.
\item Test whether cellular-sheaf Laplacians with bus/controller stalks and
line restriction maps reveal weak nodes or weak links that scalar Laplacian
sensitivity, fixed-\(D\) QEP screens, and Schur self-energy diagnostics miss.
\item Build theory-to-validation protocols that combine NEP sensitivity,
pole-residue analysis, and full DAE eigenvalue recomputation.
\item Combine Loewner or contour-integral NEP extraction with continuation
tracking so that local pole-motion maps can be followed through high-dimensional
hosting and controller-retuning sweeps.
\item Extend the dynamic hosting screen beyond first order, including
non-commuting damping, controller self-energy, nonlinear memory, structured
fragility, and uncertainty in operating points.
\item Identify \(\Pi(s)\) from black-box converter data with stability,
causality, port, and passivity constraints, then propagate identification
error to certified planning decisions.
\item Register line-reinforcement and virtual-inertia planning baselines so
that Braess-like reversals and placement gains can be compared against
meaningful controls.
\end{enumerate}

\section{Closing Statement}

The transition from synchronous-machine dominated grids to converter-rich
grids is often described as a loss of inertia. That description is incomplete.
The deeper theoretical change is a change in the strength, damping, fragility,
and memory operator. Uniform viscous damping, commuting fixed damping,
non-commuting modal damping, rational controller self-energy, Green-function
action calculus, structured fragility, and nonlinear memory are different
mathematical objects. Static hosting, dynamic hosting, virtual inertia,
controller retuning, and line reinforcement are likewise different planning
questions unless they share the same port, operator, and protected margin. A
doctoral theory of dynamic grid strength must keep those objects separate and
then reconnect them through certificates. This dissertation provides that path:
from static topology to operational graph-time resonance, from effective
resistance to dynamic effective resistance, from fixed \(D\) to \(\Pi(s)\),
from \(\Pi(s)\) to structured fragility, and from linear self-energy to
nonlinear memory.


\appendix
\chapter{Integrated Claim Ledger and Algebraic Research Program}

The previous chapters built the theory in layers. This chapter records how the
project artifacts in \texttt{temp/} and the companion drafts map into that
theory. Its purpose is traceability. A doctoral thesis cannot rely on the
statement that a formula was ``modified to taste.'' It must state where the
formula comes from, what model it assumes, what it proves, what it only
screens, and what validation gate is still missing.

The central integrated object is
\begin{equation}
  T(s;\rho)
  =
  s^2M(\rho)+sD_0(\rho)+L(\rho)+\Pi(s;\rho),
  \qquad
  \Pi(s;\rho)=\Pi_K(s;\rho)+s\Pi_D(s;\rho).
  \label{eq:integrated_operator}
\end{equation}
Here \(L\) is the operating-point stiffness inherited from power flow,
\(M\) is physical or virtual inertia, and \(\Pi(s)\) is the condensed
controller self-energy. The damping-channel specialization defines
\(\Sigma(s)=D_0+\Pi_D(s)\) when \(\Pi_K\) is absent or explicitly absorbed into
stiffness. The fixed-\(D\) theory is the special case \(\Pi(s)=0\) with direct
damping \(D_0=D\). The uniform proportional damping floor is the smaller
special case \(D=\gamma M\), while the broader decoupled fixed-\(D\) class is
\([\Lt,\Dt]=0\). This single inclusion chain is the clean way to organize the
project:
\begin{equation}
  \begin{aligned}
  D=\gamma M
  &\subset
  \{D:\ [\Lt,\Dt]=0\}\\
  &\subset
  \{D:\ [\Lt,\Dt]\neq0,\ D\ {\rm fixed}\}\\
  &\subset
  \{\Pi(s)\ {\rm rational}\}\\
  &\subset
  \{{\rm nonlinear\ memory}\}.
  \end{aligned}
  \label{eq:integrated_inclusion_chain}
\end{equation}

\thesisbox{thesisblue}{Chapter visual promise}{
This chapter turns the project archive into a controlled research program. The
reader should be able to trace each idea from source, to operator, to margin,
to action, and finally to a validation gate.}

\begin{table}[ht]
\centering
\caption{Claim-status vocabulary used in this chapter.}
\label{tab:claim_status_vocabulary}
\begin{tabular}{L{0.19\textwidth}L{0.67\textwidth}}
\toprule
\rowcolor{softblue}
\thead{Status} & \thead{Meaning}\\
\midrule
\statuspill{thesisgreen}{Algebraic result} &
The statement follows from matrix identities or perturbation formulas under
declared assumptions.\\
\statuspill{thesisblue}{Screen} &
The calculation ranks candidates or flags risk, but does not by itself certify
the protected margin.\\
\statuspill{thesisorange}{Gate} &
The result is an internal validation or rejection test that constrains what can
be claimed next.\\
\statuspill{thesisred}{Candidate} &
The mechanism is mathematically coherent but still requires a registered
benchmark before it becomes an empirical contribution.\\
\bottomrule
\end{tabular}
\end{table}

\FloatBarrier

\section{Source-to-Claim Ledger}

Tables~\ref{tab:source_claim_ledger_main} and
\ref{tab:source_claim_ledger_validation} convert the reviewed material into
thesis claims. The language is deliberately conservative: an algebraic
derivation is not a benchmark validation, and a validated internal gate is not
an industry-wide theorem.

\begin{table}[p]
\centering
\small
\caption{Ledger linking the main project sources to mathematical objects and
claim status.}
\label{tab:source_claim_ledger_main}
\begin{tabular}{L{0.17\textwidth}L{0.25\textwidth}L{0.30\textwidth}L{0.12\textwidth}}
\toprule
\rowcolor{softblue}
\thead{Source} & \thead{Mathematical object} & \thead{Incorporated thesis claim} & \thead{Status}\\
\midrule
Handoff protocol &
\(\zeta_k\simeq\delta_k/(2\sqrt{\nu_k})\), adaptive reduced QEP, line
sensitivity warnings &
Graph-frequency metrics are screens; damping-margin rankings must use
post-action \(\zeta_{\min}\) or pole recomputation. &
Calibrated protocol\\
\addlinespace
Consolidated blueprint &
Damping-region inequality
\(\delta_k\ge 2\zeta_\star\sqrt{\nu_k}\), reduced QEP selection, exact QEP
sensitivity &
Use the diagonal damping region when modes are separated, then promote only
coupled low-damping clusters to a reduced QEP. &
Theory plus surrogate gate\\
\addlinespace
Unified attribution note &
\(T(s)=s^2M+sD_0+L+\Pi(s)\), self-energy contribution, weak nodes and weak
links &
The final reduced controller object for IBR-rich models is rational in \(s\);
its damping-channel projection is \(\Sigma(s)\), not generally a constant
\(D\). &
Core thesis frame\\
\addlinespace
Hidden-resonance note &
Distance \(d^{\rm ctrl}_{kj}\), self-energy dominance \(\phi_\Sigma^{\rm abs}\), hidden
action margin &
A stable pole may be close to a controller-resonant damping loss even when its
present damping ratio is acceptable. &
Candidate diagnostic\\
\addlinespace
Line-shunt derivation &
\(\operatorname{Sym}(n)=\mathcal L_{\rm all}\oplus\mathcal D_{\rm sh}\) and
\(\partial\Lt/\partial M_i\) &
Virtual inertia has an existing-line plus shunt structural footprint in the
inertia-normalized stiffness; it does not create new-line directions. &
Algebraic result\\
\bottomrule
\end{tabular}
\end{table}

\begin{table}[p]
\centering
\small
\caption{Ledger linking internal drafts and validation records to incorporated
claims.}
\label{tab:source_claim_ledger_validation}
\begin{tabular}{L{0.17\textwidth}L{0.25\textwidth}L{0.30\textwidth}L{0.12\textwidth}}
\toprule
\rowcolor{softblue}
\thead{Source} & \thead{Mathematical object} & \thead{Incorporated thesis claim} & \thead{Status}\\
\midrule
Hosting-capacity draft PDF &
Mass-normalized Laplacian, first-order spectral hosting screen, stochastic fragility,
edge and node fragility &
Dynamic hosting capacity should be a set constrained by damping, resolvent,
and stochastic frequency-security margins, not only voltage limits. &
Internal draft, integrated cautiously\\
\addlinespace
Beyond-voltage draft PDF &
Kirchhoff-consistent AC power flow, weak nodes, hidden dynamic margins &
Voltage-safe operating points can still violate dynamic spectral margins; the
power-flow Jacobian is the bridge from electrical operating point to \(L\). &
Internal draft, integrated cautiously\\
\addlinespace
Frequency-estimator draft and ZIP &
\(T_{\rm risk}\), latency, peak-error exposure &
A dynamic hosting certificate should include what the protection and
measurement layer can actually observe. &
Protection gate\\
\addlinespace
Braess-NEP phase &
NEP line sensitivity with self-energy contribution &
Some line reinforcements are control-resonant damping risks, not generic
topological Braess theorems. &
Internal NEP gate\\
\addlinespace
Rational-filter output &
Diagonal, second-order, reduced-QEP, and adaptive errors &
The adaptive estimator is useful because diagonal screens are often accurate
but fail under clustered or strongly coupled damping. &
Surrogate validation\\
\addlinespace
Coordination-lever gate &
Equal-cost topology, damping, and inertia packages &
Mixed lever packages can outperform the best single lever, but modal and
resolvent margins need separate reporting. &
Planning gate\\
\bottomrule
\end{tabular}
\end{table}

\FloatBarrier

\section{Step-by-Step Analysis Stack}

The analysis should be read as a sequence. Each step adds one layer of physics
or mathematics.

\begin{enumerate}
\item \textbf{Start from the electrical operating point.} Bus voltages,
current injections, and admittances satisfy \(I=Y_{\rm bus}V\). Complex power
is \(S_i=P_i+jQ_i=V_iI_i^*\). This gives the nonlinear AC power-flow map.
\item \textbf{Linearize active power with respect to angle.} At a feasible
operating point, the block \(J_{P\theta}=\partial P/\partial\theta\) gives the
synchronizing stiffness. After the reference-angle direction is removed, write
the positive stiffness as \(L\). In the lossless small-angle case, \(L\) is a
weighted Laplacian.
\item \textbf{Attach nodal dynamics.} Structure-preserving power-system models
retain bus-level dynamics instead of hiding all loads inside a reduced network
\cite{bergen1981structure,kundur1994power,sauer2017power,milano2010power}.
For the electromechanical layer this gives
\[
  M\ddot\theta+D\dot\theta+L\theta=p.
\]
\item \textbf{Remove the uniform angle mode.} Work on a subspace orthogonal to
the rotational symmetry. If \(R^T\mathbf 1=0\), define
\[
  M_r=R^TMR,\qquad D_r=R^TDR,\qquad L_r=R^TLR.
\]
\item \textbf{Normalize inertia.} With
\[
  \Lt=M_r^{-1/2}L_rM_r^{-1/2},\qquad
  \Dt=M_r^{-1/2}D_rM_r^{-1/2},
\]
the QEP becomes
\[
  \left(s^2I+s\Dt+\Lt\right)v=0.
\]
\item \textbf{Diagonalize stiffness, not damping.} Let
\(\Lt Q=Q\operatorname{diag}(\nu_k)\). The modal damping matrix is
\[
  \Gamma=Q^T\Dt Q.
\]
Only when \(\Gamma\) is diagonal do graph modes decouple exactly.
\item \textbf{Use the diagonal damping screen first.} The first screen is
\[
  \hat\zeta_k=\frac{\Gamma_{kk}}{2\sqrt{\nu_k}},
  \qquad
  h_k=\Gamma_{kk}-2\zeta_\star\sqrt{\nu_k}.
\]
It is exact under proportional or commuting damping, and a first-order
diagnostic otherwise.
\item \textbf{Promote dangerous clusters to a reduced QEP.} If off-diagonal
modal damping is large or the graph frequencies are clustered, solve the QEP
on the selected modal subspace.
\item \textbf{Condense controllers only after admitting that the object
changes.} Eliminating controller states produces \(\Pi(s)\), with
\(\Sigma(s)\) only as a declared damping-channel specialization, so the problem
becomes the rational NEP \eqref{eq:integrated_operator}.
\item \textbf{Plan with margins, not labels.} A bus, line, or controller is
weak only relative to a margin and an action family: \(\zeta_{\min}\), protected
resolvent gain, control-pole distance, stochastic variance, or protection
exposure.
\end{enumerate}

\begin{figure}[ht]
\centering
\begin{tikzpicture}[
  nodebox/.style={
    rectangle,
    rounded corners=3pt,
    draw=#1!70!black,
    fill=#1!10,
    text width=0.22\textwidth,
    minimum height=1.05cm,
    align=center,
    font=\small
  },
  arr/.style={-{Latex[length=2.2mm]}, thick, draw=thesisgray}
]
\node[nodebox=thesisblue] (op) {\textbf{Operating point}\\
\(V^\star,\theta^\star\)};
\node[nodebox=thesisgreen, right=0.45cm of op] (lm) {\textbf{Graph stiffness}\\
\(L,\Lt,\nu_k,q_k\)};
\node[nodebox=thesisorange, right=0.45cm of lm] (dm) {\textbf{Damping screen}\\
\(\delta_k,\hat\zeta_k,h_k\)};
\node[nodebox=thesisred, below=0.7cm of lm] (nep) {\textbf{Controller bridge}\\
\(\Sigma(s),s\Sigma'(s)\)};
\node[nodebox=thesisgold, right=0.45cm of nep] (act) {\textbf{Action ranking}\\
\(\Delta w,\Delta M,\Delta D,\Delta\kappa\)};
\draw[arr] (op) -- (lm);
\draw[arr] (lm) -- (dm);
\draw[arr] (dm.south) -- ++(0,-0.38) -| ([xshift=0.24cm]act.north);
\draw[arr] (lm.south) -- (nep.north);
\draw[arr] (nep) -- (act);
\end{tikzpicture}
\caption{Analysis stack used by the integrated research program. The diagonal
damping screen is the fast route; the controller bridge is invoked when static
damping cannot preserve the pole-motion mechanism.}
\label{fig:integrated_analysis_stack}
\end{figure}

\FloatBarrier

\section{Adaptive Reduced-QEP Algebra}

The off-diagonal entries of \(\Gamma\) are the algebraic reason the scalar
screen can fail. A compact coupling diagnostic is
\begin{equation}
  \eta_{\rm off}
  =
  \frac{\norm{\Gamma-\operatorname{diag}(\Gamma)}_F}{\norm{\Gamma}_F},
  \label{eq:eta_off}
\end{equation}
and a mode-pair diagnostic is
\begin{equation}
  c_{k\ell}
  =
  \frac{|\Gamma_{k\ell}|}
       {|\sqrt{\nu_k}-\sqrt{\nu_\ell}|+\varepsilon_{\rm gap}}.
  \label{eq:pair_coupling}
\end{equation}
Large \(c_{k\ell}\) means that the damping coupling is not small relative to
the frequency separation. This is a direct application of matrix perturbation
logic \cite{stewart1990matrix,kato1995perturbation}: small denominators make
apparently small off-diagonal terms important.

The adaptive estimator used by the project is therefore:
\begin{enumerate}
\item Compute \((\nu_k,q_k)\) of \(\Lt\) and \(\Gamma=Q^T\Dt Q\).
\item Rank all modes by \(\hat\zeta_k=\Gamma_{kk}/(2\sqrt{\nu_k})\).
\item Select low-\(\hat\zeta_k\) modes and augment the set with neighbors whose
\(c_{k\ell}\) exceeds a declared threshold.
\item Build \(V=[q_k]_{k\in\mathcal I}\), \(L_{\mathcal I}=V^T\Lt V\), and
\(D_{\mathcal I}=V^T\Dt V\).
\item Solve
\[
  \det(s^2I_{\mathcal I}+sD_{\mathcal I}+L_{\mathcal I})=0.
\]
\item Report the reduced-QEP damping margin and the gap to the diagonal
screen. If the gap is large, the graph-only explanation is not enough.
\end{enumerate}

The \texttt{ieee39\_rational\_filter} surrogate record supports this workflow:
over 300 trials the diagonal screen had median relative error
\(4.68\times10^{-3}\), while the second-order correction and reduced-QEP
variants reduced the median to approximately \(3.2\times10^{-4}\) and
\(3.7\times10^{-4}\), respectively. This does not prove universal accuracy. It
does justify using the diagonal formula as a fast screen and the reduced QEP as
a targeted correction.

\section{Exact Lever Sensitivities}

Let \(p\) be any scalar action: line weight, local inertia, local damping,
converter gain, or hosting-penetration coordinate. For the fixed-\(D\) QEP
\[
  T_Q(s,p)=s^2I+s\Dt(p)+\Lt(p),
\]
with right and left eigenvectors \(v\) and \(y\), a simple pole satisfies
\begin{equation}
  \frac{ds}{dp}
  =
  -
  \frac{y^*\left(s\,\Dt_p+\Lt_p\right)v}
       {y^*\left(2sI+\Dt\right)v}.
  \label{eq:integrated_qep_sensitivity}
\end{equation}
This is the QEP version of the nonlinear eigenvalue sensitivity used for
matrix polynomials and NEPs \cite{tisseur2001qep,guettel2017nep,beyn2012integral}.
If \(s=\sigma+j\omega\), \(ds/dp=a+jb\), and
\(\zeta=-\sigma/|s|\), then
\begin{equation}
  \frac{d\zeta}{dp}
  =
  \frac{\sigma\omega b-\omega^2a}{|s|^3}.
  \label{eq:zeta_from_pole_sensitivity}
\end{equation}
This equation is important because a line can move a pole left in real part
and still reduce damping ratio if it increases oscillation frequency too much.

For a line \(e=(i,j)\) in the inertia-normalized graph layer,
\begin{equation}
  \frac{\partial\nu_k}{\partial w_e}
  =
  \left(
  \frac{q_k(i)}{\sqrt{M_i}}-
  \frac{q_k(j)}{\sqrt{M_j}}
  \right)^2.
  \label{eq:integrated_line_nu}
\end{equation}
But damping requires more than \(\partial\nu_k/\partial w_e\). With fixed
modal labeling,
\begin{align}
  \frac{\partial \zeta_k}{\partial w_e}
  &=
  \frac{1}{2\sqrt{\nu_k}}\frac{\partial\delta_k}{\partial w_e}
  -
  \frac{\delta_k}{4\nu_k^{3/2}}
  \frac{\partial\nu_k}{\partial w_e},
  \label{eq:integrated_line_zeta}\\
  \frac{\partial\delta_k}{\partial w_e}
  &=
  q_k^T\Dt_{w_e}q_k
  +
  2\sum_{\ell\ne k}
  \frac{
    (q_\ell^T\Lt_{w_e}q_k)(q_\ell^T\Dt q_k)
  }{\nu_k-\nu_\ell}.
  \label{eq:integrated_delta_rotation}
\end{align}
The first term in \eqref{eq:integrated_delta_rotation} is direct damping
change. The sum is modal rotation. It is zero under proportional damping and
can be decisive under heterogeneous damping.

For the controller-condensed NEP \eqref{eq:integrated_operator},
\begin{equation}
  \frac{ds}{dp}
  =
  -
  \frac{
  y^*\left(s^2M_p+sD_{0,p}+L_p+\Pi_p(s)\right)x
  }{
  y^*\left(2sM+D_0+\Pi_s(s)\right)x
  }.
  \label{eq:integrated_nep_sensitivity}
\end{equation}
Here the subscript \(p\) denotes a total derivative after re-solving the
operating point and updating the controller blocks. The additional term
\(\Pi_s(s)\), which reduces to \(\Sigma(s)+s\Sigma'(s)-D_0\) in the
damping-channel case, is the spectral memory of the condensed controller block.
It is exactly the term that a static damping bridge removes.

\subsection{Mechanism-Resolved Attribution}

Equation~\eqref{eq:integrated_nep_sensitivity} is not only a way to compute a
pole derivative. It is also an attribution identity once the total derivative
is logged by mechanism. For a declared operating-point continuation, write
\begin{equation}
  T_p
  =
  T_p^{\rm net}
  +T_p^{\rm op}
  +T_p^{\rm ports}
  +T_p^{\rm ctrl}
  +T_p^{M,D}.
  \label{eq:mechanism_tp_split}
\end{equation}
The terms mean, respectively: direct network change, equilibrium or power-flow
shift, change in controller port visibility and residues, movement of
condensed controller poles, and direct mass-damping change. This split is not
unique until the model assembly is fixed, but after that bookkeeping choice it
is algebraic.

Let \(s=\sigma+j\omega\) and define the real linear map
\begin{equation}
  \mathcal D_\zeta(s)[\eta]
  =
  \frac{\sigma\omega\,\Imag \eta-\omega^2\,\Real \eta}{|s|^3}.
  \label{eq:zeta_linear_map}
\end{equation}
For each channel \(c\in\{\rm net,op,ports,ctrl,MD\}\), set
\begin{equation}
  g_c
  =
  \mathcal D_\zeta(s)
  \left[
  \frac{-y^*T_p^c x}{y^*T_sx}
  \right].
  \label{eq:channel_zeta_gain}
\end{equation}
Then
\begin{equation}
  \frac{d\zeta}{dp}
  =
  g_{\rm net}
  +g_{\rm op}
  +g_{\rm ports}
  +g_{\rm ctrl}
  +g_{M,D}.
  \label{eq:zeta_channel_decomposition}
\end{equation}
The denominator \(y^*T_sx\) is shared by all channels. Therefore the rational
memory term \(\Pi_s(s)\) does more than add one extra numerator contribution:
it can rotate or amplify every channel simultaneously. This is the causal
language needed to turn a pole plot into a mechanism statement.

\begin{definition}[Strict controller-induced reinforcement reversal]
For a line or parameter action \(p\), a strict controller-induced reinforcement
reversal occurs when a network-strength score improves and the static
reduction does not predict damping degradation, but the controller-aware NEP
does:
\begin{equation}
  \frac{dG_{\rm graph}}{dp}>0,
  \qquad
  \frac{d\zeta_{\rm static}}{dp}\ge0,
  \qquad
  \frac{d\zeta_{\rm NEP}}{dp}<0.
  \label{eq:strict_controller_reversal}
\end{equation}
Equivalently, the static and controller-aware models disagree in sign:
\begin{equation}
  \operatorname{sign}
  \left(\frac{d\zeta_{\rm static}}{dp}\right)
  \ne
  \operatorname{sign}
  \left(\frac{d\zeta_{\rm NEP}}{dp}\right).
  \label{eq:static_nep_sign_mismatch}
\end{equation}
\end{definition}

This definition is intentionally stronger than the statement
\(d\lambda_2/dp>0\) and \(d\zeta/dp<0\). The latter can occur in an elementary
QEP because increasing stiffness may raise modal frequency faster than decay.
The stronger claim is that the usual static reduction can recommend the wrong
line or parameter action.

\begin{proposition}[Controller-induced reversal certificate]
Assume a simple protected pole and the decomposition
\eqref{eq:mechanism_tp_split}. Suppose the rational self-energy admits a local
pole-residue expansion
\begin{equation}
  \Pi(s,p)
  =
  \sum_\ell \frac{Q_\ell(s,p)}{s-\alpha_\ell(p)}
  +\Pi_{\rm rem}(s,p),
  \label{eq:pi_pole_residue_expansion}
\end{equation}
and the contribution of all neglected or grouped terms to \(d\zeta/dp\) is
bounded above by \(\bar r\). Let
\begin{equation}
  g_0=g_{\rm net}+g_{\rm op}+g_{M,D}
  \label{eq:base_static_gain}
\end{equation}
and let \(g_\ell\) be the damping-ratio contribution of one dominant
controller pole \(\alpha_\ell\):
\begin{equation}
  g_\ell
  =
  \mathcal D_\zeta(s)
  \left[
  \frac{-y^*\Pi_{p,\ell}(s)x}{y^*T_sx}
  \right],
  \qquad
  \Pi_{p,\ell}(s)
  =
  \frac{Q_{\ell,p}(s)}{s-\alpha_\ell}
  +
  \frac{Q_\ell(s)\alpha_{\ell,p}}{(s-\alpha_\ell)^2}.
  \label{eq:dominant_pole_reversal_gain}
\end{equation}
If
\begin{equation}
  g_0\ge0,
  \qquad
  g_0+g_\ell+\bar r<0,
  \label{eq:controller_induced_reversal_certificate}
\end{equation}
then the first-order degradation is certified as controller-induced rather
than merely controller-correlated.
\end{proposition}

\begin{proof}
Linearity of \eqref{eq:integrated_nep_sensitivity} and
\eqref{eq:zeta_linear_map} gives the channel sum
\eqref{eq:zeta_channel_decomposition}. The pole-residue derivative in
\eqref{eq:dominant_pole_reversal_gain} separates residue change from movement
of the condensed pole. If the non-controller baseline is nonnegative and the
dominant controller-pole term plus the worst admissible remainder is negative,
then every compatible completion of the grouped remainder has negative total
first-order damping-ratio derivative. The sign reversal is therefore forced by
the rational controller channel.
\end{proof}

\section{Hidden Control-Resonance Margin}

Let \(\mu_j\in\operatorname{spec}(A_{cc})\) be a condensed controller pole and
let \(s_k\) be a retained network-family pole. Define
\begin{equation}
  d^{\rm ctrl}_{kj}
  =
  \frac{|s_k-\mu_j|}{|s_k|+|\mu_j|},
  \qquad
  \phi_\Sigma^{\rm abs}(k)
  =
  \frac{|y_k^*s_k\Sigma'(s_k)x_k|}
  {|y_k^*2s_kMx_k|+
   |y_k^*\Sigma(s_k)x_k|+
   |y_k^*s_k\Sigma'(s_k)x_k|+
   \varepsilon_{\rm reg}}.
  \label{eq:hidden_ctrl_metrics}
\end{equation}
Small \(d^{\rm ctrl}_{kj}\) and large \(\phi_\Sigma^{\rm abs}(k)\) mean that a presently
stable pole lies in the resolvent shadow of a controller pole. If \(A_{cc}\) is
diagonalizable with eigenvector condition factor \(\kappa_c\), then outside
\(\operatorname{spec}(A_{cc})\)
\begin{equation}
  \norm{(sI-A_{cc})^{-1}}
  \le
  \frac{\kappa_c}{\Delta_c(s)},
  \qquad
  \Delta_c(s)=\min_j |s-\mu_j|,
  \label{eq:resolvent_bound}
\end{equation}
and
\begin{equation}
  \norm{\Sigma'(s)}
  \le
  \frac{\norm{C}\norm{B}\kappa_c^2}{\Delta_c(s)^2}.
  \label{eq:sigma_prime_bound}
\end{equation}
Thus the warning is algebraic, not rhetorical: as a retained pole approaches a
condensed controller pole, the self-energy slope can become large.
\(\phi_\Sigma^{\rm abs}\) is a dominance index, not a strict fraction; the
absolute denominator avoids artificial inflation or suppression caused by
complex cancellation in \(y^*T_sx\).

For an admissible action \(p\) and target margin \(\zeta_\star\), define
\begin{equation}
  m_p^\zeta(k)
  =
  \frac{\zeta_k-\zeta_\star}
       {\left[-\partial\zeta_k/\partial p\right]_+},
  \qquad
  m_p^d(k,j)
  =
  \frac{d^{\rm ctrl}_{kj}-d_\star}
       {\left[-\partial d^{\rm ctrl}_{kj}/\partial p\right]_+}.
  \label{eq:hidden_action_margins}
\end{equation}
The hidden margin is
\begin{equation}
  m_p^{\rm hid}(k,j)=\min\{m_p^\zeta(k),m_p^d(k,j)\}.
  \label{eq:hidden_margin}
\end{equation}
This quantity says how much of action \(p\) can be applied before either the
damping threshold or the controller-distance threshold is locally exhausted.
The \texttt{phase\_braess\_nep} record gives one internal mechanism gate: 460
line-mode perturbations produced 40 strict resonant cases, and the clearest
line action moved a network-family mode toward a PLL-dominated condensed pole
with a self-energy-dominated damping decrease. The correct thesis claim is
therefore control-resonant reinforcement risk, not a universal Braess paradox.
The next extraction must count the smaller but stronger subset satisfying
\eqref{eq:strict_controller_reversal}; that subset is the empirical target for
the controller-induced reversal paper.

\subsection{Static Adequacy Gate and Resolvent Shadow}

The practical question is when the fixed-\(D\) or static-QEP screen is
sufficient. Let \(\Omega_k\) be a disk or contour neighborhood of a protected
pole \(s_k\) with radius \(r_k\). A first-order staticization error indicator is
\begin{equation}
  \varepsilon_{{\rm static},k}
  =
  \kappa_k r_k
  \sup_{\xi\in\Omega_k}\norm{\Pi_s(\xi)}_2,
  \qquad
  \kappa_k
  =
  \frac{\norm{x_k}_2\norm{y_k}_2}
       {|y_k^*T_s(s_k)x_k|+\varepsilon_{\rm reg}}.
  \label{eq:static_adequacy_gate}
\end{equation}
If \(\varepsilon_{{\rm static},k}<\tau m_k^{\rm margin}\) for a declared
tolerance \(\tau\), then the static screen is acceptable for that mode-action
pair. Otherwise the candidate lies in the resolvent shadow and must be checked
with the Schur NEP or the full DAE model. This turns the hierarchy into a
computational policy: use the cheap QEP broadly, escalate only shadowed modes,
and reserve full-model recomputation for finalists.

Distance to a condensed pole is not enough when \(A_{cc}\) is non-normal. The
screen should combine port-visible resolvent gain, modal projection, rational
memory, and pole conditioning:
\begin{align}
  \chi_c(s_k)
  &=
  \norm{C(s_kI-A_{cc})^{-1}(B_\theta+s_kB_\omega)}_2,
  \label{eq:port_visible_resolvent_gain}\\
  \chi_{k\ell}^{\rm proj}
  &=
  \frac{|y_k^*Q_\ell(s_k)x_k|}
       {|s_k-\alpha_\ell|\,\norm{y_k}_2\norm{x_k}_2+\varepsilon_{\rm reg}},
  \label{eq:projected_resolvent_shadow}\\
  \rho_{{\rm mem},k}
  &=
  \frac{|y_k^*\Pi_s(s_k)x_k|}
       {|y_k^*(2s_kM+D_0)x_k|+|y_k^*\Pi_s(s_k)x_k|
        +\varepsilon_{\rm reg}}.
  \label{eq:memory_dominance_index}
\end{align}
The desired empirical result is not that every large value produces a
reversal. It is that strict sign reversals concentrate where
\(\chi_{k\ell}^{\rm proj}\), \(\rho_{{\rm mem},k}\), and \(\kappa_k\) are large.
That statement can be tested with precision, recall, false-negative rate, and
the share of strict reversals captured after reviewing only the top decile of
line actions by resolvent-shadow score.

\section{Stochastic Fragility and Protection Exposure}

The hosting draft also contains a stochastic interpretation. In first-order
state form
\begin{equation}
  \dot x=Ax+B\xi,\qquad z=Cx,
  \label{eq:stochastic_state}
\end{equation}
with white input \(\xi\), the stationary covariance \(X\) satisfies
\begin{equation}
  AX+XA^T+BB^T=0,
  \qquad
  \mathbb E\norm{z}^2=\operatorname{trace}(CXC^T).
  \label{eq:lyapunov_variance}
\end{equation}
This is the same Gramian object used in network coherence analysis
\cite{bamieh2012coherence,tyloo2018robustness}. In a homogeneous proportional
limit, the graph contribution reduces to Kirchhoff-index-like quantities such
as \(\operatorname{trace}(L^+)\), which explains why effective resistance can
describe average vulnerability. In heterogeneous damping or rational
self-energy models, however, the covariance must be computed from the actual
linearization or an explicitly justified reduction.

The protection layer adds a separate constraint. Let \(e[m]\) be a frequency
estimator error and \(e_{\rm th}\) a protection threshold. The estimator draft
uses
\begin{equation}
  T_{\rm risk}
  =
  \sum_{m=0}^{N-1}\Delta t\,
  \mathbf 1_{\{|e[m]|>e_{\rm th}\}},
  \label{eq:integrated_trip_risk}
\end{equation}
which measures how long a protection-relevant error is exceeded. A complete
dynamic hosting set can therefore be written as
\begin{align}
  \mathcal H_{\rm dyn}^{\rm full}
  =
  \{\,\rho\in\mathcal H_{\rm stat}:&
  \ \zeta_{\min}(\rho)\ge\zeta_\star,\ 
  g_{\rm prot}(\rho)\le g_\star,\nonumber\\
  &
  \operatorname{trace}(C X(\rho) C^T)\le \sigma_\star^2,\ 
  T_{\rm risk}(\rho)\le T_\star,\ 
  d^{\rm ctrl}_{kj}(\rho)\ge d_\star
  \,\}.
  \label{eq:full_dynamic_hosting_set}
\end{align}
The point is not that every study must include every constraint. The point is
that each omitted constraint must be named.

\section{Pole-Motion Cone and Planning Optimization}

Collect line, inertia, damping, and controller actions in
\begin{equation}
  u=
  \begin{bmatrix}
    \Delta w & \Delta M & \Delta D & \Delta\kappa
  \end{bmatrix}^T.
  \label{eq:integrated_action_vector}
\end{equation}
For each monitored pole \(s_k\), NEP sensitivity gives
\begin{equation}
  \Delta s_k\simeq J_k u,
  \qquad
  J_k=
  \begin{bmatrix}
    \partial s_k/\partial p_1 & \cdots & \partial s_k/\partial p_r
  \end{bmatrix}.
  \label{eq:pole_motion_cone}
\end{equation}
Through \eqref{eq:zeta_from_pole_sensitivity}, this induces
\[
  \Delta\zeta_k\simeq g_k^Tu.
\]
Feasible actions therefore form a cone of pole motions in the complex plane.
Safe actions lie in the damping-improving half-space. Hidden-resonant actions
may still improve a graph metric while moving a pole toward a dangerous
controller resolvent region.

\begin{figure}[ht]
\centering
\begin{tikzpicture}[
  scale=0.95,
  axis/.style={->, thick, draw=thesisgray},
  safe/.style={fill=softgreen, draw=thesisgreen!70!black, line width=0.45pt},
  risk/.style={fill=softred, draw=thesisred!70!black, line width=0.45pt},
  arrow/.style={-{Latex[length=2.2mm]}, thick},
  lab/.style={font=\small}
]
\draw[axis] (-4.2,0) -- (2.2,0) node[right,lab] {\(\Real(s)\)};
\draw[axis] (0,-2.6) -- (0,2.6) node[above,lab] {\(\Imag(s)\)};
\draw[dashed, thesisgray] (-0.95,-2.3) -- (-0.95,2.3)
  node[above,lab] {target damping boundary};
\path[safe] (-2.6,0.35) -- (-3.65,1.45) -- (-3.65,-0.75) -- cycle;
\path[risk] (-2.6,0.35) -- (-1.05,1.15) -- (-0.95,-0.25) -- cycle;
\fill[thesisblue] (-2.6,0.35) circle (2.1pt)
  node[above left,lab] {current pole \(s_k\)};
\fill[thesisred] (-1.2,1.5) circle (2.1pt)
  node[above right,lab] {controller pole \(\mu_j\)};
\draw[arrow, thesisgreen] (-2.6,0.35) -- (-3.25,0.95)
  node[midway,above left,lab] {safe action};
\draw[arrow, thesisred] (-2.6,0.35) -- (-1.45,1.05)
  node[midway,below right,lab] {hidden-resonant action};
\node[lab, align=center, text width=0.26\textwidth] at (-3.35,-1.85)
  {larger damping ratio\\and larger controller distance};
\node[lab, align=center, text width=0.26\textwidth] at (1.0,-1.65)
  {rightward motion or\\controller approach is rejected};
\end{tikzpicture}
\caption{Pole-motion cone induced by the sensitivity matrix \(J_k\). A graph
action is acceptable only if its pole-motion vector improves the protected
margin and does not move the mode into a controller-resonant region.}
\label{fig:pole_motion_cone_visual}
\end{figure}

A local planning problem follows:
\begin{align}
  \max_{u,t}\quad & t \label{eq:planning_problem}\\
  \text{subject to}\quad
  & \zeta_k+g_k^Tu\ge t,\qquad k\in\mathcal K, \nonumber\\
  & d^{\rm ctrl}_{kj}+h_{kj}^Tu\ge d_\star,\qquad (k,j)\in\mathcal R,
    \nonumber\\
  & c^T|u|\le B,\qquad u_{\min}\le u\le u_{\max}, \nonumber\\
  & \text{AC, thermal, voltage, and protection feasibility gates hold.}
    \nonumber
\end{align}
This problem states exactly what the project can do with the new analysis:
rank SG-to-IBR conversions by dynamic headroom, identify where virtual damping
or inertia is honest, find line upgrades that improve damping rather than only
stiffness, detect control-resonant line risks, and compare mixed packages of
topology, damping, inertia, and control under equal cost.

\section{What This Analysis Enables}

The framework is useful because it converts broad planning questions into
declared algebraic tasks. Table~\ref{tab:analysis_enables} summarizes the
main actions.

\begin{table}[ht]
\centering
\small
\caption{Planning and research tasks enabled by the integrated analysis.}
\label{tab:analysis_enables}
\begin{tabular}{L{0.22\textwidth}L{0.25\textwidth}L{0.25\textwidth}L{0.16\textwidth}}
\toprule
\rowcolor{softblue}
\thead{Task} & \thead{Algebraic object} & \thead{Output} & \thead{Required gate}\\
\midrule
SG-to-IBR hosting direction &
\(h_k(\rho)=\delta_k-2\zeta_\star\sqrt{\nu_k}\) &
Maximum conversion direction before damping headroom is exhausted &
Full QEP or NEP recomputation\\
\addlinespace
Virtual inertia placement &
\(\partial h_k/\partial M_i\), RoCoF and damping weights &
Bus ranking that separates frequency support from damping support &
Post-action \(\zeta_{\min}\) and time-domain response\\
\addlinespace
Line reinforcement without batteries &
\(\partial\nu_k/\partial w_e\), \(\partial\zeta_k/\partial w_e\) &
Line list that improves margin, not only graph stiffness &
Flow, voltage, and damping checks\\
\addlinespace
Control-resonant risk audit &
\(d^{\rm ctrl}_{kj}\), \(\phi_\Sigma^{\rm abs}(k)\), \(s\Sigma'(s)\) &
Flag for stable poles that are close to controller-mediated damping loss &
Schur NEP or full DAE eigenanalysis\\
\addlinespace
Mixed lever package design &
Pole-motion cone \(J_ku\) and budget \(c^T|u|\le B\) &
Equal-cost packages of topology, damping, inertia, and control &
Registered baseline comparison\\
\addlinespace
Protection-aware hosting &
\(T_{\rm risk}\), estimator latency, protected gain &
Hosting set that includes what protection can observe &
Estimator and relay exposure test\\
\addlinespace
Self-energy passivity retuning &
\(\mu_\Sigma(\Omega)=\inf_{\omega\in\Omega}
\lambda_{\min}H_\Sigma(\omega)\) and \(ds_\star/d\rho\) &
Controller-gain update that improves both a passivity screen and the protected
pole-motion direction &
Schur NEP plus full-model retuning gate\\
\addlinespace
Sheaf weak-element research &
Cellular-sheaf Laplacian with bus/controller stalks and line restriction maps &
Candidate weak-node or weak-link certificate that is richer than scalar
Laplacian sensitivity &
New benchmark showing information not recovered by \(L\), QEP, or
\(\Sigma(s)\) alone\\
\bottomrule
\end{tabular}
\end{table}

\FloatBarrier

The last two rows are deliberately not ordinary planning claims. The passivity
row is a near-term extension because it uses the existing Schur self-energy
operator and adds a controller-respected dissipativity screen. The sheaf row is
a longer-term research program. It may become a powerful structural language
only if it produces a weak-element certificate that scalar graph theory and the
present self-energy diagnostics cannot already reproduce.

\section{Three Go--No-Go Studies}

The next work should be staged sequentially. First, the sign-characteristic
bridge must be grounded in the exact hypotheses of the matrix-polynomial
literature. Second, the passivity-retuning result should be made operational
because it extends the existing \(\Sigma(s)\) framework and can fail cleanly.
Third, the cellular-sheaf idea should be tested as a falsifiable structural
extension, not adopted as a new language by default.

\subsection{Sign-Characteristic Gate}

The strongest possible theoretical headline is not that controller condensation
makes a rational operator. That is Schur-complement mathematics. The stronger
claim is that a planning action can move a network-family pole through the
resolvent of a condensed controller and produce a sign collision that a
constant-\(D\) model cannot predict.

This claim needs a gate before it becomes a theorem. The gate has four steps:
\begin{enumerate}
\item Read the sign-characteristic definitions in the Hermitian
matrix-polynomial and self-adjoint meromorphic literature and list their
assumptions: regularity, symmetry, leading coefficient, elementary divisors, and
structure-preserving transformations.
\item State exactly which assumptions fail for
\(T(s)=s^2M+sD_0+L+\Pi(s)\): rationality, dissipation, non-normality, and loss
of a fixed symplectic or Hermitian product.
\item Build a two-bus SCR benchmark with a synchronous source and a
grid-following converter. The reduced coefficients must come from the exact
linearized algebraic-differential Jacobian; placeholders such as ``or similar''
are not allowed.
\item Track three quantities while varying SCR or PLL bandwidth:
\(\chi_k^{\rm op}\), \(\phi_\Sigma^{\rm abs}\), and the full-model eigenvalues.
The gate passes only if the same event appears as signature change or
singularity, self-energy dominance, and damping loss.
\end{enumerate}

If this gate passes, the first Transactions-style paper should be narrowly
framed around controller-resolvent reinforcement reversals:
\[
  \frac{dG_{\rm graph}}{dp}>0,
  \qquad
  \frac{d\zeta_{\rm static}}{dp}\ge0,
  \qquad
  \frac{d\zeta_{\rm NEP}}{dp}<0,
\]
with the degradation attributed to \(\Pi_p(s)\) or \(\Pi_s(s)\), not to the
static graph term alone. That paper should not mix frequency-estimator latency,
stochastic fragility, sheaves, and broad hosting optimization. The minimum
package is a structured Schur derivation, total NEP sensitivity, a
channel-resolved attribution, finite-difference validation, and ablation of
\(L_p\), \(\Pi_p\), and \(\Pi_s\). The empirical headline should be the rate
of wrong-sign static recommendations,
\[
  \Pr\left[
  \operatorname{sign}(\zeta_p^{\rm static})
  \ne
  \operatorname{sign}(\zeta_p^{\rm full})
  \right],
\]
plus the recall achieved by the resolvent-shadow filter when only a small
fraction of line actions are escalated to the full model.

\begin{table}[ht]
\centering
\caption{Minimal result package for the controller-induced reinforcement
reversal paper.}
\label{tab:controller_reversal_paper_package}
\begin{tabular}{L{0.18\textwidth}L{0.34\textwidth}L{0.34\textwidth}}
\toprule
\rowcolor{softblue}
\thead{Result} & \thead{Claim} & \thead{Evidence required}\\
\midrule
A & Exact Schur derivation and mechanism-resolved sensitivity attribution &
Derive \(\Pi(s)\), \(T_s\), and the channel gains \(g_c\).\\
\addlinespace
B & Robust sufficient certificate for controller-induced sign reversal &
Show \(g_0\ge0\) and \(g_0+g_\ell+\bar r<0\) under a registered remainder
bound.\\
\addlinespace
C & Static model wrong-sign rate &
Report strict cases where \(d\zeta_{\rm static}/dp\ge0\) but
\(d\zeta_{\rm full}/dp<0\).\\
\addlinespace
D & Resolvent-shadow computational filter &
Measure precision, recall, false negatives, and cases captured by the top
5--10\% of candidate line actions.\\
\bottomrule
\end{tabular}
\end{table}

\subsection{Passivity Retuning Study}

For a controller parameter \(\kappa\), such as a PLL bandwidth or integral gain,
the core object is
\[
  \mu_\Sigma(\Omega;\kappa)=
  \inf_{\omega\in\Omega}
  \lambda_{\min} H_\Sigma(\omega;\kappa).
\]
The computational question is whether \(\mu_\Sigma\) can be moved from negative
to positive after a stressed operating condition or an \(N-1\) contingency, while
the protected eigenvalue also moves toward the required damping region. The
study is deliberately narrow:
\begin{enumerate}
\item Select a contingency, IBR penetration level, and weak-grid operating
point; re-solve the power flow instead of freezing the pre-contingency
linearization.
\item Sweep \(\kappa\) over the admissible controller range and compute
\(\mu_\Sigma(\Omega;\kappa)\) from the Schur self-energy.
\item Record the active frequency \(\omega_\star\), active Hermitian eigenvector
\(q_\star\), and the envelope derivative
\(q_\star^H(\partial H_\Sigma/\partial\kappa)q_\star\).
\item Identify whether the margin is monotone, locally unimodal, or nonregular
because active frequencies switch.
\item At the best candidate \(\kappa_\star\), recompute the full DAE eigenvalues
and time-domain response. The passivity screen is a gate; the final claim still
needs the full-model check.
\item Compare against a static nodal-\(D\) screen. The negative control is
important: if a scalar damping change predicts the same action, the rational
self-energy did not add design information.
\end{enumerate}

A successful result is not merely a plot where an eigenvalue moves left. It is a
three-part certificate: the self-energy passivity margin improves, the protected
pole moves in the desired direction, and the full model confirms the damping
recovery. Failure is also useful. If \(\mu_\Sigma\) is nonmonotone with multiple
active-frequency switches, the thesis has identified a controller interaction
that cannot be summarized by one PLL retuning knob.

\subsection{Cellular-Sheaf Weak-Link Test}

The sheaf extension has one binary question:
\[
\text{Does a nonconstant sheaf Laplacian reveal a weak link that scalar }
\partial\nu_k/\partial w_e\text{ does not?}
\]
The smallest useful construction is enough for this test. Let each converter bus
carry a stalk \(\mathcal F(v_i)\simeq\R^2\) representing a local PLL phase and
frequency frame. Let each edge stalk \(\mathcal F(e_{ij})\) represent the frame
discrepancy that the line must reconcile. The restriction maps
\(\mathcal F_{v_i\to e_{ij}}\) and \(\mathcal F_{v_j\to e_{ij}}\) are determined
from the solved operating angles and the converter reference frames. With
identity restrictions, this reduces to the constant sheaf and recovers the
ordinary scalar graph intuition. With nontrivial frame rotations, the sheaf
coboundary
\[
  \delta_{\mathcal F}: C^0(G;\mathcal F)\to C^1(G;\mathcal F)
\]
produces the sheaf Laplacian \(L_{\mathcal F}=\delta_{\mathcal F}^\ast
\delta_{\mathcal F}\) on node stalks.

The go--no-go protocol is:
\begin{enumerate}
\item Build \(L_{\mathcal F}\) on the same 39-bus operating point used by the
QEP and Schur studies. The solved angles, not the bare topology, define the
restriction maps.
\item Compare the low spectrum of \(L_{\mathcal F}\) with the scalar
inertia-normalized Laplacian. The comparison should be made under the same
loading and IBR placement assumptions.
\item Remove or down-weight each candidate line and measure the collapse of the
lowest informative sheaf eigenvalues. This gives a sheaf weak-link ranking.
\item Compare that ranking against \(\partial\nu_k/\partial w_e\) and the
existing damping-risk vector. A material ranking difference is the first success
condition.
\item Compute the first harmonic obstruction
\(\dim\ker\delta_{\mathcal F}^\ast\) on high-angle subgraphs. A nonzero value
would indicate a frame-consistency obstruction, i.e., a synchronization
frustration not visible in a scalar Laplacian screen.
\item Validate any sheaf-only weak link in the full model. A sheaf ranking that
does not predict a changed damping, synchronization, or protection-relevant
response is only a mathematical artifact.
\end{enumerate}

The decision rule should be explicit. The sheaf program is a \emph{GO} only if
it produces a weak-link or weak-node ranking that differs materially from the
scalar ranking in at least one sweep case, and the difference has a physical
explanation in frame rotation, PLL alignment, or controller-coordinate
compatibility. It is a \emph{NO-GO} if the sheaf ranking reproduces the scalar
ranking up to scale. In that case the sheaf remains a future mathematical
extension, while the passivity-retuning paper remains the safer near-term
contribution.

\section{Mathematical Audit Gates Before Publication}

The following gates prevent the thesis from overclaiming:
\begin{enumerate}
\item \textbf{Structured Schur gate.} A condensed model must state whether the
rational correction is \(\Pi(s)\), \(\Pi_K(s)\), or \(s\Pi_D(s)\). Calling the
whole Schur correction ``dynamic damping'' is not allowed unless the stiffness
channel has been eliminated or justified.
\item \textbf{Sign-convention gate.} The active-power Jacobian must be tied to
one convention. This thesis uses exported electrical power, so \(L=J_{P\theta}\)
on retained angles.
\item \textbf{Invariant diagnostic gate.} Mode-family labels must use physical
weights, port visibility, left-right participation, or projectors. Raw
controller-state norm ratios are not admissible evidence.
\item \textbf{Curvature gate.} A first-order hosting polytope is a tangent
screen. An inner hosting certificate needs a trust region and a Hessian or model
error bound.
\item \textbf{Cluster gate.} Near repeated or clustered modes, use spectral
projectors and eigenvalue conditioning rather than derivatives of a named mode.
\item \textbf{Resolvent gate.} Controller-pole distance is only a proxy.
Non-normal \(A_{cc}\) requires resolvent gain or port-visible self-energy
dominance.
\item \textbf{Passivity gate.} A Hermitian part of \(\Sigma(j\omega)\) has a
passivity interpretation only after the conjugate ports and supply rate are
declared.
\item \textbf{Total-derivative gate.} A line, inertia, or controller action must
be differentiated through the re-solved operating point whenever the claim is a
planning claim rather than a frozen-linearization screen.
\end{enumerate}

\section{Validation Ladder}

Every claim should climb the same ladder:
\begin{enumerate}
\item \textbf{Algebraic identity.} Verify formulas such as
\(\partial\Lt/\partial M_i\) and the line-shunt decomposition to numerical
precision.
\item \textbf{Surrogate QEP.} Compare diagonal, second-order, reduced-QEP, and
full-QEP damping margins over controlled synthetic systems.
\item \textbf{IEEE topology QEP.} Use realistic network topology and operating
points, while still keeping the reduced second-order abstraction explicit.
\item \textbf{Schur NEP.} Condense controller states and test whether
\(\Pi(s)\), its damping-channel projection \(\Sigma(s)\), and resolvent
dominance explain pole motion.
\item \textbf{Full ANDES eigenanalysis.} Recompute the complete DAE
linearization after each selected action.
\item \textbf{Time-domain and protection gate.} Check RoCoF, nadir,
oscillation decay, estimator exposure, and protection-relevant thresholds.
\end{enumerate}

\begin{figure}[ht]
\centering
\begin{tikzpicture}[
  rung/.style={
    rectangle,
    rounded corners=3pt,
    draw=#1!70!black,
    fill=#1!10,
    text width=0.25\textwidth,
    minimum height=0.78cm,
    align=center,
    font=\small
  },
  arr/.style={-{Latex[length=2.1mm]}, thick, draw=thesisgray}
]
\node[rung=thesisblue] (r1) {1. Algebraic identity};
\node[rung=thesisgreen, right=0.36cm of r1] (r2) {2. Surrogate QEP};
\node[rung=thesisorange, right=0.36cm of r2] (r3) {3. IEEE topology QEP};
\node[rung=thesisred, below=0.55cm of r3] (r4) {4. Schur NEP};
\node[rung=thesisgold, left=0.36cm of r4] (r5) {5. Full ANDES};
\node[rung=thesisblue, left=0.36cm of r5] (r6) {6. Time-domain and protection};
\draw[arr] (r1) -- (r2);
\draw[arr] (r2) -- (r3);
\draw[arr] (r3) -- (r4);
\draw[arr] (r4) -- (r5);
\draw[arr] (r5) -- (r6);
\end{tikzpicture}
\caption{Validation ladder for this research program. A result becomes more
credible as it moves from algebraic identity to full linearization and
time-domain/protection checks.}
\label{fig:validation_ladder_visual}
\end{figure}

Negative gates are part of the thesis. The registered inertia-control cross
term is mathematically real in the local expansion, but the full-ANDES
singleton gate found no measurable effect in 13 tested cases. The correct use
is therefore as a possible second-order diagnostic, not as a headline. The same
discipline applies to hidden margins, line reinforcement, and virtual inertia:
the theory names the mechanism; validation decides whether it is a planning
claim in a given benchmark.

\chapter{Reproducibility of the Theoretical Manuscript}

This appendix records how to rebuild the thesis artifact. It intentionally does
not define a simulation protocol. The present manuscript is the theoretical
part of the dissertation.

\section{Source Layout}

The LaTeX source lives under:

\begin{verbatim}
reports/thesis/
\end{verbatim}

The main file is:

\begin{verbatim}
reports/thesis/main.tex
\end{verbatim}

The chapter files are:

\begin{verbatim}
reports/thesis/chapters/
\end{verbatim}

The bibliography is:

\begin{verbatim}
reports/thesis/references.bib
\end{verbatim}

\section{Compile the Thesis}

From the repository root, the manual build is:

\begin{verbatim}
cd reports/thesis
pdflatex -interaction=nonstopmode -halt-on-error main.tex
bibtex main
pdflatex -interaction=nonstopmode -halt-on-error main.tex
pdflatex -interaction=nonstopmode -halt-on-error main.tex
\end{verbatim}

The final PDF is copied to:

\begin{verbatim}
outputs/pdf/thesis_draft.pdf
\end{verbatim}

\section{Evidence Boundary}

This theoretical manuscript does not use plots or time-domain traces as
evidence. Numerical studies, when added later, should be treated as a separate
validation layer. A future validation artifact should record:

\begin{enumerate}
\item the full dynamic model and operating point;
\item the retained and condensed state partition;
\item the construction of \(\Pi(s)\), its channel decomposition, and any
\(\Sigma(s)\) approximation;
\item full-model eigenvalue targets;
\item accepted error tolerances;
\item failure cases and rejected hypotheses.
\end{enumerate}

This boundary protects the thesis from mixing mathematical derivation with
uncalibrated simulation evidence.

\chapter{Notation}

This appendix records the principal notation used in the theoretical thesis.

\begin{description}[leftmargin=3.2cm,style=nextline]
\item[\(s\)] Complex spectral parameter or eigenvalue.
\item[\(x\)] Dynamic state vector in a first-order DAE or state-space model.
\item[\(y\)] Algebraic variable vector in a DAE model.
\item[\(A\)] Linearized first-order state matrix after algebraic elimination.
\item[\(A_{ee}\)] Retained-state block of a partitioned state matrix.
\item[\(A_{cc}\)] Condensed controller or auxiliary-state block.
\item[\(A_{ec},A_{ce}\)] Coupling blocks between retained and condensed states.
\item[\(M\)] Diagonal inertia or mass matrix in a second-order retained model.
\item[\(D\)] Fixed damping matrix. The case \(D=\gamma M\) is the uniform
damping-to-inertia case, not the full decouplable class.
\item[\(\gamma_i\)] Local damping-to-inertia ratio \(D_i/M_i\) used in graph
roughness and commutation diagnostics.
\item[\(L\)] Synchronizing stiffness or active-power Jacobian Laplacian.
\item[\(B\)] Oriented branch-bus incidence matrix of the retained network.
\item[\(\eta\)] Branch angle-difference vector, \(\eta=B^T\theta\).
\item[\(\kappa_e\)] Nominal or parameter-dependent branch stiffness
coefficient before the operating-point cosine factor.
\item[\(f_e\)] Active power flow on branch \(e\), \(f_e=\kappa_e\sin\eta_e\)
in the lossless branch model.
\item[\(w_e\)] Operational branch stiffness \(w_e=\kappa_e\cos\eta_e\).
\item[\(W\)] Diagonal matrix of operational branch stiffnesses, so that
\(L=BWB^T\) on the angle subspace.
\item[\(\ones\)] Uniform angle-shift vector.
\item[\(\Lt\)] Inertia-normalized stiffness \(M^{-1/2}LM^{-1/2}\).
\item[\(\Dt\)] Inertia-normalized damping \(M^{-1/2}DM^{-1/2}\).
\item[\(\operatorname{comm}(\Lt,\Dt)\)] Commutator \(\Lt\Dt-\Dt\Lt\). For symmetric fixed
operators, \([\Lt,\Dt]=0\) is the exact modal decoupling condition.
\item[\(Q\)] Orthogonal modal basis of \(\Lt\) on the nonzero angle subspace.
\item[\(\nu_k\)] Nonzero eigenvalue of \(\Lt\), interpreted as squared modal
frequency in the undamped model.
\item[\(\phi_k\)] Generalized graph mode satisfying \(L\phi_k=\nu_kM\phi_k\)
and \(\phi_k^TM\phi_\ell=\delta_{k\ell}\).
\item[\(\delta_k\)] Frozen-mode scalar damping \(\phi_k^TD\phi_k\).
\item[\(\Lambda\)] Diagonal matrix of nonzero \(\Lt\) eigenvalues.
\item[\(\Gamma\)] Modal damping matrix \(Q^T\Dt Q\).
\item[\(E\)] Off-diagonal part of \(\Gamma\), responsible for modal damping
coupling.
\item[\(a_e\)] Incidence vector of line \(e=(i,j)\), usually \(a_e=e_i-e_j\).
\item[\(b_e\)] Susceptance-like or stiffness weight of line \(e\).
\item[\(L^+\)] Moore--Penrose pseudoinverse of the Laplacian \(L\).
\item[\(K_f\)] Kirchhoff index or total effective resistance, \(n\mathrm{trace}(L^+)\).
\item[\(P(s)\)] Quadratic eigenvalue operator \(s^2M+sD+L\).
\item[\(T(s)\)] Nonlinear retained second-order operator.
\item[\(G_T(s)\)] Dynamic Green function or retained compliance,
\(G_T(s)=T(s)^{-1}\) on the resolvent set.
\item[\(H_G(\omega)\)] Graph-time transfer obtained by projecting the
inertia-normalized operator into the graph-mode basis and evaluating at
\(s=j\omega\).
\item[\(\mathcal R_e(s)\)] Dynamic effective resistance of line \(e\),
\(a_e^TG_T(s)a_e\).
\item[\(S(s)\)] First-order Schur complement operator for retained states.
\item[\(\Pi(s)\)] General rational self-energy obtained after Schur
condensation of controller or auxiliary states.
\item[\(\widehat\Pi(s)\)] Identified or approximate self-energy obtained from
frequency scans, EMT data, PMU data, or black-box models.
\item[\(\Pi_K(s),\Pi_D(s)\)] Stiffness-like and damping-like channels in
\(\Pi(s)=\Pi_K(s)+s\Pi_D(s)\).
\item[\(\Sigma(s)\)] Damping-channel specialization
\(\Sigma(s)=D_0+\Pi_D(s)\), used only after the channel decomposition is
declared.
\item[\(\mathcal P_\Sigma(\omega)\)] Weighted directional passivity matrix of
the damping-channel self-energy.
\item[\(U,K_u,V\)] Low-rank action factors in
\(T_u(s)=T(s)+U(s)K_u(s)V(s)^*\).
\item[\(\mathcal G_u(s)\)] Action interaction matrix
\(V(s)^*T(s)^{-1}U(s)\).
\item[\(\eta_\Gamma\)] Contour adequacy error
\(\sup_{s\in\Gamma}\|\widehat T(s)^{-1}(T(s)-\widehat T(s))\|_2\).
\item[\(\Lambda_\varepsilon^{\mathcal A}(T)\)] Structured pseudospectrum under
the declared physical action or uncertainty family \(\mathcal A\).
\item[\(r_{\rm frag}\)] Structured dynamic fragility radius, the smallest
normalized admissible perturbation that moves a root to a protected boundary.
\item[\(r_{\rm top}\)] Structured topology fragility radius over admissible
line or edge-weight perturbations.
\item[\(r_e^{\C},r_e^{\R}\)] Complex and real frozen-line fragility radii for
line \(e\) on a declared protected spectral boundary.
\item[\(r_{\rm forced}\)] Smallest admissible forcing or input amplitude that
violates a declared protected response margin.
\item[\(r_{\rm ROA}\)] Nonlinear region-of-attraction fragility radius around
the protected equilibrium.
\item[\(D_0\)] Direct retained damping term inside the damping-channel
specialization.
\item[\(\zeta\)] Damping ratio of a complex pole.
\item[\(\zetaReq\)] Required damping ratio.
\item[\(S_\zeta\)] Log damping-ratio margin.
\item[\(g_{\rm prot}\)] Protected resolvent gain over selected frequency,
input, and output channels.
\item[\(S_g\)] Log protected-resolvent margin.
\item[\(\rho\)] Planning or hosting action vector, such as IBR conversion,
line reinforcement, damping addition, or virtual-inertia placement.
\item[\(\mathcal H_{\rm stat}\)] Static hosting set defined by voltage,
thermal, power-quality, protection, or operating constraints.
\item[\(\mathcal H_{\rm dyn}\)] Dynamic spectral hosting set after adding
damping-ratio, protected-gain, and forbidden-root constraints.
\item[\(h_{\rm dyn}(\alpha)\)] Directional dynamic hosting capacity in
penetration direction \(\alpha\).
\item[\(h_k\)] Modal damping-region headroom
\(\delta_k-2\zetaReq\sqrt{\nu_k}\).
\item[\(\mathcal P_{\rm host}^{(1)}\)] First-order spectral hosting polytope.
\item[\(\gamma_i^v\)] Virtual-damping-to-virtual-inertia coefficient at bus
\(i\) in a local fixed-\(D\) screen.
\item[\(J_i^{\rm VI}\)] Virtual-inertia placement score with declared
frequency-security and damping-headroom weights.
\item[\(\mathcal L_{\rm all}\)] Linear space spanned by all edge Laplacian
directions \(b_{ij}b_{ij}^T\).
\item[\(\mathcal L_{\rm ex}\)] Linear space spanned by existing-line Laplacian
directions.
\item[\(\mathcal D_{\rm sh}\)] Diagonal shunt space used in the line-shunt
decomposition.
\item[\(s_i\)] Modal substitutability score for replacing local inertia with
existing-line actions.
\item[\(T_{\rm risk}\)] Frequency-estimation trip-risk exposure duration.
\item[\(\pi_c\)] Fraction of eigenvector norm associated with condensed
controller states.
\item[\(\mathcal D_i\)] Fixed-\(D\) damping placement score at bus \(i\).
\item[\(\mathcal M_i\)] Fixed-\(D\) inertia-dilution exposure at bus \(i\).
\item[\(G_e^{(2)}\)] Algebraic-connectivity gain score for line \(e\).
\item[\(G_e^{(K)}\)] Effective-resistance gain score for line \(e\).
\item[\(H_i,H_e\)] Hidden damping-margin flags for bus- or line-related actions.
\item[\(r(a)\)] Risk vector for a candidate planning action \(a\).
\item[\(\mathcal R(a)\)] Scalarized risk score \(w^Tr(a)\) after vector
components have been reported.
\item[\texttt{GFL}] Grid-following converter control family.
\item[\texttt{GFM}] Grid-forming converter control family.
\end{description}


\backmatter
\bibliographystyle{ieeetr}
\bibliography{references}

\end{document}
